% Special Discrete Probability Distributions — Pure Mathematics with Statistics Upper Sixth — Cours signé OnBuch+
% Official MINESEC syllabus. Compiler avec : tectonic PMS12-special-discrete-probability-distributions.tex
\newcommand{\DOCMATIERE}{Pure Mathematics with Statistics}
\newcommand{\DOCNIVEAU}{Upper Sixth}
\newcommand{\DOCMODULE}{Upper Sixth — Pure mathematics with statistics}
\newcommand{\DOCLECON}{12}
\newcommand{\DOCTITRE}{Special Discrete Probability Distributions}
% OnBuch+ — préambule « cours signés » ENRICHI (compilé avec tectonic / XeLaTeX)
% Système visuel « Pop » d'OnBuch+ : fond crème (jamais blanc), bordures encre
% épaisses, ombres « dures » décalées, titres Archivo Black en capitales.
% Version MATHÉMATIQUES : graphes (pgfplots/tikz), boîtes pédagogiques
% (définition, propriété, méthode, attention, le savais-tu, exercice résolu,
% exercice, TP, bilan), page de garde et signature OnBuch+ ancrée en bas.
%
% Macros attendues dans main.tex AVANT \input{preamble} :
%   \DOCMATIERE  \DOCNIVEAU  \DOCMODULE  \DOCLECON (numéro)  \DOCTITRE
\documentclass[11pt]{article}

\usepackage{fontspec}
\usepackage[english]{babel}
\usepackage[a4paper,margin=2cm,top=3.4cm,bottom=2.8cm,headheight=1.4cm,footskip=1.3cm]{geometry}
\usepackage{amsmath,amssymb,amsfonts}
\usepackage{mathtools} % \xmapsto, \coloneqq... souvent utilisés par les modèles
\usepackage{cancel} % simplifications barrées (bilans de forces)
% \abs{z} (module, valeur absolue) et \norm{u} (norme), utilisés par les modèles
\providecommand{\abs}{}\let\abs\relax\DeclarePairedDelimiter{\abs}{\lvert}{\rvert}
\providecommand{\norm}{}\let\norm\relax\DeclarePairedDelimiter{\norm}{\lVert}{\rVert}
\usepackage[table]{xcolor} % [table] : fournit aussi \rowcolors (alternance de lignes)
\usepackage{pagecolor}
\usepackage{tikz}
\usetikzlibrary{arrows.meta,positioning,calc,shapes.geometric,shapes.misc,shapes.arrows,decorations.pathmorphing,decorations.pathreplacing,decorations.markings,patterns,shadows,mindmap,trees,fit,backgrounds,matrix,intersections,angles,quotes}
\usepackage{pgfplots}
\usepackage[version=4]{mhchem} % équations nucléaires \ce{^{235}_{92}U}
\usepackage[european,siunitx]{circuitikz} % schémas électriques
\pgfplotsset{compat=1.18}
\usepgfplotslibrary{fillbetween,groupplots} % aires entre courbes (calcul intégral)
\usepackage{enumitem}
\usepackage{fancyhdr}
% [most] charge toutes les bibliothèques tcolorbox (listings, minted, raster,
% documentation...) et gonfle inutilement la mémoire TeX sur un document long
% (~300 boîtes) : on ne charge que ce qui est réellement utilisé.
\usepackage[skins,breakable]{tcolorbox}
\usepackage{graphicx}
\usepackage{colortbl}
\usepackage{booktabs}
\usepackage{tabularx}
\usepackage{diagbox} % case d'en-tête diagonale des tableaux à double entrée
% Colonnes extensibles centrée / gauche / droite (C, L, R), souvent utilisées par les modèles
\newcolumntype{C}{>{\centering\arraybackslash}X}
\newcolumntype{L}{>{\raggedright\arraybackslash}X}
\newcolumntype{R}{>{\raggedleft\arraybackslash}X}
\usepackage{array}
\usepackage{multicol}
\usepackage{siunitx}
% parse-numbers=false : accepte aussi \SI{2,5 \times 10^{-3}}{...}
\sisetup{locale=UK,output-decimal-marker={.},group-minimum-digits=4,parse-numbers=false}
\DeclareSIUnit\ohm{\ensuremath{\symup{\Omega}}} % U+2126 absent de la police
\DeclareSIUnit\division{div} % réglages d'oscilloscope (ms/div, V/div)
\DeclareSIUnit\tr{tr} % tours (vitesse de rotation, ex. \SI{1500}{\tr\per\minute})
\usepackage{qrcode}
% \degree : commande fréquemment utilisée par les modèles pour un angle ou une
% température en dehors de \SI{}{} (non fournie par les paquets chargés).
\providecommand{\degree}{^{\circ}}
% \qed (fin de démonstration), utilisé par les modèles sans amsthm
\providecommand{\qed}{\hfill$\square$}
% \barre : double barre des valeurs interdites dans un tableau de variations (tkz-tab)
\providecommand{\barre}{\ensuremath{\|}}
% environnement proof (amsthm non chargé) utilisé par les modèles
\ifdefined\proof\else\newenvironment{proof}[1][Proof]{\par\noindent\textit{#1.}\ }{\qed\par}\fi
\usepackage{lastpage}
\usepackage{hyperref}
\hypersetup{hidelinks}

% ============================================================
% Palette « Pop » OnBuch+ — mêmes hex que la classe P (plus_ui.dart)
% ============================================================
\definecolor{poporange}{HTML}{F59321}
\definecolor{popdark}{HTML}{E07A0C}
\definecolor{popcream}{HTML}{FFF6E8}
\definecolor{popink}{HTML}{1C1714}
\definecolor{poppurple}{HTML}{7A5AE0}
\definecolor{popgreen}{HTML}{2E9E6A}
\definecolor{poppink}{HTML}{E0457B}
\definecolor{popblue}{HTML}{2F7DE1}
\definecolor{popgold}{HTML}{C98A12}
\definecolor{popmuted}{HTML}{8A7F76}
\definecolor{popline}{HTML}{EADFD2}
\definecolor{popgray}{HTML}{8A7F76}
\colorlet{popgrey}{popgray}
% Alias tolérants (le modèle nomme parfois les couleurs différemment)
\colorlet{popwhite}{white}
\colorlet{popblack}{popink}
\colorlet{popred}{poppink}
\colorlet{popyellow}{popgold}
% Teintes claires pour fonds de tableaux / zones
\colorlet{popblueL}{popblue!10!white}
\colorlet{poporangeL}{poporange!14!white}
\colorlet{popgreenL}{popgreen!12!white}
\colorlet{poppurpleL}{poppurple!12!white}
\colorlet{poppinkL}{poppink!12!white}
\colorlet{popgoldL}{popgold!14!white}

\pagecolor{popcream}

% ============================================================
% Typographie
% ============================================================
\defaultfontfeatures{Ligatures=TeX}
\setmainfont{PlusJakartaSans-Medium.ttf}[Path=../../../fonts/,BoldFont=PlusJakartaSans-Bold.ttf]
% Maths en sans-serif (Fira Math) avec chiffres et lettres droites de Plus
% Jakarta, pour que les formules se fondent dans le texte.
\usepackage[math-style=ISO,bold-style=ISO]{unicode-math}
\setmathfont{FiraMath-Regular.otf}[Path=../../../fonts/]
\setmathfont{PlusJakartaSans-Medium.ttf}[Path=../../../fonts/,range={up/num,up/{Latin,latin},\mathup}]
% (icomma retiré : décimales avec point en anglais)
% Caractères mathématiques Unicode absents des polices de texte (Plus Jakarta,
% Archivo Black) : rendus par leur équivalent mathématique, en texte comme en
% formule — sinon ils s'affichent en carré vide (ex. « Arithmétique dans ℤ »).
\usepackage{newunicodechar}
\newunicodechar{ℝ}{\ensuremath{\mathbb{R}}}
\newunicodechar{ℤ}{\ensuremath{\mathbb{Z}}}
\newunicodechar{ℂ}{\ensuremath{\mathbb{C}}}
\newunicodechar{ℕ}{\ensuremath{\mathbb{N}}}
\newunicodechar{ℚ}{\ensuremath{\mathbb{Q}}}
\newunicodechar{𝔻}{\ensuremath{\mathbb{D}}}
\newunicodechar{α}{\ensuremath{\alpha}}
\newunicodechar{β}{\ensuremath{\beta}}
\newunicodechar{γ}{\ensuremath{\gamma}}
\newunicodechar{δ}{\ensuremath{\delta}}
\newunicodechar{ε}{\ensuremath{\varepsilon}}
\newunicodechar{θ}{\ensuremath{\theta}}
\newunicodechar{λ}{\ensuremath{\lambda}}
\newunicodechar{μ}{\ensuremath{\mu}}
\newunicodechar{σ}{\ensuremath{\sigma}}
\newunicodechar{τ}{\ensuremath{\tau}}
\newunicodechar{φ}{\ensuremath{\varphi}}
\newunicodechar{ψ}{\ensuremath{\psi}}
\newunicodechar{ω}{\ensuremath{\omega}}
\newunicodechar{Σ}{\ensuremath{\Sigma}}
% majuscules produites par \MakeUppercase dans les titres (α-aminés -> Α)
\newunicodechar{Α}{\ensuremath{\alpha}}
\newunicodechar{Β}{\ensuremath{\beta}}
\newunicodechar{Γ}{\ensuremath{\Gamma}}
\newunicodechar{Δ}{\ensuremath{\Delta}}
\newunicodechar{Ε}{\ensuremath{\varepsilon}}
\newunicodechar{Θ}{\ensuremath{\Theta}}
\newunicodechar{Λ}{\ensuremath{\Lambda}}
\newunicodechar{Μ}{\ensuremath{\mu}}
\newunicodechar{Τ}{\ensuremath{\tau}}
\newunicodechar{Φ}{\ensuremath{\Phi}}
\newunicodechar{Ψ}{\ensuremath{\Psi}}
\newunicodechar{Ω}{\ensuremath{\Omega}}
\newunicodechar{↦}{\ensuremath{\mapsto}}
\newunicodechar{∈}{\ensuremath{\in}}
\newunicodechar{∉}{\ensuremath{\notin}}
\newunicodechar{∧}{\ensuremath{\wedge}}
\newunicodechar{⇔}{\ensuremath{\Leftrightarrow}}
\newunicodechar{⟺}{\ensuremath{\Longleftrightarrow}}
\newunicodechar{⇒}{\ensuremath{\Rightarrow}}
\newunicodechar{⟹}{\ensuremath{\Longrightarrow}}
\newunicodechar{∘}{\ensuremath{\circ}}
\newunicodechar{∪}{\ensuremath{\cup}}
\newunicodechar{∩}{\ensuremath{\cap}}
\newunicodechar{≡}{\ensuremath{\equiv}}
\newunicodechar{⊂}{\ensuremath{\subset}}
\newunicodechar{⊥}{\ensuremath{\perp}}
\newunicodechar{∃}{\ensuremath{\exists}}
\newunicodechar{∀}{\ensuremath{\forall}}
\newunicodechar{∛}{\ensuremath{\sqrt[3]{\,}}}
\newunicodechar{∜}{\ensuremath{\sqrt[4]{\,}}}
\newunicodechar{⃗}{\ensuremath{{}^{\to}}}
\newunicodechar{ⁿ}{\ifmmode^{n}\else\textsuperscript{\ensuremath{n}}\fi}
\newunicodechar{ˣ}{\ifmmode^{x}\else\textsuperscript{\ensuremath{x}}\fi}
\newunicodechar{ᵇ}{\ifmmode^{b}\else\textsuperscript{\ensuremath{b}}\fi}
\newunicodechar{ʳ}{\ifmmode^{r}\else\textsuperscript{\ensuremath{r}}\fi}
\newunicodechar{ᵘ}{\ifmmode^{u}\else\textsuperscript{\ensuremath{u}}\fi}
\newunicodechar{ʸ}{\ifmmode^{y}\else\textsuperscript{\ensuremath{y}}\fi}
\newunicodechar{ᵃ}{\ifmmode^{a}\else\textsuperscript{\ensuremath{a}}\fi}
\newunicodechar{ᵉ}{\ifmmode^{e}\else\textsuperscript{\ensuremath{e}}\fi}
\newunicodechar{ᵏ}{\ifmmode^{k}\else\textsuperscript{\ensuremath{k}}\fi}
\newunicodechar{ᵖ}{\ifmmode^{p}\else\textsuperscript{\ensuremath{p}}\fi}
\newunicodechar{ᴺ}{\ifmmode^{N}\else\textsuperscript{\ensuremath{N}}\fi}
\newunicodechar{ᶜ}{\ifmmode^{c}\else\textsuperscript{\ensuremath{c}}\fi}
\newunicodechar{ʰ}{\ifmmode^{h}\else\textsuperscript{\ensuremath{h}}\fi}
\newunicodechar{ᵠ}{\ifmmode^{q}\else\textsuperscript{\ensuremath{q}}\fi}
\newunicodechar{ₙ}{\ifmmode_{n}\else\textsubscript{\ensuremath{n}}\fi}
\newunicodechar{ₐ}{\ifmmode_{a}\else\textsubscript{\ensuremath{a}}\fi}
\newunicodechar{ᵢ}{\ifmmode_{i}\else\textsubscript{\ensuremath{i}}\fi}
\newunicodechar{ᵣ}{\ifmmode_{r}\else\textsubscript{\ensuremath{r}}\fi}
\newunicodechar{ₖ}{\ifmmode_{k}\else\textsubscript{\ensuremath{k}}\fi}
\newunicodechar{ᵦ}{\ifmmode_{\beta}\else\textsubscript{\ensuremath{\beta}}\fi}
\newunicodechar{ₓ}{\ifmmode_{x}\else\textsubscript{\ensuremath{x}}\fi}
\newfontfamily\popdisplay{ArchivoBlack-Regular.ttf}[Path=../../../fonts/]
\newfontfamily\poplogo{SpaceGrotesk-Bold.ttf}[Path=../../../fonts/]
\newfontfamily\popsemi{PlusJakartaSans-SemiBold.ttf}[Path=../../../fonts/]
\newfontfamily\popxbold{PlusJakartaSans-ExtraBold.ttf}[Path=../../../fonts/]
\newfontfamily\popxboldls{PlusJakartaSans-ExtraBold.ttf}[Path=../../../fonts/,LetterSpace=3.0]

\setlist[enumerate]{leftmargin=1.7em,itemsep=5pt,topsep=4pt}
\setlist[itemize]{leftmargin=1.7em,itemsep=4pt,topsep=4pt,label={\color{poporange}\textbullet}}
\setlength{\parindent}{0pt}
\setlength{\parskip}{5pt}
\renewcommand{\arraystretch}{1.25}

% pgfplots : style Pop par défaut
\pgfplotsset{
  every axis/.append style={axis line style={popink,line width=1pt},tick style={popink},
    grid style={popline,line width=0.5pt},label style={font=\small},tick label style={font=\footnotesize}},
  popcurve/.style={poporange,line width=1.8pt,no markers,smooth},
}
% Même style disponible dans un \draw TikZ simple (hors pgfplots)
\tikzset{popcurve/.style={poporange,line width=1.8pt}}
\tikzset{popgrid/.style={popline,very thin}}
\colorlet{popcurve}{poporange} % le nom du style est parfois utilisé comme couleur

% ============================================================
% Pill / squiggle
% ============================================================
\newcommand{\poppill}[3]{%
  \tikz[baseline=-0.55ex]\node[draw=popink,line width=1.2pt,rounded corners=2.6pt,%
    fill=#2,inner xsep=6.5pt,inner ysep=3pt]{%
    \popxboldls\fontsize{7}{7}\selectfont\color{#3}\MakeUppercase{#1}};%
}
\newcommand{\popsquiggle}[1]{%
  \tikz[baseline=-0.4ex]{
    \draw[#1,line width=2.6pt,line cap=round]
      plot [smooth] coordinates {(0,0.09) (0.34,0.01) (0.68,0.21) (1.02,0.01) (1.36,0.10)};
  }%
}
% Mot-clé surligné (marqueur orange sous le texte)
\newcommand{\cle}[1]{{\popxbold\color{popdark}#1}}

% ============================================================
% En-tête / pied
% ============================================================
\pagestyle{fancy}
\fancyhf{}
\renewcommand{\headrulewidth}{0pt}
\renewcommand{\footrulewidth}{0pt}
\lhead{\raisebox{-0.15ex}{\poplogo\fontsize{13}{13}\selectfont\color{popink}OnBuch\color{poporange}+}}
\rhead{\poppill{\DOCMATIERE\ \textbullet\ \DOCNIVEAU\ \textbullet\ Chapter \DOCLECON}{white}{popink}}
\lfoot{\popxbold\fontsize{7.6}{7.6}\selectfont\color{popmuted}\MakeUppercase{Signed course OnBuch+}\ \textbullet\ {\popsemi\DOCTITRE}}
\rfoot{\tikz[baseline=-0.6ex]\node[rounded corners=8.5pt,fill=popink,minimum height=17pt,minimum width=17pt,inner xsep=5pt,inner sep=0]{\popxbold\fontsize{8}{8}\selectfont\color{popcream}\thepage\,/\,\pageref*{LastPage}};}

% ============================================================
% Titres
% ============================================================
\newcommand{\chaptitre}[2]{%
  \begin{tcolorbox}[enhanced,colback=poporange,colframe=popink,boxrule=2.6pt,arc=11pt,
      drop shadow=popink,left=15pt,right=15pt,top=12pt,bottom=11pt,before skip=3pt,after skip=18pt]
    \poppill{#2}{popink}{popcream}\par\vspace{9pt}
    {\raggedright\hyphenpenalty=10000\exhyphenpenalty=10000\popdisplay\fontsize{22}{24}\selectfont\color{popcream}\MakeUppercase{#1}\par}
  \end{tcolorbox}
}
% Section numérotée : pastille orange avec le numéro + titre Archivo
\newcounter{popsec}
\newcommand{\coursec}[1]{%
  \stepcounter{popsec}\setcounter{popsub}{0}%
  \par\vspace{16pt}\noindent
  \begin{minipage}{\linewidth}
  \tikz[baseline=-0.7ex]\node[draw=popink,line width=1.6pt,fill=poporange,rounded corners=4pt,minimum size=22pt,inner sep=0]{\popdisplay\fontsize{12}{12}\selectfont\color{popcream}\thepopsec};%
  \hspace{8pt}{\popdisplay\fontsize{15}{17}\selectfont\color{popink}\MakeUppercase{#1}}\par
  \vspace{4pt}\noindent{\color{poporange}\rule{36pt}{3.6pt}}
  \end{minipage}\par\nopagebreak\vspace{8pt}
}
\newcounter{popsub}
\newcommand{\courssub}[1]{%
  \stepcounter{popsub}%
  \par\vspace{9pt}\noindent{\popxbold\fontsize{11.5}{13}\selectfont\color{popink}{\color{poporange}\thepopsec.\thepopsub}\hspace{6pt}#1}\par\nopagebreak\vspace{4pt}
}

% ============================================================
% Boîtes pédagogiques — même recette : fond blanc, bordure épaisse colorée,
% ombre dure encre, badge plein. Titre optionnel [#1].
% ============================================================
\tcbset{popbase/.style={breakable,enhanced,colback=white,boxrule=2.3pt,arc=9pt,
  shadow={3pt}{-3pt}{0pt}{popink},left=14pt,right=14pt,top=11pt,bottom=11pt,before skip=11pt,after skip=15pt,
  fontupper=\popsemi\fontsize{10.6}{14.5}\selectfont\color{popink},
  fontlower=\popsemi\fontsize{10.6}{14.5}\selectfont\color{popink}}}
% Coche dessinée (le glyphe ✓ n'existe pas dans les polices du document)
\newcommand{\popcheck}[1]{\tikz[baseline=-0.5ex]\draw[#1,line width=1.6pt,line cap=round,line join=round](-0.13,0.0)--(-0.04,-0.09)--(0.13,0.1);}
\providecommand{\coche}{\popcheck{popgreen}}
\providecommand{\resultat}[1]{\par\smallskip\noindent\fcolorbox{popgreen}{popgreenL}{\parbox{\dimexpr\linewidth-2\fboxsep-2\fboxrule\relax}{\textbf{Result.} #1}}\par\smallskip}
\newcommand{\popboxhead}[3]{\poppill{#1}{#2}{white}\ifx\relax#3\relax\else\hspace{6pt}{\popxbold\fontsize{10.6}{12}\selectfont\color{popink}#3}\fi\par\vspace{7pt}}

\newtcolorbox{aretenir}[1][]{popbase,colframe=popgreen,before upper={\popboxhead{Key points}{popgreen}{#1}}}
\newtcolorbox{exemplebox}[1][]{popbase,colframe=poporange,before upper={\popboxhead{Exemple}{poporange}{#1}}}
\newtcolorbox{definition}[1][]{popbase,colframe=popblue,before upper={\popboxhead{Definition}{popblue}{#1}}}
\newtcolorbox{propriete}[1][]{popbase,colframe=poppurple,before upper={\popboxhead{Property}{poppurple}{#1}}}
\newtcolorbox{methode}[1][]{popbase,colframe=popgold,before upper={\popboxhead{Method}{popgold}{#1}}}
\newtcolorbox{attention}[1][]{popbase,colframe=poppink,before upper={\popboxhead{Warning}{poppink}{#1}}}
\newtcolorbox{experience}[1][]{popbase,colframe=popblue,colback=popblueL,before upper={\popboxhead{Experiment}{popblue}{#1}}}
% « Le savais-tu ? » : carte violette pleine (contexte, culture, Cameroun)
\newtcolorbox{savaistu}[1][]{popbase,colback=poppurple,colframe=popink,
  fontupper=\popsemi\fontsize{10.4}{14.2}\selectfont\color{white},
  before upper={\poppill{Did you know?}{white}{poppurple}\ifx\relax#1\relax\else\hspace{6pt}{\popxbold\color{white}#1}\fi\par\vspace{7pt}}}
% Exercice résolu : énoncé en haut, \tcblower puis la solution sur fond crème
\newcounter{popexo}\newcounter{popexr}
\newtcolorbox{exoresolu}[1][]{popbase,colframe=poporange,colbacklower=popcream,
  segmentation style={popink,line width=1.2pt,dashed},
  before upper={\stepcounter{popexr}\popboxhead{Worked example \thepopexr}{poporange}{#1}},
  before lower={\poppill{Solution}{popink}{popcream}\par\vspace{6pt}}}
% Exercice (non résolu en place) — la correction est dans la partie Corrigés
\newtcolorbox{exercice}[1][]{popbase,colframe=popink,
  before upper={\stepcounter{popexo}\popboxhead{Exercise \thepopexo}{popink}{#1}}}
% Corrigé
\newtcolorbox{corrige}[1][]{popbase,colframe=popgreen,colback=popgreenL,
  before upper={\popboxhead{Solution}{popgreen}{#1}}}
% Objectifs / prérequis / bilan
\newtcolorbox{objectifs}{popbase,colframe=popink,colback=poporangeL,
  before upper={\popboxhead{Chapter objectives}{popink}{}}}
\newtcolorbox{prerequis}{popbase,colframe=popink,colback=popblueL,
  before upper={\popboxhead{Prerequisites}{popblue}{}}}
\newtcolorbox{bilan}{popbase,colframe=popink,colback=white,boxrule=2.8pt,
  before upper={\popboxhead{Fiche bilan}{poporange}{L'essentiel à connaître pour l'examen}}}
% Figure encadrée avec légende
\newtcolorbox{popfigure}[1][]{enhanced,breakable=false,colback=white,colframe=popline,boxrule=1.4pt,arc=8pt,
  left=10pt,right=10pt,top=10pt,bottom=8pt,before skip=10pt,after skip=12pt,halign=center,
  lower separated=false,halign lower=center,
  fontlower=\popsemi\fontsize{9}{11.5}\selectfont\color{popmuted},
  #1}
\newcommand{\legende}[1]{\par\vspace{6pt}{\popsemi\fontsize{9}{11.5}\selectfont\color{popmuted}#1}}

% ============================================================
% Signature OnBuch+
% ============================================================
\newcommand{\poplogoinline}[1]{{\poplogo\fontsize{#1}{#1}\selectfont\color{popink}OnBuch\color{poporange}+}}
\newcommand{\signatureonbuch}{%
  \begin{tcolorbox}[enhanced,colback=popink,colframe=popink,boxrule=2.2pt,arc=10pt,
      drop shadow=poporange,left=14pt,right=14pt,top=10pt,bottom=10pt,before skip=0pt,after skip=0pt]
    \tikz[baseline=-0.7ex]\node[circle,fill=popgreen,draw=popcream,line width=1.3pt,minimum size=22pt,inner sep=0]{\popcheck{white}};%
    \hspace{9pt}\begin{minipage}[c]{0.62\linewidth}
      {\popxbold\fontsize{12}{13}\selectfont\color{popcream}Signed\ \textbullet\ {\poplogo OnBuch\color{poporange}+}}\par\vspace{2pt}
      {\raggedright\popsemi\fontsize{8}{10.5}\selectfont\color{popline}\MakeUppercase{OnBuch+ teaching team}\\\MakeUppercase{Follows the official MINESEC syllabus}\par}
    \end{minipage}\hfill
    {\poplogo\fontsize{22}{22}\selectfont\color{popcream}OnBuch\color{poporange}+}
  \end{tcolorbox}%
}

% ============================================================
% Page de garde
% #1 titre · #2 matière · #3 niveau · #4 module · #5 n° leçon · #6 durée ·
% #7 description / objectifs (texte ou itemize)
% ============================================================
\newcommand{\pagegarde}[7]{%
  \thispagestyle{empty}
  \newgeometry{a4paper,margin=1.7cm,top=1.6cm,bottom=1.4cm}%
  \begin{tikzpicture}[remember picture,overlay]
    \fill[poporange!18] ([xshift=-3.2cm,yshift=-3.6cm]current page.north east) circle (4.6cm);
    \fill[poppurple!14] ([xshift=2.4cm,yshift=7.6cm]current page.south west) circle (3.8cm);
  \end{tikzpicture}%
  {\poplogo\fontsize{30}{30}\selectfont\color{popink}OnBuch\color{poporange}+}\hfill
  \raisebox{0.6ex}{\poppill{Signed course \textbullet\ Premium}{popink}{popcream}}\par
  \vspace{1pt}\popsquiggle{poppurple}\par
  \vspace{8pt}
  \begin{tcolorbox}[enhanced,colback=poporange,colframe=popink,boxrule=2.8pt,arc=13pt,
      drop shadow=popink,left=18pt,right=18pt,top=12pt,bottom=13pt,after skip=10pt]
    \poppill{#2 \textbullet\ #3}{popink}{popcream}\hspace{5pt}\poppill{#4}{white}{popink}\par\vspace{8pt}
    {\popdisplay\fontsize{14}{14}\selectfont\color{popink}CHAPTER #5}\par\vspace{4pt}
    {\raggedright\hyphenpenalty=10000\exhyphenpenalty=10000\popdisplay\fontsize{25}{27}\selectfont\color{popcream}\MakeUppercase{#1}\par}
    \vspace{8pt}
    {\popxbold\fontsize{9.5}{9.5}\selectfont\color{popink}\MakeUppercase{Suggested time: #6}}
  \end{tcolorbox}
  \begin{tcolorbox}[enhanced,colback=white,colframe=popink,boxrule=2.2pt,arc=10pt,
      drop shadow=popink,left=16pt,right=16pt,top=10pt,bottom=10pt,after skip=8pt]
    \poppill{In this chapter}{poppurple}{white}\par\vspace{5pt}
    \popsemi\fontsize{9.3}{12.6}\selectfont\color{popink}#7
  \end{tcolorbox}
  \noindent\begin{minipage}[t]{0.487\linewidth}
    \begin{tcolorbox}[enhanced,colback=popgreen,colframe=popink,boxrule=2.2pt,arc=10pt,
        drop shadow=popink,left=11pt,right=11pt,top=10pt,bottom=10pt,height=3.1cm,valign=center]
      \tcbox[colback=white,colframe=popink,boxrule=1.3pt,arc=5pt,boxsep=0pt,nobeforeafter,
        left=3pt,right=3pt,top=3pt,bottom=3pt,valign=center]{\qrcode[height=1.9cm]{https://play.google.com/store/apps/details?id=com.luvvix.onbuch&hl=en}}%
      \hspace{8pt}\begin{minipage}[c]{0.5\linewidth}
        {\popdisplay\fontsize{11}{12.5}\selectfont\color{white}\MakeUppercase{Download OnBuch}}\par\vspace{4pt}
        {\raggedright\popsemi\fontsize{8.4}{11}\selectfont\color{white}Free \textbullet\ Google Play\\AI tutor, past papers, results\par}
      \end{minipage}
    \end{tcolorbox}
  \end{minipage}\hfill
  \begin{minipage}[t]{0.487\linewidth}
    \begin{tcolorbox}[enhanced,colback=poppurple,colframe=popink,boxrule=2.2pt,arc=10pt,
        drop shadow=popink,left=11pt,right=11pt,top=10pt,bottom=10pt,height=3.1cm,valign=center]
      \tikz[baseline=-0.7ex]\node[circle,fill=white,draw=popink,line width=1.3pt,minimum size=1.6cm,inner sep=0]{\popdisplay\fontsize{24}{24}\selectfont\color{poppurple}+};%
      \hspace{8pt}\begin{minipage}[c]{0.6\linewidth}
        {\popdisplay\fontsize{11}{12.5}\selectfont\color{white}\MakeUppercase{Go OnBuch+}}\par\vspace{4pt}
        {\raggedright\popsemi\fontsize{8.4}{11}\selectfont\color{white}All signed courses, unlimited AI tutor, PDF sheets — OnBuch+ tab in the app\par}
      \end{minipage}
    \end{tcolorbox}
  \end{minipage}
  \vspace{8pt}
  \popcontenu
  \vfill
  \signatureonbuch
  \restoregeometry
  \newpage
}

% Rangée « Ce cours contient » (page de garde)
\newcommand{\poptile}[3]{% #1 couleur #2 gros texte #3 légende
  \begin{tcolorbox}[enhanced,colback=#1,colframe=popink,boxrule=2pt,arc=9pt,shadow={2.5pt}{-2.5pt}{0pt}{popink},
      left=6pt,right=6pt,top=9pt,bottom=9pt,halign=center,valign=center,height=1.75cm,width=\linewidth,nobeforeafter]
    {\popdisplay\fontsize{12.5}{13}\selectfont\color{white}\MakeUppercase{#2}}\par\vspace{3pt}
    {\centering\popsemi\fontsize{7.8}{9.5}\selectfont\color{white}#3\par}
  \end{tcolorbox}}
\newcommand{\popcontenu}{%
  \par\noindent{\popxboldls\fontsize{7.5}{7.5}\selectfont\color{popmuted}THIS SIGNED COURSE CONTAINS}\par\vspace{7pt}
  \noindent\begin{minipage}[t]{0.235\linewidth}\poptile{popblue}{Course}{complete, illustrated, step by step}\end{minipage}\hfill
  \begin{minipage}[t]{0.235\linewidth}\poptile{poporange}{Methods}{and worked examples}\end{minipage}\hfill
  \begin{minipage}[t]{0.235\linewidth}\poptile{poppink}{Exercises}{exam style + solutions}\end{minipage}\hfill
  \begin{minipage}[t]{0.235\linewidth}\poptile{popgreen}{Summary sheet}{the essentials to remember}\end{minipage}\par}

% Sommaire
\newtcolorbox{sommaire}{enhanced,colback=white,colframe=popink,boxrule=2.2pt,arc=10pt,
  drop shadow=popink,left=18pt,right=18pt,top=13pt,bottom=13pt,before skip=6pt,after skip=20pt,
  before upper={\poppill{Contents}{poppurple}{white}\par\vspace{9pt}\popsemi\fontsize{10.6}{14.5}\selectfont\color{popink}}}

% Signature de fin de cours — poussée tout en bas de la dernière page
\newcommand{\findecours}{%
  \par\vfill\nopagebreak
  \begin{center}\popsquiggle{poporange}\end{center}\vspace{4pt}
  \signatureonbuch
}
\AtBeginDocument{\def\llbracket{\mathopen{[\mkern-3.5mu[}}\def\rrbracket{\mathclose{]\mkern-3.5mu]}}} % intervalles d'entiers (glyphe ⟦ absent de la police)
% Glyphes absents de Fira Math : redéfinis à partir de glyphes disponibles
\AtBeginDocument{%
  \def\setminus{\mathbin{\backslash}}\let\smallsetminus\setminus
  \def\vdots{\mathord{\vbox{\baselineskip3.2pt\lineskiplimit0pt\kern4pt\hbox{.}\hbox{.}\hbox{.}}}}%
  \def\ddots{\mathinner{\mkern1mu\raise6.5pt\hbox{.}\mkern2mu\raise3.6pt\hbox{.}\mkern2mu\raise0.7pt\hbox{.}\mkern1mu}}%
  \def\triangle{\mathord{\scalebox{0.9}{$\symup{\Delta}$}}}%
  \def\nabla{\mathord{\raisebox{\depth}{\scalebox{1}[-1]{$\symup{\Delta}$}}}}%
  \def\checkmark{\popcheck{popdark}}%
  \def\star{\mathbin{\ast}}\def\bigstar{\mathord{\ast}}%
  \def\diamond{\mathbin{\scalebox{0.6}{$\Diamond$}}}%
  \def\bigcup{\mathop{\mathchoice{\raisebox{-0.3em}{\scalebox{1.7}{$\cup$}}}{\scalebox{1.25}{$\cup$}}{\cup}{\cup}}\displaylimits}%
  \def\bigcap{\mathop{\mathchoice{\raisebox{-0.3em}{\scalebox{1.7}{$\cap$}}}{\scalebox{1.25}{$\cap$}}{\cap}{\cap}}\displaylimits}%
  \def\lesseqqgtr{\mathrel{\lessgtr}}\def\bigoplus{\mathop{\oplus}\displaylimits}%
}
\providecommand{\ssi}{if and only if} % abréviation fréquente
\AtBeginDocument{\providecommand{\wideparen}{\overparen}} % arc de cercle

% Fonctions usuelles que les modèles utilisent sans que amsmath les fournisse
\providecommand{\arsinh}{\operatorname{arsinh}}\providecommand{\arcosh}{\operatorname{arcosh}}
\providecommand{\artanh}{\operatorname{artanh}}\providecommand{\arsech}{\operatorname{arsech}}
\providecommand{\arcsch}{\operatorname{arcsch}}\providecommand{\arcoth}{\operatorname{arcoth}}
\providecommand{\arcsinh}{\operatorname{arsinh}}\providecommand{\arccosh}{\operatorname{arcosh}}
\providecommand{\arctanh}{\operatorname{artanh}}\providecommand{\sech}{\operatorname{sech}}
\providecommand{\csch}{\operatorname{csch}}\providecommand{\arcsec}{\operatorname{arcsec}}
\providecommand{\arccsc}{\operatorname{arccsc}}\providecommand{\arccot}{\operatorname{arccot}}
\providecommand{\sgn}{\operatorname{sgn}}\providecommand{\lcm}{\operatorname{lcm}}
\providecommand{\Var}{\operatorname{Var}}\providecommand{\Cov}{\operatorname{Cov}}

%% FIGFIX-BEGIN v5 — correctifs globaux des figures (docs/onbuch-generation/tools/figfix)
\usepackage{adjustbox}\usepackage{etoolbox}\tcbuselibrary{raster}
% toute figure TikZ/pgfplots plus large que la ligne est réduite à la largeur disponible
\BeforeBeginEnvironment{tikzpicture}{\begin{adjustbox}{max width=\linewidth}}
\AfterEndEnvironment{tikzpicture}{\end{adjustbox}}
% légendes pgfplots lisibles par-dessus les courbes
\pgfplotsset{every axis/.append style={legend style={fill=white,fill opacity=0.92,text opacity=1,draw=black!15,inner sep=3pt}}}
% étiquettes d'arêtes (syntaxe edge["..."]) avec fond blanc
\tikzset{every edge quotes/.append style={fill=white,inner sep=1.2pt}}
%% FIGFIX-END
\begin{document}

\pagegarde{Special Discrete Probability Distributions}{Pure Mathematics with Statistics}{Upper Sixth}{Upper Sixth — Pure mathematics with statistics}{12}{20 periods}{This chapter studies the three special discrete distributions of the A-Level syllabus: the binomial, the geometric and the Poisson distributions. You will learn their characteristics, probability formulae, expectations, variances, recurrence relations, cumulative tables, and how to fit, approximate, compare and apply them to real-life problems.
\vspace{4pt}
\begin{multicols}{2}\small
\begin{itemize}
\item By the end of this chapter I can recognise a binomial situation and state its parameters n and p.
\item By the end of this chapter I can compute binomial probabilities using the formula P(X = r) = C(n,r) p\^{}r q\^{}(n-r).
\item By the end of this chapter I can find the expectation and variance of a binomial distribution and use recurrence formulae to generate probabilities.
\item By the end of this chapter I can use the cumulative distribution function and cumulative binomial probability tables.
\item By the end of this chapter I can fit a binomial distribution to given data and compare observed and expected frequencies.
\item By the end of this chapter I can model waiting-time situations with the geometric distribution and compute its expectation and variance.
\end{itemize}\end{multicols}}

\begin{sommaire}
\begin{multicols}{2}
\begin{enumerate}
\item The Binomial Distribution: Characteristics and Probabilities
\item Expectation, Variance and Recurrence Formulae for the Binomial
\item Cumulative Distribution Function and Binomial Tables
\item Fitting a Binomial Distribution to Data and Applications
\item The Geometric Distribution
\item The Poisson Distribution: Characteristics, Expectation and Variance
\item Random Events, Unit Interval, Recurrence Formula and Mode of the Poisson Distribution
\item Poisson Approximation to the Binomial and Applications of the Poisson Distribution
\item Comparing Distributions and Real-Life Problem Solving
\item Integration activity
\item Methods and worked examples
\item Exercises
\item Solutions to the exercises
\item Summary sheet
\end{enumerate}
\end{multicols}
\end{sommaire}

\begin{objectifs}
\begin{itemize}
\item By the end of this chapter I can recognise a binomial situation and state its parameters n and p.
\item By the end of this chapter I can compute binomial probabilities using the formula P(X = r) = C(n,r) p\^{}r q\^{}(n-r).
\item By the end of this chapter I can find the expectation and variance of a binomial distribution and use recurrence formulae to generate probabilities.
\item By the end of this chapter I can use the cumulative distribution function and cumulative binomial probability tables.
\item By the end of this chapter I can fit a binomial distribution to given data and compare observed and expected frequencies.
\item By the end of this chapter I can model waiting-time situations with the geometric distribution and compute its expectation and variance.
\item By the end of this chapter I can recognise Poisson situations, compute Poisson probabilities, and use its expectation, variance, recurrence formula and mode.
\item By the end of this chapter I can use the Poisson distribution as an approximation to the binomial distribution when appropriate.
\item By the end of this chapter I can compare distributions and solve real-life problems involving binomial, geometric and Poisson models.
\end{itemize}
\end{objectifs}

\begin{prerequis}
\begin{itemize}
\item Probability of events, independent events and the binomial expansion (Lower Sixth Pure Mathematics).
\item Discrete random variables, probability distributions, expectation E(X) and variance Var(X) (PMS chapter on random variables).
\item Combinations C(n,r) and factorial notation.
\item Sigma notation and manipulation of series, including the exponential series for e\^{}x.
\item Basic statistical measures: mean and variance of a frequency distribution.
\end{itemize}
\end{prerequis}

\newpage

\coursec{The Binomial Distribution: Characteristics and Probabilities}

At the Yaoundé central vaccination post, a nurse knows from records that $8\%$ of children develop a mild fever after a given vaccine. If $20$ children are vaccinated this morning, how many should she expect to react, and what is the chance that at most two do? Meanwhile, at the Mfoundi call centre, an average of $5$ calls arrive per minute: what is the probability of receiving exactly $3$ calls in the next minute? These two everyday situations look different, yet each is governed by a special discrete probability distribution — the \cle{binomial} and the \cle{Poisson} distributions. This chapter gives you the tools to model, compute and compare them, exactly as required at the GCE Advanced Level.

We begin with the binomial distribution, which answers the nurse's question: it counts the number of ``successes'' in a fixed number of repeated, identical, independent experiments.

\courssub{Bernoulli Trials: The Building Block}

\begin{experience}[One trial, two outcomes]
Toss a coin once. Either it lands Heads (call it \emph{success}) or Tails (\emph{failure}). Now imagine repeating the toss $10$ times and counting the Heads. Each single toss is a \cle{Bernoulli trial}; the total count of Heads is what we shall soon call a \cle{binomial random variable}. Notice three features of each toss: it has exactly two outcomes, the probability of Heads never changes, and one toss does not influence the next. These features are the heart of the whole chapter.
\end{experience}

\begin{definition}[Bernoulli trial]
A \cle{Bernoulli trial} is a single random experiment with exactly two possible outcomes, conventionally called \emph{success} ($S$) and \emph{failure} ($F$). We write
\[ P(S) = p, \qquad P(F) = q = 1 - p, \qquad 0 \le p \le 1. \]
\end{definition}

Everyday Bernoulli trials are everywhere:
\begin{itemize}
\item Tossing a coin: success $=$ Heads, $p = 0.5$.
\item A student sitting a pass/fail driving test in Buea: success $=$ pass.
\item Inspecting one bottle on a production line at a Douala brewery: success $=$ the bottle is defective, say $p = 0.02$.
\item The nurse's situation: one vaccinated child either develops a mild fever ($p = 0.08$) or does not ($q = 0.92$).
\end{itemize}

\begin{attention}[``Success'' need not be pleasant]
\emph{Success} is just a label for the outcome we are counting. Counting defective bottles or children with fever is perfectly legitimate: the ``success'' is the event of interest, however undesirable it may be in real life.
\end{attention}

\courssub{From Repeated Trials to the Binomial Distribution}

Suppose we repeat a Bernoulli trial a \emph{fixed} number $n$ of times, under \emph{identical} and \emph{independent} conditions, and we count the total number $X$ of successes. Then $X$ takes values in $\{0, 1, 2, \dots, n\}$ and is called a binomial random variable.

\begin{definition}[Binomial distribution]
A discrete random variable $X$ follows a \cle{binomial distribution} with parameters $n$ and $p$, written
\[ X \sim B(n, p), \]
if $X$ counts the number of successes in $n$ independent Bernoulli trials, each with constant probability $p$ of success. Here $n \in \mathbb{N}$ is the \emph{number of trials} and $p$ is the \emph{probability of success on each trial}; we write $q = 1 - p$.
\end{definition}

\begin{methode}[The four conditions to check before using $B(n,p)$]
Before modelling a situation by a binomial distribution, verify all four conditions (often summarised as \textbf{BINS}):
\begin{enumerate}
\item \textbf{B}inary: each trial has exactly two outcomes (success/failure).
\item \textbf{I}ndependent: the outcome of one trial does not affect the others.
\item \textbf{N}umber fixed: the number of trials $n$ is fixed in advance.
\item \textbf{S}ame probability: $p$ is constant from trial to trial.
\end{enumerate}
If any condition fails (e.g.\ sampling without replacement from a small population, or a probability that changes with time), the binomial model is not appropriate.
\end{methode}

\begin{exemplebox}[Testing the conditions]
A fair die is rolled $6$ times and $X$ counts the number of sixes obtained. Binary? Yes (six / not six). Independent? Yes. Fixed $n = 6$? Yes. Constant $p = \tfrac{1}{6}$? Yes. Hence $X \sim B\!\left(6, \tfrac{1}{6}\right)$.

Counter-example: drawing $5$ cards \emph{without replacement} from a deck and counting hearts — the probability of a heart changes after each draw, so the model is not binomial.
\end{exemplebox}

\courssub{Derivation of the Probability Formula}

Let us build the formula from a concrete case. Take $n = 3$ trials with $P(S) = p$. The full tree of outcomes is shown below.

\begin{popfigure}
\begin{tikzpicture}[grow=right, level distance=2.6cm,
level 1/.style={sibling distance=3.6cm},
level 2/.style={sibling distance=1.8cm},
level 3/.style={sibling distance=0.9cm},
every node/.style={font=\small},
edge from parent/.style={draw=popmuted}]
\node {Start}
  child { node {$F\ (q)$}
    child { node {$F$}
      child { node {$F$} edge from parent node[below, font=\scriptsize] {$q$} }
      child { node {$S$} edge from parent node[above, font=\scriptsize] {$p$} }
      edge from parent node[below, font=\scriptsize] {$q$} }
    child { node {$S$}
      child { node {$F$} edge from parent node[below, font=\scriptsize] {$q$} }
      child { node {$S$} edge from parent node[above, font=\scriptsize] {$p$} }
      edge from parent node[above, font=\scriptsize] {$p$} }
    edge from parent node[below] {$q$} }
  child { node {$S\ (p)$}
    child { node {$F$}
      child { node {$F$} edge from parent node[below, font=\scriptsize] {$q$} }
      child { node {$S$} edge from parent node[above, font=\scriptsize] {$p$} }
      edge from parent node[below, font=\scriptsize] {$q$} }
    child { node {$S$}
      child { node {$F$} edge from parent node[below, font=\scriptsize] {$q$} }
      child { node {$S$} edge from parent node[above, font=\scriptsize] {$p$} }
      edge from parent node[above, font=\scriptsize] {$p$} }
    edge from parent node[above] {$p$} };
\end{tikzpicture}
\legende{Tree diagram for $n = 3$ Bernoulli trials: the $8$ outcome paths.}
\end{popfigure}

Reading the $8$ paths and grouping them by the number of successes:
\begin{center}
\begin{tabular}{|c|c|c|}
\hline
\rowcolor{popblueL} Number of successes $r$ & Paths & Probability of each path \\
\hline
$0$ & $FFF$ & $q^3$ \\
\hline
$1$ & $SFF,\ FSF,\ FFS$ & $pq^2$ \\
\hline
$2$ & $SSF,\ SFS,\ FSS$ & $p^2q$ \\
\hline
$3$ & $SSS$ & $p^3$ \\
\hline
\end{tabular}
\end{center}

Because the trials are independent, each path with exactly $r$ successes and $n - r$ failures has probability $p^r q^{\,n-r}$ (multiplication rule). The number of such paths is the number of ways of choosing which $r$ of the $n$ trials are successes, namely $\dbinom{n}{r}$. Adding the probabilities of these mutually exclusive paths gives the general law.

\begin{propriete}[Binomial probability formula — proof examinable]
If $X \sim B(n, p)$ and $q = 1 - p$, then for $r = 0, 1, 2, \dots, n$,
\[ P(X = r) = \binom{n}{r} p^r q^{\,n-r} = \binom{n}{r} p^r (1-p)^{\,n-r}. \]
\emph{Proof.} Any sequence of $n$ trial results containing exactly $r$ successes and $n - r$ failures has probability $p^r q^{\,n-r}$ by independence. There are $\binom{n}{r}$ such sequences (choose the positions of the $r$ successes), and they are mutually exclusive, so their probabilities add. $\square$
\end{propriete}

A reassuring check: the probabilities must sum to $1$. Indeed, by the binomial theorem,
\[ \sum_{r=0}^{n} P(X = r) = \sum_{r=0}^{n} \binom{n}{r} p^r q^{\,n-r} = (q + p)^n = 1^n = 1. \]
So the formula defines a genuine probability distribution, and the numbers $\binom{n}{r} p^r q^{n-r}$ are precisely the terms of the expansion of $(q+p)^n$ — this is why the distribution is called \emph{binomial}.

\begin{aretenir}[Key facts]
\begin{itemize}
\item $X \sim B(n,p)$: $X$ counts successes in $n$ independent trials with constant $p$.
\item $P(X = r) = \dbinom{n}{r} p^r q^{\,n-r}$, with $q = 1 - p$.
\item $P(X = 0) = q^n$ and $P(X = n) = p^n$ (only one path each).
\item $\displaystyle\sum_{r=0}^{n} P(X=r) = (q+p)^n = 1$.
\end{itemize}
\end{aretenir}

\courssub{Direct Computation of Probabilities}

\begin{exoresolu}[Computing binomial probabilities by hand]
A fair coin is tossed $5$ times. Let $X$ be the number of Heads. Find (a) $P(X = 3)$; (b) $P(X = 0)$; (c) $P(X = 5)$; (d) verify the total is $1$.
\tcblower
\textbf{Solution.} Each toss is a Bernoulli trial with $p = 0.5$, $q = 0.5$, and the $n = 5$ tosses are independent, so $X \sim B(5, 0.5)$.

(a) $\displaystyle P(X = 3) = \binom{5}{3}(0.5)^3(0.5)^2 = 10 \times (0.5)^5 = \frac{10}{32} = 0.3125.$

(b) $\displaystyle P(X = 0) = q^5 = (0.5)^5 = \frac{1}{32} = 0.03125.$

(c) $\displaystyle P(X = 5) = p^5 = (0.5)^5 = 0.03125.$

(d) The full distribution: $P(X=r) = \binom{5}{r}/32$ gives $\dfrac{1, 5, 10, 10, 5, 1}{32}$, and $1+5+10+10+5+1 = 32$, so the total is $1$. $\blacksquare$
\end{exoresolu}

\begin{exoresolu}[The nurse's vaccination problem — first part]
At the Yaoundé vaccination post, $8\%$ of children develop a mild fever after the vaccine. Out of $20$ children vaccinated this morning, find the probability that (a) exactly $2$ react; (b) none reacts.
\tcblower
\textbf{Solution.} Assuming the children react independently, let $X$ be the number who develop a fever. Then $X \sim B(20, 0.08)$, with $q = 0.92$.

(a) $\displaystyle P(X = 2) = \binom{20}{2}(0.08)^2(0.92)^{18} = 190 \times 0.0064 \times (0.92)^{18}.$

Now $(0.92)^{18} \approx 0.2229$, so $P(X = 2) \approx 190 \times 0.0064 \times 0.2229 \approx 0.2711.$

(b) $\displaystyle P(X = 0) = (0.92)^{20} \approx 0.1887.$

So there is roughly a $19\%$ chance that no child reacts, and about a $27\%$ chance that exactly two do. $\blacksquare$
\end{exoresolu}

\courssub{Cumulative Distribution Function and Binomial Tables}

Most examination questions ask not for a single value but for an \emph{accumulated} probability. Since the events $\{X = r\}$ are mutually exclusive, we simply add:
\[ P(X \ge r) = \sum_{k=r}^{n} P(X = k), \qquad P(X \le r) = \sum_{k=0}^{r} P(X = k), \qquad P(a \le X \le b) = \sum_{k=a}^{b} P(X = k). \]

\begin{definition}[Cumulative distribution function (CDF)]
For a discrete random variable $X$, the \cle{cumulative distribution function} $F(x)$ is defined by
\[ F(x) = P(X \le x) = \sum_{k \le x} P(X = k). \]
For $X \sim B(n,p)$, $F(r) = \sum_{k=0}^{r} \binom{n}{k} p^k q^{n-k}$ for $r = 0,1,\dots,n$.
\end{definition}

A very useful complement trick avoids long sums:
\[ P(X \ge 1) = 1 - P(X = 0) = 1 - q^n. \]

\begin{attention}[Strict versus non-strict inequalities]
$P(X \ge 2) = 1 - P(X = 0) - P(X = 1)$, but $P(X > 2) = 1 - P(X=0) - P(X=1) - P(X=2)$. Read the question twice: ``at least'', ``at most'', ``more than'' and ``fewer than'' all mean different sums. Also, never forget the coefficient $\binom{n}{r}$, and never swap $p$ and $q$: the power of $p$ is the number of \emph{successes} $r$.
\end{attention}

\begin{exoresolu}[Compound probabilities]
A marksman hits a target with probability $0.7$ on each shot, independently. He fires $4$ shots. Find the probability that he hits (a) at least once; (b) at least $3$ times; (c) between $1$ and $2$ times inclusive.
\tcblower
\textbf{Solution.} Let $X$ be the number of hits. Then $X \sim B(4, 0.7)$, $q = 0.3$.

(a) $P(X \ge 1) = 1 - P(X = 0) = 1 - (0.3)^4 = 1 - 0.0081 = 0.9919.$

(b) $P(X \ge 3) = P(X=3) + P(X=4) = \dbinom{4}{3}(0.7)^3(0.3) + (0.7)^4$
$= 4 \times 0.343 \times 0.3 + 0.2401 = 0.4116 + 0.2401 = 0.6517.$

(c) $P(1 \le X \le 2) = \dbinom{4}{1}(0.7)(0.3)^3 + \dbinom{4}{2}(0.7)^2(0.3)^2$
$= 4 \times 0.7 \times 0.027 + 6 \times 0.49 \times 0.09 = 0.0756 + 0.2646 = 0.3402.$ $\blacksquare$
\end{exoresolu}

\begin{methode}[Using cumulative binomial probability tables]
When $n$ is large, summing many terms by hand is tedious. In examinations, \cle{cumulative binomial tables} giving $P(X \le r)$ for selected $n$ and $p$ are often provided.
\begin{enumerate}
\item Identify $n$ and $p$; locate the correct page/block for that $n$.
\item Find the column for your $p$ (or $q$ if the table uses $q$).
\item Read off the value opposite $r$ for $P(X \le r)$.
\item For $P(X \ge r)$ use $1 - P(X \le r-1)$; for $P(X = r)$ use $P(X \le r) - P(X \le r-1)$.
\end{enumerate}
If tables are not provided, you must compute using the formula or a calculator (if allowed).
\end{methode}

\begin{exemplebox}[Reading a cumulative table]
Suppose $X \sim B(10, 0.3)$ and the table gives $P(X \le 3) = 0.6496$ and $P(X \le 2) = 0.3828$. Then:
$P(X = 3) = 0.6496 - 0.3828 = 0.2668$,
$P(X \ge 4) = 1 - 0.6496 = 0.3504$,
$P(2 \le X \le 4) = P(X \le 4) - P(X \le 1)$.
\end{exemplebox}

\courssub{The Shape of the Binomial Distribution}

The value of $p$ controls the shape of the distribution. When $p = 0.5$ the distribution is symmetric about its centre; when $p < 0.5$ it is skewed to the right (a long tail of large values); when $p > 0.5$ it is skewed to the left. The bar charts below compare $B(10, 0.3)$ and $B(10, 0.5)$.

\begin{popfigure}
\begin{tikzpicture}
\begin{axis}[
    ybar, bar width=6pt,
    width=12cm, height=6.2cm,
    xlabel={$r$}, ylabel={$P(X = r)$},
    xmin=-0.7, xmax=10.7, ymin=0, ymax=0.32,
    xtick={0,1,2,3,4,5,6,7,8,9,10},
    legend style={at={(0.97,0.95)}, anchor=north east, font=\small},
    axis lines=left]
\addplot+[ybar, fill=poporange, draw=popdark] coordinates {
(0,0.0282) (1,0.1211) (2,0.2335) (3,0.2668) (4,0.2001) (5,0.1029)
(6,0.0368) (7,0.0090) (8,0.0014) (9,0.0001) (10,0.0000)};
\addlegendentry{$B(10,\,0.3)$}
\addplot+[ybar, fill=popblue, draw=popdark] coordinates {
(0,0.0010) (1,0.0098) (2,0.0439) (3,0.1172) (4,0.2051) (5,0.2461)
(6,0.2051) (7,0.1172) (8,0.0439) (9,0.0098) (10,0.0010)};
\addlegendentry{$B(10,\,0.5)$}
\end{axis}
\end{tikzpicture}
\legende{Effect of $p$ on the shape: $B(10,0.3)$ is right-skewed; $B(10,0.5)$ is symmetric.}
\end{popfigure}

Observe that for $p = 0.3$ the most likely value is $r = 3$ and the probabilities die away slowly to the right, while for $p = 0.5$ the bars are mirror images about $r = 5$. In general the distribution peaks near $r = np$, an observation we shall make precise when we study the expectation in the next section.

\begin{savaistu}[Exam tip — GCE Advanced Level]
Always \textbf{state the distribution before computing}: write ``Let $X$ be the number of children who react; then $X \sim B(20, 0.08)$.'' Examiners award a mark for correctly identifying the model and its parameters, and it protects you from mixing up $n$ and $p$ later. Quote the formula $P(X=r) = \binom{n}{r}p^r q^{n-r}$ once, then substitute.
\end{savaistu}

\begin{exercice}[Practice]
\begin{enumerate}
\item A fair die is thrown $8$ times and $X$ counts the sixes. State the distribution of $X$ and compute $P(X = 2)$ and $P(X \ge 1)$.
\item It is known that $30\%$ of seedlings in a nursery fail to survive transplanting. Out of $6$ transplanted seedlings, find the probability that (a) exactly $4$ survive; (b) at most $1$ survives.
\item For $X \sim B(5, 0.4)$, write out the full probability distribution and verify that the probabilities sum to $1$.
\item A multiple-choice test has $10$ questions, each with $4$ options. A student guesses all answers. Let $X$ be the number of correct answers. State the distribution of $X$ and find $P(X \le 2)$.
\end{enumerate}
\end{exercice}

\begin{corrige}[Exercise 1]
\textbf{1.} $X \sim B\!\left(8, \tfrac16\right)$. $P(X=2) = \binom{8}{2}\left(\tfrac16\right)^2\left(\tfrac56\right)^6 = 28 \times \tfrac{1}{36} \times 0.3349 \approx 0.2605$. $P(X \ge 1) = 1 - \left(\tfrac56\right)^8 \approx 1 - 0.2326 = 0.7674$.

\textbf{2.} Let $Y$ = number that survive, $Y \sim B(6, 0.7)$. (a) $P(Y = 4) = \binom{6}{4}(0.7)^4(0.3)^2 = 15 \times 0.2401 \times 0.09 = 0.3241$. (b) $P(Y \le 1) = (0.3)^6 + 6(0.7)(0.3)^5 = 0.000729 + 0.010206 = 0.0109$.

\textbf{3.} $P(X=r) = \binom{5}{r}(0.4)^r(0.6)^{5-r}$: $0.07776,\ 0.25920,\ 0.34560,\ 0.23040,\ 0.07680,\ 0.01024$. Sum: $0.07776 + 0.25920 + 0.34560 + 0.23040 + 0.07680 + 0.01024 = 1.00000$.

\textbf{4.} $X \sim B(10, 0.25)$. $P(X \le 2) = P(X=0)+P(X=1)+P(X=2) = (0.75)^{10} + 10(0.25)(0.75)^9 + 45(0.25)^2(0.75)^8 \approx 0.0563 + 0.1877 + 0.2816 = 0.5256$. $\blacksquare$
\end{corrige}

You can now recognise a binomial situation, state $X \sim B(n,p)$, and compute both single and compound probabilities. In the next section we shall measure the \emph{average} behaviour of this distribution — its expectation and variance — and discover a quick recurrence way of generating the probabilities.

\coursec{Expectation, Variance and Recurrence Formulae for the Binomial}

In the previous section we defined the binomial distribution $B(n,p)$ and learnt how to compute individual probabilities $P(X=r)=\binom{n}{r}p^{r}q^{n-r}$. Knowing the full probability distribution is excellent, but in practice an examiner, an epidemiologist or an engineer often wants quick answers to three questions: \emph{What value do we expect on average? How spread out are the results? Which outcome is the most likely?} This section answers all three: we derive the \cle{expectation} $E(X)=np$ and the \cle{variance} $\mathrm{Var}(X)=npq$, and we establish a powerful \cle{recurrence formula} linking consecutive probabilities, which leads us directly to the \cle{mode} of the distribution.

\courssub{Expectation of a Binomial Distribution}

\textbf{An intuitive situation.}\par
Suppose a nurse at a health centre in Bamenda vaccinates $n=200$ children in a week, and each child has a probability $p=0.05$ of developing a mild fever after the injection. How many cases of fever should she \emph{expect}? Common sense says: $5\%$ of $200$, that is $10$ children. Nothing more complicated is needed: the expected number of successes in $n$ independent trials is simply $n$ times the probability of success in one trial. This is exactly the content of the following property.

\begin{propriete}[Expectation of $B(n,p)$]
Let $X\sim B(n,p)$. Then
\[ E(X)=np. \]
\emph{Proof (examinable).} Write $X=X_1+X_2+\cdots+X_n$, where $X_i$ is the indicator of success on trial $i$: $X_i=1$ with probability $p$ and $X_i=0$ with probability $q=1-p$. Each $X_i$ is a Bernoulli variable, so $E(X_i)=1\times p+0\times q=p$. By linearity of expectation (which holds for any random variables, independent or not),
\[ E(X)=E(X_1)+E(X_2)+\cdots+E(X_n)=\underbrace{p+p+\cdots+p}_{n\text{ terms}}=np. \]
\end{propriete}

\begin{attention}[Do not confuse $E(X)$ with a possible value]
The expectation $np$ need not be an integer, and it need not be a value $X$ can actually take. If $X\sim B(12,0.4)$, then $E(X)=4.8$: we can never observe $4.8$ successes, but over many repetitions of the $12$ trials, the \emph{average} number of successes settles near $4.8$. The expectation is a long-run average, not a prediction of a single experiment.
\end{attention}

\courssub{Variance and Standard Deviation}

Knowing the average is not enough. Two distributions can share the same mean yet behave very differently: one may cluster tightly around the mean, the other may fluctuate wildly. The \cle{variance} measures this spread.

\textbf{Intuition first.}\par
For a single Bernoulli trial, the variance is easily computed:
\[ \mathrm{Var}(X_i)=E(X_i^2)-\big(E(X_i)\big)^2=p-p^2=p(1-p)=pq. \]
Notice that this is largest when $p=\tfrac12$ (maximum uncertainty: success and failure are equally likely) and smallest when $p$ is near $0$ or $1$ (the outcome is almost certain). Since the $n$ trials of a binomial experiment are \emph{independent}, variances add up, giving the following result.

\begin{propriete}[Variance of $B(n,p)$]
Let $X\sim B(n,p)$ with $q=1-p$. Then
\[ \mathrm{Var}(X)=npq \qquad\text{and}\qquad \sigma=\sqrt{npq}. \]
\emph{Proof (examinable).} With $X=X_1+\cdots+X_n$ as above, the $X_i$ are independent, so the variance of the sum is the sum of the variances:
\[ \mathrm{Var}(X)=\mathrm{Var}(X_1)+\cdots+\mathrm{Var}(X_n)=\underbrace{pq+pq+\cdots+pq}_{n\text{ terms}}=npq. \]
The standard deviation is by definition $\sigma=\sqrt{\mathrm{Var}(X)}=\sqrt{npq}$.
\end{propriete}

\begin{attention}[The classic error: $\mathrm{Var}(X)=np$]
Every year, candidates write $\mathrm{Var}(X)=np$. This is \textbf{wrong}: $np$ is the \emph{mean}. The variance is $\mathrm{Var}(X)=npq$ — never forget the factor $q$. A quick sanity check: if $p=1$ (certain success), the variance must be $0$ since nothing varies; $npq=n\times1\times0=0$ confirms this, while $np=n$ would absurdly claim variability. Also remember: the \emph{standard deviation} is $\sqrt{npq}$, not $npq$.
\end{attention}

\begin{aretenir}[Summary of key parameters]
For $X\sim B(n,p)$, $q=1-p$:
\[ E(X)=np,\qquad \mathrm{Var}(X)=npq,\qquad \sigma=\sqrt{npq}. \]
These three formulae must be known by heart for the GCE Advanced Level examination.
\end{aretenir}

\begin{exoresolu}[Vaccination reactions]
A vaccine produces a mild side-effect in $8\%$ of patients. In a clinic, $150$ patients are vaccinated independently. Let $X$ be the number of patients who develop the side-effect. Find the mean, variance and standard deviation of $X$.
\tcblower
\textbf{Solution.} Each patient is a Bernoulli trial with $p=0.08$, and the trials are independent, so $X\sim B(150,\,0.08)$ with $q=0.92$.
\begin{align*}
E(X)&=np=150\times0.08=12,\\
\mathrm{Var}(X)&=npq=150\times0.08\times0.92=11.04,\\
\sigma&=\sqrt{11.04}\approx3.32.
\end{align*}
The clinic should expect about $12$ cases, with typical fluctuations of roughly $\pm3$ cases around this value.
\end{exoresolu}

\begin{exoresolu}[Guessing a multiple-choice quiz]
A GCE candidate guesses all $20$ questions of a multiple-choice paper, each question having $4$ options of which exactly one is correct. Let $X$ be the number of correct answers. Find $E(X)$, $\mathrm{Var}(X)$ and $\sigma$, and comment.
\tcblower
\textbf{Solution.} Guessing gives $p=\tfrac14=0.25$, $q=0.75$, $n=20$, so $X\sim B(20,0.25)$.
\begin{align*}
E(X)&=20\times0.25=5,\\
\mathrm{Var}(X)&=20\times0.25\times0.75=3.75,\\
\sigma&=\sqrt{3.75}\approx1.94.
\end{align*}
A pure guesser scores on average $5$ out of $20$, and will rarely deviate from $5$ by more than about $2$ correct answers. Guessing is not a strategy!
\end{exoresolu}

\courssub{The Recurrence Formula}

\textbf{The problem.}\par
Computing $P(X=r)=\binom{n}{r}p^{r}q^{n-r}$ for $r=0,1,2,\dots,n$ one by one is painful: for $n=12$ you would evaluate $13$ different binomial coefficients, each with powers of $p$ and $q$. Under examination conditions this is slow and error-prone. Fortunately, consecutive probabilities are linked by a remarkably simple ratio.

\begin{propriete}[Recurrence formula for the binomial]
Let $X\sim B(n,p)$, $q=1-p$, and write $P(r)=P(X=r)$. Then for $r=0,1,\dots,n-1$,
\[ \frac{P(X=r+1)}{P(X=r)}=\frac{(n-r)\,p}{(r+1)\,q},
\qquad\text{so}\qquad
P(X=r+1)=\frac{(n-r)\,p}{(r+1)\,q}\,P(X=r). \]
\emph{Proof (examinable and instructive).} Take the ratio of two consecutive terms:
\begin{align*}
\frac{P(X=r+1)}{P(X=r)}
&=\frac{\dbinom{n}{r+1}p^{r+1}q^{n-r-1}}{\dbinom{n}{r}p^{r}q^{n-r}}\\[2pt]
&=\frac{n!}{(r+1)!\,(n-r-1)!}\times\frac{r!\,(n-r)!}{n!}\times\frac{p^{r+1}}{p^{r}}\times\frac{q^{n-r-1}}{q^{n-r}}\\[2pt]
&=\frac{(n-r)}{(r+1)}\times p\times\frac{1}{q}
=\frac{(n-r)\,p}{(r+1)\,q}.
\end{align*}
Here we used $\dfrac{r!}{(r+1)!}=\dfrac{1}{r+1}$ and $\dfrac{(n-r)!}{(n-r-1)!}=n-r$.
\end{propriete}

\begin{exemplebox}[Remark: why this formula is a gift in exams]
The recurrence formula means you compute \textbf{only one} binomial probability from scratch — usually the easiest, $P(X=0)=q^{n}$ — and then generate the whole table by repeated multiplication. No more large coefficients $\binom{n}{r}$, no more fresh powers of $p$ and $q$ at each step. As a bonus, any arithmetic slip shows up immediately because the final probabilities must sum to $1$: this gives you a free checking mechanism.
\end{exemplebox}

\begin{methode}[Building a full probability table with the recurrence formula]
\begin{enumerate}
\item Compute the starting value $P(0)=q^{n}$.
\item For each $r=0,1,2,\dots,n-1$, compute the multiplier $m_r=\dfrac{(n-r)p}{(r+1)q}$.
\item Set $P(r+1)=m_r\times P(r)$, keeping at least $4$ decimal places throughout.
\item Stop at $P(n)$, then \textbf{check}: $\sum_{r=0}^{n}P(r)$ must equal $1$ (up to rounding).
\end{enumerate}
\end{methode}

\begin{exoresolu}[Full table of $B(3,0.4)$ by recurrence]
Let $X\sim B(3,0.4)$. Using the recurrence formula, construct the complete probability distribution of $X$.
\tcblower
\textbf{Solution.} Here $n=3$, $p=0.4$, $q=0.6$.
\textbf{Step 1:} $P(0)=q^{3}=0.6^{3}=0.216$.
\textbf{Step 2:} apply $P(r+1)=\dfrac{(3-r)\times0.4}{(r+1)\times0.6}P(r)$:
\begin{align*}
P(1)&=\frac{3\times0.4}{1\times0.6}\times0.216=2\times0.216=0.432,\\
P(2)&=\frac{2\times0.4}{2\times0.6}\times0.432=\frac{2}{3}\times0.432=0.288,\\
P(3)&=\frac{1\times0.4}{3\times0.6}\times0.288=\frac{2}{9}\times0.288=0.064.
\end{align*}
\textbf{Check:} $0.216+0.432+0.288+0.064=1.000$. \checkmark\ The distribution is
\begin{center}
\begin{tabular}{|l|c|c|c|c|}
\hline
\rowcolor{popblueL} $r$ & $0$ & $1$ & $2$ & $3$ \\
\hline
$P(X=r)$ & $0.216$ & $0.432$ & $0.288$ & $0.064$ \\
\hline
\end{tabular}
\end{center}
\end{exoresolu}

\courssub{Locating the Mode with the Recurrence Ratio}

The \cle{mode} of a discrete distribution is the value of $r$ with the greatest probability — the \emph{most likely} outcome. The recurrence formula gives an elegant way to find it without computing any probability at all.

\textbf{The key idea.}\par
The probabilities $P(0),P(1),P(2),\dots$ first increase, reach a peak (the mode), then decrease. The sequence is \emph{increasing} exactly when the next term is at least the current one, i.e. when the recurrence ratio is at least $1$:
\[ \frac{P(X=r+1)}{P(X=r)}=\frac{(n-r)p}{(r+1)q}\geq1. \]
Let us solve this inequality:
\begin{align*}
(n-r)p\geq(r+1)q
&\iff np-rp\geq rq+q\\
&\iff np-q\geq r(p+q)\\
&\iff r\leq np-q=(n+1)p-1 \quad(\text{since }p+q=1).
\end{align*}
So the probabilities \emph{increase} while $r\leq(n+1)p-1$ and \emph{decrease} afterwards. The peak — the mode — is therefore the first integer $r$ for which the increase stops.

\begin{methode}[Finding the mode of $B(n,p)$]
\begin{enumerate}
\item Compute $(n+1)p$.
\item If $(n+1)p$ is \textbf{not an integer}, the mode is the unique integer
\[ m=\lfloor (n+1)p\rfloor \quad(\text{the integer part of }(n+1)p). \]
\item If $(n+1)p$ \textbf{is an integer} $k$, there are \textbf{two equal modes}: $r=k-1$ and $r=k$, because then $P(k)/P(k-1)=1$.
\end{enumerate}
\end{methode}

\begin{exemplebox}[Mode of $B(12,0.4)$]
Here $(n+1)p=13\times0.4=5.2$, not an integer, so the mode is $m=\lfloor5.2\rfloor=5$: the most likely number of successes is $5$. Note the mean is $np=4.8$: mode and mean are close but not identical — the distribution is slightly skewed, as the bar chart below shows.
\end{exemplebox}

\begin{popfigure}
\begin{center}
\begin{tikzpicture}
\begin{axis}[
    ybar, bar width=8pt,
    width=13cm, height=7cm,
    xlabel={$r$ (number of successes)}, ylabel={$P(X=r)$},
    xmin=-0.7, xmax=12.7, ymin=0, ymax=0.26,
    xtick={0,1,2,3,4,5,6,7,8,9,10,11,12},
    ytick={0,0.05,0.10,0.15,0.20,0.25},
    axis lines=left,
    tick label style={font=\small},
    label style={font=\small},
]
\addplot+[ybar, fill=popblue, draw=popdark] coordinates {
(0,0.0022) (1,0.0174) (2,0.0639) (3,0.1419) (4,0.2128)
(6,0.1766) (7,0.1009) (8,0.0420) (9,0.0125) (10,0.0025)
(11,0.0003) (12,0.0000)};
\addplot+[ybar, fill=poporange, draw=popdark] coordinates {(5,0.2270)};
\draw[popgreen, very thick, dashed] (axis cs:4.8,0) -- (axis cs:4.8,0.245);
\node[popgreen, font=\small, anchor=west] at (axis cs:4.9,0.24) {mean $np=4.8$};
\node[popdark, font=\small, anchor=south] at (axis cs:5,0.2270) {mode $=5$};
\end{axis}
\end{tikzpicture}
\end{center}
\legende{Bar chart of $B(12,0.4)$. The orange bar marks the mode $r=5=\lfloor(n+1)p\rfloor$; the dashed green line marks the mean $E(X)=np=4.8$.}
\end{popfigure}

\begin{exoresolu}[Mode and most likely outcome]
A fair coin is tossed $10$ times and $X$ counts the number of heads. Using the recurrence ratio, find the mode of $X$ and verify it directly.
\tcblower
\textbf{Solution.} $X\sim B(10,0.5)$, so $(n+1)p=11\times0.5=5.5$, not an integer: the mode is $m=\lfloor5.5\rfloor=5$.
\emph{Verification via the ratio:} with $p=q=0.5$, the ratio is $\dfrac{P(r+1)}{P(r)}=\dfrac{10-r}{r+1}$. This is $\geq1$ iff $10-r\geq r+1$, i.e. $r\leq4.5$, i.e. $r\leq4$. So $P(0)<P(1)<\cdots<P(5)$ and $P(5)>P(6)>\cdots>P(10)$: the peak is at $r=5$, with $P(X=5)=\binom{10}{5}(0.5)^{10}=\dfrac{252}{1024}\approx0.246$. \checkmark
\end{exoresolu}

\begin{attention}[Mode versus mean]
Do not write ``the mode is $np$''. The mean is $np$; the mode is $\lfloor(n+1)p\rfloor$ (with a double mode when $(n+1)p$ is an integer). For $B(12,0.4)$ these are $4.8$ and $5$ respectively — close, but conceptually different quantities, and examiners test the distinction.
\end{attention}

\begin{exercice}[Practice]
\begin{enumerate}
\item A machine produces defective items with probability $0.02$. A batch of $500$ items is inspected. Let $X$ be the number of defectives. Find $E(X)$, $\mathrm{Var}(X)$ and $\sigma$.
\item For $X\sim B(8,0.3)$, use the recurrence formula starting from $P(0)=0.7^{8}$ to compute $P(1)$ and $P(2)$, giving $4$ decimal places.
\item Find the mode of $B(9,0.5)$. What do you notice?
\end{enumerate}
\end{exercice}

\begin{corrige}[Exercise 1]
\begin{enumerate}
\item $X\sim B(500,0.02)$: $E(X)=np=10$; $\mathrm{Var}(X)=npq=500\times0.02\times0.98=9.8$; $\sigma=\sqrt{9.8}\approx3.13$.
\item $P(0)=0.7^{8}=0.0576$. Ratio for $r=0$: $\dfrac{8\times0.3}{1\times0.7}=\dfrac{24}{7}$, so $P(1)=\dfrac{24}{7}\times0.0576\approx0.1977$. Ratio for $r=1$: $\dfrac{7\times0.3}{2\times0.7}=1.5$, so $P(2)=1.5\times0.1977\approx0.2965$.
\item $(n+1)p=10\times0.5=5$, an integer: there are \textbf{two modes}, $r=4$ and $r=5$, with equal probabilities $P(4)=P(5)\approx0.2461$.
\end{enumerate}
\end{corrige}

\begin{aretenir}[What you must remember from this section]
\begin{itemize}
\item For $X\sim B(n,p)$: $E(X)=np$, $\mathrm{Var}(X)=npq$, $\sigma=\sqrt{npq}$ — the factor $q$ in the variance is the most frequently forgotten detail in the exam.
\item The recurrence formula $\dfrac{P(X=r+1)}{P(X=r)}=\dfrac{(n-r)p}{(r+1)q}$ generates the whole distribution from $P(0)=q^{n}$ and provides a built-in check ($\sum P(r)=1$).
\item The probabilities increase while $r\leq(n+1)p-1$; hence the mode is $\lfloor(n+1)p\rfloor$, with a double mode $k-1$ and $k$ when $(n+1)p=k$ is an integer.
\end{itemize}
\end{aretenir}

These parameters and tools will be essential in the next section, where we move from individual probabilities to \emph{cumulative} probabilities $P(X\leq r)$ and learn to read the cumulative binomial tables provided in the examination.

\coursec{Cumulative Distribution Function and Binomial Tables}

In the previous sections, we mastered the calculation of individual probabilities $P(X = r)$ for a binomial random variable $X \sim B(n, p)$, using the formula, the recurrence relation, and the properties of expectation and variance. However, in many practical situations — quality control in a factory in Douala, medical screening in Yaoundé, or crop-yield estimates in the Northwest Region — we are not interested in the probability of \emph{exactly} $r$ successes, but rather in the probability of \emph{at most} $r$ successes, or \emph{at least} $r$ successes. This calls for the \cle{cumulative distribution function} (CDF) and the use of standard statistical tables, which are provided in the GCE examination.

\courssub{The Cumulative Distribution Function (CDF)}

\begin{definition}[Cumulative Distribution Function of a Discrete Random Variable]
Let $X$ be a discrete random variable taking values $x_1, x_2, \dots$ with probabilities $P(X = x_i)$. The \cle{cumulative distribution function}, denoted by $F(x)$, is defined for all real $x$ as:
\[
F(x) = P(X \le x) = \sum_{x_i \le x} P(X = x_i).
\]
For a binomial variable $X \sim B(n, p)$, the CDF at an integer $r$ ($0 \le r \le n$) is:
\[
F(r) = P(X \le r) = \sum_{k=0}^{r} \binom{n}{k} p^k (1-p)^{n-k}.
\]
By convention, $F(x) = 0$ for $x < 0$ and $F(x) = 1$ for $x \ge n$.
\end{definition}

\begin{aretenir}[Properties of the CDF for a Discrete Variable]
\begin{itemize}
\item $F(x)$ is a non-decreasing step function, with jumps exactly at the values taken by $X$.
\item $F(x)$ is right-continuous: the value at a jump point $r$ is the value \emph{after} the jump, i.e., $F(r) = P(X \le r)$.
\item $F(x) = 0$ for all $x$ below the smallest value of $X$, and $F(x) = 1$ for all $x$ at or above the largest value of $X$.
\item $P(X = r) = F(r) - F(r-1)$ for integer $r$.
\item $P(a < X \le b) = F(b) - F(a)$, and $P(a \le X \le b) = F(b) - F(a-1)$.
\end{itemize}
\end{aretenir}

\begin{exemplebox}[Building a Cumulative Table for $B(5, 0.4)$]
Let $X \sim B(5, 0.4)$. We first compute the individual probabilities $P(X=r)$ (using the formula $P(X=r)=\binom{5}{r}(0.4)^r(0.6)^{5-r}$ or the recurrence relation), then accumulate them row by row.

\begin{center}
\begin{tabular}{|c|cccccc|}
\hline
$r$ & 0 & 1 & 2 & 3 & 4 & 5 \\
\hline
$P(X=r)$ & 0.0778 & 0.2592 & 0.3456 & 0.2304 & 0.0768 & 0.0102 \\
\hline
$F(r)=P(X\le r)$ & 0.0778 & 0.3370 & 0.6826 & 0.9130 & 0.9898 & 1.0000 \\
\hline
\end{tabular}
\end{center}
For instance, $F(2) = 0.0778 + 0.2592 + 0.3456 = 0.6826$.
\textbf{Check:} the final cumulative probability must equal $1$ (subject to rounding).
\end{exemplebox}

\begin{popfigure}
\begin{tikzpicture}
\begin{axis}[
    width=0.9\linewidth,
    height=7cm,
    axis lines=middle,
    xlabel=$r$, ylabel=$F(r)$,
    xmin=-0.5, xmax=8.5,
    ymin=-0.05, ymax=1.05,
    xtick={0,1,2,3,4,5,6,7,8},
    ytick={0,0.2,0.4,0.6,0.8,1.0},
    grid=major,
    grid style={popline},
    tick label style={font=\small},
    label style={font=\small},
    clip=false
]
% Step function for B(8, 0.25)
\draw[popblue, very thick] (-0.5, 0) -- (0, 0);
\draw[popblue, very thick] (0, 0.1001) -- (1, 0.1001);
\draw[popblue, very thick] (1, 0.3671) -- (2, 0.3671);
\draw[popblue, very thick] (2, 0.6785) -- (3, 0.6785);
\draw[popblue, very thick] (3, 0.8862) -- (4, 0.8862);
\draw[popblue, very thick] (4, 0.9727) -- (5, 0.9727);
\draw[popblue, very thick] (5, 0.9958) -- (6, 0.9958);
\draw[popblue, very thick] (6, 0.9996) -- (7, 0.9996);
\draw[popblue, very thick] (7, 1.0000) -- (8, 1.0000);
\draw[popblue, very thick] (8, 1.0000) -- (8.5, 1.0000);

% Vertical jumps (dashed)
\draw[popblue, dashed, thin] (0, 0) -- (0, 0.1001);
\draw[popblue, dashed, thin] (1, 0.1001) -- (1, 0.3671);
\draw[popblue, dashed, thin] (2, 0.3671) -- (2, 0.6785);
\draw[popblue, dashed, thin] (3, 0.6785) -- (3, 0.8862);
\draw[popblue, dashed, thin] (4, 0.8862) -- (4, 0.9727);
\draw[popblue, dashed, thin] (5, 0.9727) -- (5, 0.9958);
\draw[popblue, dashed, thin] (6, 0.9958) -- (6, 0.9996);
\draw[popblue, dashed, thin] (7, 0.9996) -- (7, 1.0000);

% Filled circles (included)
\fill[popblue] (axis cs:0, 0.1001) circle (2pt);
\fill[popblue] (axis cs:1, 0.3671) circle (2pt);
\fill[popblue] (axis cs:2, 0.6785) circle (2pt);
\fill[popblue] (axis cs:3, 0.8862) circle (2pt);
\fill[popblue] (axis cs:4, 0.9727) circle (2pt);
\fill[popblue] (axis cs:5, 0.9958) circle (2pt);
\fill[popblue] (axis cs:6, 0.9996) circle (2pt);
\fill[popblue] (axis cs:7, 1.0000) circle (2pt);
\fill[popblue] (axis cs:8, 1.0000) circle (2pt);

% Open circles (excluded)
\draw[popblue, fill=white] (axis cs:0, 0) circle (2pt);
\draw[popblue, fill=white] (axis cs:1, 0.1001) circle (2pt);
\draw[popblue, fill=white] (axis cs:2, 0.3671) circle (2pt);
\draw[popblue, fill=white] (axis cs:3, 0.6785) circle (2pt);
\draw[popblue, fill=white] (axis cs:4, 0.8862) circle (2pt);
\draw[popblue, fill=white] (axis cs:5, 0.9727) circle (2pt);
\draw[popblue, fill=white] (axis cs:6, 0.9958) circle (2pt);
\draw[popblue, fill=white] (axis cs:7, 0.9996) circle (2pt);
\end{axis}
\end{tikzpicture}
\legende{Cumulative distribution function $F(r) = P(X \le r)$ for $X \sim B(8, 0.25)$. The function jumps at each integer value $r$; the solid dot indicates the value actually taken by $F$ at the jump.}
\end{popfigure}

\courssub{Using Cumulative Binomial Probability Tables}

In the GCE examination, you are provided with \cle{cumulative binomial probability tables}. These tables give $P(X \le r)$ directly for various values of $n$ (usually up to 50), $p$ (usually from 0.05 to 0.50 in steps of 0.05), and $r$.

\begin{aretenir}[Structure of the Tables]
The tables are organised by $n$ (number of trials). For each $n$, there are columns for different values of $p$. The rows correspond to $r = 0, 1, 2, \dots, n$. The entry in the cell $(r, p)$ is $P(X \le r)$ for $X \sim B(n, p)$.
\end{aretenir}

\begin{popfigure}
\begin{tikzpicture}
% Annotated table extract for n=10, p=0.2
\begin{scope}[scale=0.85, transform shape]
\draw[popdark, thick] (0,0) grid [xstep=2.5cm, ystep=0.8cm] (5,7.2);
% Header
\fill[popblueL] (0,6.4) rectangle (5,7.2);
\node[font=\bfseries\small, text=popdark] at (1.25, 6.8) {$n = 10$};
\node[font=\bfseries\small, text=popdark] at (3.75, 6.8) {$p = 0.20$};
% Column headers
\fill[popblueL] (0,5.6) rectangle (2.5,6.4);
\node[font=\bfseries\small] at (1.25, 6.0) {$r$};
\fill[popblueL] (2.5,5.6) rectangle (5,6.4);
\node[font=\bfseries\small] at (3.75, 6.0) {$P(X \le r)$};
% Rows
\foreach \r/\val/\y in {
    0/0.1074/5.6,
    1/0.3758/4.8,
    2/0.6778/4.0,
    3/0.8791/3.2,
    4/0.9672/2.4,
    5/0.9936/1.6,
    6/0.9991/0.8,
    7/0.9999/0.0
} {
    \node[font=\small] at (1.25, \y+0.4) {\r};
    \node[font=\small] at (3.75, \y+0.4) {\val};
}
% Highlight row r=2
\draw[poporange, very thick] (0, 4.0) rectangle (5, 4.8);
\node[poporange, font=\bfseries\small, anchor=west] at (5.2, 4.4) {$\leftarrow P(X \le 2) = 0.6778$};
% Title
\node[popdark, font=\bfseries\small, anchor=south] at (2.5, 7.4) {Extract from Cumulative Binomial Table};
\end{scope}
\end{tikzpicture}
\legende{Reading the table: for $n=10$, $p=0.20$, locate the block for $n=10$, the column $p=0.20$, and the row $r=2$. The entry is $0.6778$, so $P(X \le 2) = 0.6778$.}
\end{popfigure}

\courssub{Three Essential Table-Reading Rules}

Since the tables only give $P(X \le r)$, we must derive all other probabilities from this single quantity.

\begin{methode}[From Cumulative Tables to Any Probability]
Let $X \sim B(n, p)$ and let $F(r) = P(X \le r)$ be the value read from the table.
\begin{enumerate}
    \item \textbf{Exactly $r$ successes:} $P(X = r) = F(r) - F(r-1)$.
    \item \textbf{At least $r$ successes:} $P(X \ge r) = 1 - P(X \le r-1) = 1 - F(r-1)$.
    \item \textbf{More than $r$ successes:} $P(X > r) = 1 - P(X \le r) = 1 - F(r)$.
\end{enumerate}
\end{methode}

\begin{attention}[The Classic Trap: $P(X \ge r)$ versus $P(X > r)$]
A very common error is to write $P(X \ge r) = 1 - P(X \le r)$.
\emph{This is wrong.} The event $X \ge r$ includes $X = r$, so its complement is $X \le r-1$, not $X \le r$:
\[
P(X \ge r) = 1 - F(r-1) \qquad \text{and NOT} \qquad 1 - F(r).
\]
Similarly, $P(X > r) = 1 - F(r)$ because the event $X > r$ excludes $r$. Always ask yourself: \emph{is $r$ itself included or excluded?}
\end{attention}

\begin{exoresolu}[Using Tables for $X \sim B(15, 0.3)$]
Using the cumulative binomial tables for $n=15$, $p=0.30$, find:
\begin{enumerate}
    \item $P(X \le 3)$
    \item $P(X = 3)$
    \item $P(X \ge 3)$
    \item $P(X > 3)$
\end{enumerate}
\tcblower
\textbf{Solution.}
From the table ($n=15$, $p=0.30$):
\[
F(2) = P(X \le 2) = 0.1268, \quad F(3) = P(X \le 3) = 0.2969, \quad F(4) = P(X \le 4) = 0.5155.
\]
\begin{enumerate}
    \item \textbf{Direct reading:} $P(X \le 3) = F(3) = \mathbf{0.2969}$.
    \item \textbf{Subtraction:} $P(X = 3) = F(3) - F(2) = 0.2969 - 0.1268 = \mathbf{0.1701}$.
    \item \textbf{At least 3:} $P(X \ge 3) = 1 - F(2) = 1 - 0.1268 = \mathbf{0.8732}$.
    \item \textbf{More than 3:} $P(X > 3) = 1 - F(3) = 1 - 0.2969 = \mathbf{0.7031}$.
\end{enumerate}
\textbf{Check:} $P(X \ge 3) = P(X=3) + P(X>3) = 0.1701 + 0.7031 = 0.8732$. Consistent.
\end{exoresolu}

\courssub{Inverse Reading: Finding Critical Values}

Sometimes we are given a probability $\alpha$ (e.g., $0.95$ or $0.99$) and asked to find the smallest integer $r$ such that $P(X \le r) \ge \alpha$. This is finding a \cle{critical value} (or percentile). This skill is essential for hypothesis testing (covered later) and for quality control.

\begin{methode}[Finding the Smallest $r$ such that $P(X \le r) \ge \alpha$]
\begin{enumerate}
    \item Locate the correct block for $n$ and the column for $p$.
    \item Scan down the column of cumulative probabilities $F(r)$.
    \item Find the \textbf{first} row where the value is greater than or equal to $\alpha$.
    \item The corresponding value of $r$ is the answer.
\end{enumerate}
\end{methode}

\begin{exemplebox}[Finding a Critical Value for $B(20, 0.25)$]
Find the smallest integer $r$ such that $P(X \le r) \ge 0.90$ for $X \sim B(20, 0.25)$.

\textbf{Solution.}
Scan the column $n=20$, $p=0.25$:
\[
F(6) = 0.7858 \ (\text{too small}), \qquad F(7) = 0.8982 \ (\text{too small}), \qquad F(8) = 0.9591 \ (\ge 0.90).
\]
The smallest such $r$ is $\mathbf{8}$.
\end{exemplebox}

\courssub{Handling $p > 0.5$: The Symmetry Trick}

Standard binomial tables usually only provide values for $p \le 0.5$. If $p > 0.5$, we exploit the symmetry of the binomial distribution by counting \emph{failures} instead of successes.

\begin{propriete}[Symmetry of the Binomial Distribution]
If $X \sim B(n, p)$, then the number of failures $Y = n - X$ follows $B(n, 1-p)$, and
\[
P(X \le r) = P(n - Y \le r) = P(Y \ge n - r) = 1 - P(Y \le n - r - 1).
\]
Since $1 - p < 0.5$ when $p > 0.5$, we can look up $Y \sim B(n, 1-p)$ in the tables.
\end{propriete}

\begin{exoresolu}[Using Tables for $p > 0.5$: $X \sim B(10, 0.7)$]
Find $P(X \le 8)$ for $X \sim B(10, 0.7)$.
\tcblower
\textbf{Solution.}
Here $p = 0.7 > 0.5$. Define $Y = \text{number of failures} = 10 - X$. Then $Y \sim B(10, 0.3)$.
\[
P(X \le 8) = P(10 - Y \le 8) = P(Y \ge 2) = 1 - P(Y \le 1).
\]
From the tables for $n=10$, $p=0.30$: $F_Y(1) = P(Y \le 1) = 0.1493$.
\[
P(X \le 8) = 1 - 0.1493 = \mathbf{0.8507}.
\]
\end{exoresolu}

\courssub{Limitations of Tables and When to Use the Formula}

While tables are fast, they have limitations you must be aware of for the GCE:

\begin{itemize}
    \item \textbf{Limited $n$:} tables typically stop at $n = 50$. If $n > 50$, you cannot use the tables directly (use the Poisson approximation if $p$ is small, or the Normal approximation covered later).
    \item \textbf{Limited $p$:} tables usually give $p = 0.05, 0.10, \dots, 0.50$. If $p = 0.37$, you must use the formula (interpolation is not expected in the examination).
    \item \textbf{Rounding:} tables give 4 decimal places. If the question asks for an \emph{exact} probability or an answer \emph{in terms of $p$}, you \textbf{must} use the formula $\binom{n}{r}p^r(1-p)^{n-r}$.
    \item \textbf{Accumulated rounding:} computing $P(X=r)$ via $F(r)-F(r-1)$ subtracts two rounded values. For a single individual probability, the formula (or the recurrence relation) is often more accurate.
\end{itemize}

\begin{aretenir}[GCE Exam Technique: Quoting Tables]
\begin{itemize}
    \item Always state the distribution: ``$X \sim B(n, p)$''.
    \item State the values you are reading: ``From tables, $P(X \le 2) = 0.6778$''.
    \item Show the subtraction step explicitly for $P(X=r)$, $P(X \ge r)$ and $P(X > r)$.
    \item Give final answers to \textbf{4 decimal places} (or 3 significant figures) unless an exact answer is requested.
    \item If $p > 0.5$, clearly define $Y = n - X$ and state $Y \sim B(n, 1-p)$ before using the tables.
\end{itemize}
\end{aretenir}

\courssub{Worked Examination-Style Example}

\begin{exoresolu}[GCE Style: Cumulative Probabilities and Critical Value]
The probability that a cocoa pod harvested in a certain region is of export quality is $0.65$. A random sample of $12$ pods is selected.
\begin{enumerate}
    \item Model the number of export-quality pods with a suitable distribution, stating your assumptions.
    \item Using tables, find the probability that:
    \begin{enumerate}
        \item at most 7 pods are of export quality;
        \item exactly 7 pods are of export quality;
        \item at least 9 pods are of export quality.
    \end{enumerate}
    \item Find the smallest number $k$ such that the probability of having at most $k$ export-quality pods is at least $0.95$.
\end{enumerate}
\tcblower
\textbf{Solution.}
\begin{enumerate}
    \item Let $X$ be the number of export-quality pods in the sample of 12. There is a fixed number $n=12$ of trials, each trial has two outcomes (export quality or not), the probability $p = 0.65$ is constant, and the pods are selected at random so the trials are independent. Hence
    \[
    X \sim B(12, 0.65).
    \]
    \item Since $p = 0.65 > 0.5$, we use the symmetry trick. Let $Y = \text{number of non-export pods} = 12 - X$. Then $Y \sim B(12, 0.35)$, and we use the tables for $n=12$, $p=0.35$:
    \[
    F_Y(3) = 0.3467, \qquad F_Y(4) = 0.5833, \qquad F_Y(5) = 0.7873.
    \]
    \begin{enumerate}
        \item $P(X \le 7) = P(12 - Y \le 7) = P(Y \ge 5) = 1 - P(Y \le 4) = 1 - 0.5833 = \mathbf{0.4167}$.
        \item $P(X = 7) = P(Y = 5) = F_Y(5) - F_Y(4) = 0.7873 - 0.5833 = \mathbf{0.2040}$.
        \item $P(X \ge 9) = P(12 - Y \ge 9) = P(Y \le 3) = F_Y(3) = \mathbf{0.3467}$.
    \end{enumerate}
    \item We need the smallest $k$ such that $P(X \le k) \ge 0.95$.
    \[
    P(X \le k) = P(Y \ge 12 - k) = 1 - P(Y \le 11 - k) \ge 0.95
    \iff P(Y \le 11 - k) \le 0.05.
    \]
    Scan the table for $Y \sim B(12, 0.35)$:
    \[
    F_Y(0) = 0.0057 \le 0.05, \qquad F_Y(1) = 0.0424 \le 0.05, \qquad F_Y(2) = 0.1513 > 0.05.
    \]
    The largest value of $r$ with $F_Y(r) \le 0.05$ is $r = 1$, so $11 - k = 1$, giving $k = 10$.
    \textbf{Check:} $P(X \le 10) = 1 - F_Y(1) = 1 - 0.0424 = 0.9576 \ge 0.95$, while $P(X \le 9) = 1 - F_Y(2) = 1 - 0.1513 = 0.8487 < 0.95$.
    The smallest $k$ is $\mathbf{10}$.
\end{enumerate}
\end{exoresolu}

\courssub{Summary of Key Formulas for Table Use}

\begin{center}
\begin{tabular}{|l|l|l|}
\hline
\rowcolor{popblueL}
\textbf{Required probability} & \textbf{Expression using $F(r)=P(X\le r)$} & \textbf{Table action} \\
\hline
$P(X \le r)$ & $F(r)$ & Direct reading \\
\hline
$P(X < r) = P(X \le r-1)$ & $F(r-1)$ & Direct reading (row $r-1$) \\
\hline
$P(X = r)$ & $F(r) - F(r-1)$ & Read two rows, subtract \\
\hline
$P(X \ge r)$ & $1 - F(r-1)$ & Read row $r-1$, subtract from 1 \\
\hline
$P(X > r)$ & $1 - F(r)$ & Read row $r$, subtract from 1 \\
\hline
$P(a \le X \le b)$ & $F(b) - F(a-1)$ & Read two rows, subtract \\
\hline
\end{tabular}
\end{center}

\begin{savaistu}[Historical Note: Tables in the Digital Age]
Before electronic calculators and software, printed statistical tables were the \emph{only} practical way to obtain binomial probabilities for moderate $n$, and they remain a standard feature of the GCE Advanced Level examination. Modern scientific calculators have built-in binomial cumulative functions, but the syllabus still requires you to read printed tables. Why? Because table reading tests your understanding of the cumulative structure $P(X \le r)$ and of the complement rules, rather than ``black-box'' button pressing. Master both: use a calculator to check your answers while revising, but show full table-reading working in the examination.
\end{savaistu}

\begin{exercice}[Practice: Cumulative Tables and Critical Values]
Let $X \sim B(20, 0.3)$. Using cumulative binomial tables, find:
\begin{enumerate}
    \item $P(X \le 5)$
    \item $P(X = 5)$
    \item $P(X \ge 5)$
    \item $P(4 \le X \le 7)$
    \item The smallest $r$ such that $P(X \le r) > 0.99$.
    \item The largest $r$ such that $P(X \ge r) > 0.90$.
\end{enumerate}
\end{exercice}

\begin{corrige}[Exercise 1]
From the tables for $n=20$, $p=0.30$:
\[
F(2)=0.0355,\ F(3)=0.1071,\ F(4)=0.2375,\ F(5)=0.4164,\ F(7)=0.7723,\ F(10)=0.9830,\ F(11)=0.9949.
\]
\begin{enumerate}
    \item $P(X \le 5) = F(5) = \mathbf{0.4164}$.
    \item $P(X=5) = F(5) - F(4) = 0.4164 - 0.2375 = \mathbf{0.1789}$.
    \item $P(X \ge 5) = 1 - F(4) = 1 - 0.2375 = \mathbf{0.7625}$.
    \item $P(4 \le X \le 7) = F(7) - F(3) = 0.7723 - 0.1071 = \mathbf{0.6652}$.
    \item Scanning the column: $F(10) = 0.9830 < 0.99$ and $F(11) = 0.9949 > 0.99$. The smallest $r$ is $\mathbf{11}$.
    \item $P(X \ge r) > 0.90 \iff 1 - F(r-1) > 0.90 \iff F(r-1) < 0.10$.
    Now $F(2) = 0.0355 < 0.10$ but $F(3) = 0.1071 > 0.10$, so $r - 1 = 2$, giving $r = 3$.
    \textbf{Check:} $P(X \ge 3) = 1 - F(2) = 0.9645 > 0.90$, while $P(X \ge 4) = 1 - F(3) = 0.8929 < 0.90$. The largest $r$ is $\mathbf{3}$.
\end{enumerate}
\end{corrige}

\coursec{Fitting a Binomial Distribution to Data and Applications}

In the previous section we learnt how to read cumulative binomial probabilities from tables, so that for a \emph{given} model $X\sim B(n,p)$ we can answer any question about $P(X\le r)$. But real life rarely hands us the model directly. A quality-control officer in a Douala factory, a biologist studying families, or a teacher analysing multiple-choice tests first collects \emph{data}, and only then asks: \emph{does the binomial distribution describe this data well, and if so, with which value of $p$?} This reverse process is called \cle{fitting a distribution}, and it is one of the most frequently examined skills in the GCE Advanced Level paper (Lesson 146).

\courssub{What Does ``Fitting a Distribution'' Mean?}

\begin{experience}[Counting boys in families]
A social researcher in Yaoundé selects $100$ families, each having exactly $4$ children, and records the number of boys in each family. She obtains:
\begin{center}
\begin{tabular}{|l|c|c|c|c|c|}
\hline
\rowcolor{popblueL} Number of boys $r$ & $0$ & $1$ & $2$ & $3$ & $4$ \\
\hline Number of families $f$ & $5$ & $28$ & $38$ & $22$ & $7$ \\
\hline
\end{tabular}
\end{center}
If the sex of each child is independent of the others and the probability of a boy is the same constant $p$ for every birth, then the number of boys $X$ in a family of $4$ should follow $B(4,p)$. The question is: \emph{which} $p$, and \emph{how well} does the model match the data?
\end{experience}

\begin{definition}[Fitting a binomial distribution]
To \cle{fit} a binomial distribution $B(n,p)$ to observed data means to:
\begin{enumerate}
\item choose the parameters $n$ and $p$ (usually $n$ is known from the context and $p$ is estimated from the data);
\item compute the \cle{expected frequencies} $E_r = N\times P(X=r)$, where $N$ is the total number of observations;
\item compare the expected frequencies with the observed frequencies $O_r$ to judge whether the model is adequate.
\end{enumerate}
\end{definition}

The intuition is simple: if the model is correct, then in $N$ repetitions of the experiment, the value $r$ should occur \emph{on average} $N\times P(X=r)$ times. So the fitted model ``predicts'' the frequencies, and we check the prediction against reality.

\courssub{Estimating $p$ from the Data}

Since $n$ is fixed by the situation (here $n=4$ children), the only unknown is $p$. We estimate it using the \emph{mean} of the data, because for $X\sim B(n,p)$ we proved that $E(X)=np$.

\begin{propriete}[Estimating $p$]
If a sample of $N$ observations from an alleged $B(n,p)$ distribution has mean $\bar{x}$, then the natural estimate of $p$ is
\[ \hat{p}=\frac{\bar{x}}{n}. \]
\end{propriete}

\textbf{Justification.} The theoretical mean is $np$. Fitting consists of making the theoretical model match the data, so we demand that the theoretical mean equal the observed mean: $n\hat{p}=\bar{x}$, giving $\hat{p}=\bar{x}/n$. \hfill$\square$

\begin{attention}[The classic trap]
Do \textbf{not} estimate $p$ as the relative frequency of one particular value. For instance, in the family data, $\dfrac{38}{100}=0.38$ is the proportion of families with $2$ boys --- it is \emph{not} an estimate of $p$. The estimate must come from the \textbf{mean} of the whole distribution: $\hat{p}=\bar{x}/n$. Using a single relative frequency is one of the most common errors at the GCE and costs several marks.
\end{attention}

\courssub{Computing Expected Frequencies}

\begin{methode}[The 4-step fitting procedure]
\begin{enumerate}
\item \textbf{Estimate $p$.} Compute the sample mean $\bar{x}=\dfrac{\sum f_r\, r}{\sum f_r}$, then set $\hat{p}=\bar{x}/n$. If $\hat{p}$ is not a tabulated value (e.g. $0.05,0.10,\dots,0.50$), round it to the nearest tabulated value \emph{only if} you intend to use cumulative binomial tables; otherwise keep the exact value and compute probabilities directly from the formula.
\item \textbf{Compute the probabilities.} Using $P(X=r)=\dbinom{n}{r}\hat{p}^{\,r}(1-\hat{p})^{n-r}$, find $P(X=r)$ for $r=0,1,\dots,n$. Keep at least $4$ decimal places during calculations.
\item \textbf{Multiply by $N$.} The expected frequencies are $E_r = N\times P(X=r)$. Round each $E_r$ to $1$ decimal place (or as instructed) and verify that $\sum E_r = N$ (adjust the last decimal if necessary due to rounding).
\item \textbf{Compare.} Place observed $O_r$ and expected $E_r$ side by side in a table (and ideally a bar chart) and comment on the goodness of fit.
\end{enumerate}
\end{methode}

\begin{aretenir}[Expected frequencies need not be integers]
An expected frequency is a \emph{mean} number of occurrences, so values like $21.4$ or $6.83$ are perfectly correct. Keep at least $3$--$4$ decimal places in the probabilities during the calculation and \textbf{round only at the very end} (usually to $1$ decimal place). Rounding too early can make the expected frequencies fail to add up to $N$ and destroys accuracy.
\end{aretenir}

\courssub{Worked Example: Fitting $B(4,p)$ to the Family Data}

\begin{exoresolu}[Fitting a binomial distribution to family data]
For the $100$ families of $4$ children above, fit a binomial distribution and comment on the fit.
\tcblower
\textbf{Step 1: Estimate $p$.}
\[ \bar{x}=\frac{0(5)+1(28)+2(38)+3(22)+4(7)}{100}=\frac{0+28+76+66+28}{100}=\frac{198}{100}=1.98. \]
Hence $\hat{p}=\dfrac{\bar{x}}{n}=\dfrac{1.98}{4}=0.495$. Since $0.495$ is very close to $0.5$ and $0.5$ is a tabulated value, we use $\hat{p}=0.5$ for the fit (this also allows us to use the cumulative binomial tables if needed).

\textbf{Step 2: Probabilities for $X\sim B(4,0.5)$.} Since $\binom{4}{r}=(1,4,6,4,1)$ and $p=q=\tfrac12$:
\[ P(X=r)=\binom{4}{r}\left(\tfrac12\right)^{4}=\binom{4}{r}\times\frac{1}{16}. \]
\begin{center}
\begin{tabular}{|l|c|c|c|c|c|c|}
\hline
\rowcolor{popblueL} $r$ & $0$ & $1$ & $2$ & $3$ & $4$ & Total \\
\hline $P(X=r)$ & $0.0625$ & $0.25$ & $0.375$ & $0.25$ & $0.0625$ & $1$ \\
\hline $E_r=100\,P(X=r)$ & $6.25$ & $25.0$ & $37.5$ & $25.0$ & $6.25$ & $100$ \\
\hline Observed $O_r$ & $5$ & $28$ & $38$ & $22$ & $7$ & $100$ \\
\hline
\end{tabular}
\end{center}

\textbf{Step 3: Comment.} The expected frequencies are very close to the observed ones: the largest discrepancy is $3$ families (for $r=1$ and $r=3$). The binomial model $B(4,0.5)$ fits the data \emph{well}, supporting the hypothesis that boys and girls are equally likely and that births are independent.
\end{exoresolu}

\begin{popfigure}
\begin{center}
\begin{tikzpicture}
\begin{axis}[
    ybar, bar width=0.32cm,
    width=12.5cm, height=7.2cm,
    xlabel={Number of boys $r$},
    ylabel={Frequency},
    symbolic x coords={0,1,2,3,4},
    xtick=data,
    ymin=0, ymax=45,
    legend style={at={(0.5,1.02)}, anchor=south, legend columns=2, draw=none},
    ymajorgrids=true, grid style={popline!40},
    every axis plot/.append style={draw=none},
]
\addplot[fill=popblue] coordinates {(0,5) (1,28) (2,38) (3,22) (4,7)};
\addplot[fill=poporange] coordinates {(0,6.25) (1,25) (2,37.5) (3,25) (4,6.25)};
\legend{Observed $O_r$, {Expected $E_r$ (fitted $B(4,0.5)$)}}
\end{axis}
\end{tikzpicture}
\end{center}
\legende{Grouped bar chart: observed frequencies versus expected frequencies from the fitted binomial distribution $B(4,0.5)$ for $100$ families of $4$ children.}
\end{popfigure}

The bars are almost the same height everywhere: this visual agreement is exactly what a good fit looks like.

\courssub{Commenting on the Goodness of Fit}

At GCE level, the comparison is made \emph{by inspection}:

\begin{itemize}
\item compute the differences $|O_r-E_r|$;
\item if all differences are small relative to the frequencies, conclude that \emph{``the binomial model is a good fit''};
\item if one or more differences are large, conclude that the fit is poor and suggest why (e.g.\ $p$ not constant, trials not independent).
\end{itemize}

\begin{savaistu}[Measuring the fit precisely: the chi-squared test]
Statisticians measure goodness of fit with a single number, the \cle{chi-squared statistic} $\chi^2=\sum \dfrac{(O_r-E_r)^2}{E_r}$, and compare it with critical values. For our family data, $\chi^2 = \frac{1.5625}{6.25}+\frac{9}{25}+\frac{0.25}{37.5}+\frac{9}{25}+\frac{0.5625}{6.25}\approx 1.07$, a very small value, confirming an excellent fit. The formal chi-squared test is \emph{beyond the GCE syllabus} --- you only need to compare and comment --- but you will meet it at university.
\end{savaistu}

\begin{attention}[Exam tip --- GCE]
Fitting questions typically carry \textbf{6 to 8 marks}. To collect them all:
\begin{enumerate}
\item show the calculation of $\bar{x}$ and state $\hat{p}=\bar{x}/n$ explicitly;
\item present the \textbf{full probability table} $P(X=r)$ for every $r$ from $0$ to $n$, not just the values you find interesting;
\item give expected frequencies to at least $1$ decimal place and check they sum to $N$;
\item finish with a clear sentence of conclusion (``good fit'' / ``poor fit, because\ldots'').
\end{enumerate}
\end{attention}

\courssub{Applications of the Binomial Distribution (Lesson 149)}

The binomial model appears whenever we repeat a fixed number of independent trials with the same probability of success. Here are four standard fields of application.

\textbf{1. Quality control.} A factory in Douala produces sachets of detergent, and past experience shows that $2\%$ of sachets are underweight. An inspector samples $20$ sachets at random. If $X$ is the number of underweight sachets, then $X\sim B(20,0.02)$, and
\[ P(X\ge 2)=1-P(X=0)-P(X=1)=1-(0.98)^{20}-20(0.02)(0.98)^{19}\approx 0.0599. \]
If the inspector finds $2$ or more defective sachets (an event of probability only about $6\%$ under the model), the batch is rejected and the machine is checked. The binomial model thus provides an objective rule for accepting or rejecting production.

\textbf{2. Guessing in multiple-choice tests.} A candidate answers $10$ MCQ questions, each with $4$ options, purely at random. The number of correct answers is $X\sim B(10,0.25)$. Then
\[ P(X\ge 5)=1-P(X\le 4)\approx 1-0.9219=0.0781 \]
(from cumulative tables). So a pure guesser scores $5$ or more only about $8\%$ of the time: high scores are genuine evidence of knowledge, which is exactly why MCQ marking schemes are fair.

\textbf{3. Penalty kicks in sports.} Suppose a striker for the Indomitable Lions converts $80\%$ of his penalties. In a training session he takes $5$ kicks; the number scored is $X\sim B(5,0.8)$. Then
\[ P(X=5)=(0.8)^5=0.3277,\qquad P(X\ge 4)=P(X=4)+P(X=5)=5(0.8)^4(0.2)+(0.8)^5\approx 0.7373. \]
Coaches use such calculations to compare players and to decide who should take the decisive kick.

\textbf{4. Medical screening.} A disease affects $5\%$ of a population. In a random sample of $15$ people, the number infected is $X\sim B(15,0.05)$, so
\[ P(X=0)=(0.95)^{15}\approx 0.4633, \]
meaning that in nearly half of such samples nobody is infected even though the disease is present in the population --- a crucial warning against concluding ``the village is disease-free'' from one small sample.

\courssub{When Is the Binomial Model Appropriate?}

\begin{aretenir}[Conditions for a binomial model]
$X\sim B(n,p)$ is appropriate only if \textbf{all four} conditions hold:
\begin{enumerate}
\item a \textbf{fixed} number $n$ of trials;
\item each trial has exactly \textbf{two outcomes} (success/failure);
\item the probability $p$ of success is \textbf{constant} from trial to trial;
\item the trials are \textbf{independent}.
\end{enumerate}
\end{aretenir}

The model \emph{fails} when, for example:
\begin{itemize}
\item \textbf{$p$ is not constant:} a basketball player's probability of scoring rises as he ``warms up''; a machine wears out during the day so late items are more often defective.
\item \textbf{Trials are dependent:} drawing balls from an urn \emph{without replacement}; the health of members of the same family (infections spread between them); questions in an exam that build on each other.
\item \textbf{$n$ is not fixed:} counting the number of phone calls received in an hour (no fixed number of trials --- this will lead us to the Poisson distribution later in the chapter).
\end{itemize}

A poor fit when fitting $B(n,p)$ to data is often the \emph{symptom} that one of these conditions is violated. For instance, if observed data show far more extreme values ($0$ and $n$) than the fitted binomial predicts, this suggests that $p$ varies between groups of trials.

\begin{exercice}[Fitting and judging a claim]
A seed company claims that $75\%$ of its tomato seeds germinate. A farmer plants $200$ rows of $5$ seeds and records the number germinating per row:
\begin{center}
\begin{tabular}{|l|c|c|c|c|c|c|}
\hline
\rowcolor{popblueL} Number germinating $r$ & $0$ & $1$ & $2$ & $3$ & $4$ & $5$ \\
\hline Number of rows & $1$ & $7$ & $30$ & $62$ & $72$ & $28$ \\
\hline
\end{tabular}
\end{center}
\begin{enumerate}
\item Estimate $p$ from the data and fit a binomial distribution $B(5,\hat{p})$. Compute the expected frequencies (keep $\hat{p}$ exact for the calculation).
\item The company claims $p=0.75$. Using $p=0.75$, compute the expected frequencies for $B(5,0.75)$.
\item Compare both fits with the observed data. Comment on the company's claim.
\end{enumerate}
\end{exercice}

\begin{corrige}[Exercise 1]
\textbf{1. Fitting with estimated $\hat{p}$.}
\[ \bar{x}=\frac{0(1)+1(7)+2(30)+3(62)+4(72)+5(28)}{200}=\frac{0+7+60+186+288+140}{200}=\frac{681}{200}=3.405. \]
Hence $\hat{p}=\dfrac{\bar{x}}{n}=\dfrac{3.405}{5}=0.681$.
Using $p=0.681$, $q=0.319$, we compute $P(X=r)=\binom{5}{r}(0.681)^r(0.319)^{5-r}$:
\begin{align*}
P(X=0) &= (0.319)^5 \approx 0.0033, & E_0 &= 200\times0.0033 = 0.66 \approx 0.7,\\
P(X=1) &= 5(0.681)(0.319)^4 \approx 0.0353, & E_1 &= 200\times0.0353 = 7.06 \approx 7.1,\\
P(X=2) &= 10(0.681)^2(0.319)^3 \approx 0.1507, & E_2 &= 200\times0.1507 = 30.14 \approx 30.1,\\
P(X=3) &= 10(0.681)^3(0.319)^2 \approx 0.3217, & E_3 &= 200\times0.3217 = 64.34 \approx 64.3,\\
P(X=4) &= 5(0.681)^4(0.319) \approx 0.3435, & E_4 &= 200\times0.3435 = 68.70 \approx 68.7,\\
P(X=5) &= (0.681)^5 \approx 0.1455, & E_5 &= 200\times0.1455 = 29.10 \approx 29.1.
\end{align*}
Sum of $E_r = 0.7+7.1+30.1+64.3+68.7+29.1 = 200.0$. Table:
\begin{center}
\begin{tabular}{|l|c|c|c|c|c|c|}
\hline
\rowcolor{popblueL} $r$ & $0$ & $1$ & $2$ & $3$ & $4$ & $5$ \\
\hline $E_r$ (fitted $\hat{p}=0.681$) & $0.7$ & $7.1$ & $30.1$ & $64.3$ & $68.7$ & $29.1$ \\
\hline $O_r$ & $1$ & $7$ & $30$ & $62$ & $72$ & $28$ \\
\hline
\end{tabular}
\end{center}
The fit is \emph{very good}: all expected frequencies are close to observed ones (largest difference $\approx 3.3$ for $r=4$).

\textbf{2. Testing the claim $p=0.75$.}
Using $B(5,0.75)$: $P(X=r)=\binom{5}{r}(0.75)^r(0.25)^{5-r}$ gives probabilities $0.0010,\ 0.0146,\ 0.0879,\ 0.2637,\ 0.3955,\ 0.2373$, hence expected frequencies (multiplying by $200$):
\begin{center}
\begin{tabular}{|l|c|c|c|c|c|c|}
\hline
\rowcolor{popblueL} $r$ & $0$ & $1$ & $2$ & $3$ & $4$ & $5$ \\
\hline $E_r$ (claim $p=0.75$) & $0.2$ & $2.9$ & $17.6$ & $52.7$ & $79.1$ & $47.5$ \\
\hline $O_r$ & $1$ & $7$ & $30$ & $62$ & $72$ & $28$ \\
\hline
\end{tabular}
\end{center}

\textbf{3. Conclusion.} The fit using the estimated $\hat{p}=0.681$ is excellent, while the fit using the company's claimed $p=0.75$ is \emph{poor}: the model predicts about $48$ rows with all $5$ seeds germinating, but only $28$ were observed, and it underestimates the low counts ($30$ rows observed with $2$ germinating against $17.6$ expected). The data suggest a true germination rate near $0.68$, so the company's claim of $75\%$ is not supported.
\end{corrige}

This section closes our study of the binomial distribution as a complete toolkit: probabilities, expectation, recurrence, cumulative tables, and now fitting to real data. We next turn to a close relative --- the \emph{geometric distribution} --- which asks not ``how many successes in $n$ trials?'' but ``how many trials until the first success?''

\coursec{The Geometric Distribution}

In the previous section we fitted a binomial distribution to observed data and applied $B(n,p)$ to situations where the \emph{number of trials is fixed in advance}. But many real questions are of a different kind: we do not ask ``how many successes in $n$ trials?'' but rather ``\emph{how long must we wait for the first success?}'' How many throws of a die until the first six appears? How many items from a production line must be inspected until the first defective one is found? How many penalties until the first goal is scored? These questions lead to a new discrete distribution: the \cle{geometric distribution}. This topic corresponds to \textbf{Lessons 147 and 148} of the official syllabus.

\courssub{From an Experiment to a Model}

\begin{experience}[Waiting for the first six]
Take an ordinary fair die and throw it repeatedly, recording the number of the throw on which the \emph{first} six appears. Sometimes the six comes immediately ($X=1$); sometimes you wait 4, 7, or even 20 throws. Repeat the experiment many times and draw a bar chart of the frequencies of $X$. You will observe a characteristic shape: the value $1$ is the most frequent, and the frequencies decrease steadily, with an occasional very long wait. This decreasing shape is the signature of the geometric distribution.
\end{experience}

Let us compute the probabilities for this experiment. Let $p=\tfrac16$ be the probability of a six on each throw and $q=1-p=\tfrac56$. Since the throws are independent:
\begin{align*}
P(X=1)&=p,\\
P(X=2)&=P(\text{failure then success})=qp,\\
P(X=3)&=P(\text{failure, failure, success})=q^2p,\\
P(X=r)&=q^{\,r-1}p.
\end{align*}
The probabilities form a \emph{geometric progression} with common ratio $q$ --- hence the name of the distribution.

\courssub{Definition and Conditions}

\begin{definition}[Geometric distribution]
A discrete random variable $X$ follows a \cle{geometric distribution} with parameter $p$ ($0<p<1$), written
\[ X\sim \mathrm{Geo}(p), \]
if $X$ counts the \textbf{number of independent Bernoulli trials up to and including the first success}, where each trial has the same probability $p$ of success. Writing $q=1-p$,
\[ \boxed{P(X=r)=q^{\,r-1}p, \qquad r=1,2,3,\dots} \]
\end{definition}

The conditions for the model are exactly those of a sequence of Bernoulli trials, but with \emph{no upper limit} on the number of trials:
\begin{itemize}
\item each trial results in \cle{success} or \cle{failure};
\item the trials are \cle{independent};
\item the probability of success $p$ is \cle{constant} from trial to trial;
\item trials continue until the first success occurs (the number of trials is \emph{unlimited}).
\end{itemize}

\begin{attention}[Check the formula is a genuine distribution]
Note that $\displaystyle\sum_{r=1}^{\infty}q^{r-1}p=p\cdot\frac{1}{1-q}=\frac{p}{p}=1$, using the sum of the geometric series $1+q+q^2+\cdots=\dfrac{1}{1-q}$. So the probabilities do add up to $1$, as they must.
\end{attention}

\begin{attention}[Alternative definition]
Some textbooks define the geometric distribution as the number of \emph{failures} before the first success (support $r=0,1,2,\dots$). The GCE Advanced Level syllabus uses the definition above: $X$ is the \textbf{trial number of the first success} (support $r=1,2,3,\dots$). Always check which definition is being used; the formulas for expectation and variance differ.
\end{attention}

\courssub{The Shape of the Distribution}

\begin{popfigure}
\begin{center}
\begin{tikzpicture}
\begin{axis}[
    ybar, bar width=8pt,
    width=13cm, height=7cm,
    xlabel={$r$ (trial of first success)}, ylabel={$P(X=r)$},
    xmin=0.3, xmax=16.7, ymin=0, ymax=0.28,
    xtick={1,2,3,4,5,6,7,8,9,10,11,12,13,14,15,16},
    ytick={0,0.05,0.10,0.15,0.20,0.25},
    axis lines=left,
    every axis x label/.style={at={(ticklabel* cs:1.02)}, anchor=west},
]
\addplot[fill=popblue, draw=popdark] coordinates {
(1,0.25) (2,0.1875) (3,0.1406) (4,0.1055) (5,0.0791) (6,0.0593)
(7,0.0445) (8,0.0334) (9,0.0250) (10,0.0188) (11,0.0141) (12,0.0106)
(13,0.0079) (14,0.0059) (15,0.0045) (16,0.0033)};
\end{axis}
\end{tikzpicture}
\legende{Bar chart of $X\sim\mathrm{Geo}(0.25)$: the probabilities $P(X=r)=0.75^{\,r-1}(0.25)$ decrease geometrically, with a long right tail.}
\end{center}
\end{popfigure}

The bar chart shows two essential features: the distribution is \emph{always} decreasing (so the most likely value of $X$ is $1$), and it has a \emph{long right tail}: very long waits are unlikely but never impossible.

\begin{propriete}[Mode of the geometric distribution]
If $X\sim\mathrm{Geo}(p)$, the \cle{mode} of $X$ is always $1$, because
\[ P(X=r)=q\,P(X=r-1)<P(X=r-1)\quad\text{for all } r\ge 2. \]
\end{propriete}

\courssub{Cumulative Probabilities}

A very useful trick: the event ``$X>r$'' means \emph{the first $r$ trials are all failures}. By independence this has probability $q^r$.

\begin{propriete}[Cumulative probabilities]
If $X\sim\mathrm{Geo}(p)$ and $q=1-p$, then for every positive integer $r$,
\[ \boxed{P(X>r)=q^{r}} \qquad\text{and}\qquad \boxed{P(X\le r)=1-q^{r}}. \]
\end{propriete}

\begin{popfigure}
\begin{center}
\begin{tikzpicture}[scale=1.05]
\draw[->, thick] (-0.3,0) -- (9.6,0) node[right] {\small trial number};
\foreach \x in {1,...,9} \draw (\x,0.08) -- (\x,-0.08) node[below] {\small \x};
\foreach \x in {1,...,6} {
  \draw[fill=poporange, draw=popdark] (\x-0.28,0.55) rectangle (\x+0.28,1.15);
  \node at (\x,0.85) {\small F};
}
\draw[fill=popgreen, draw=popdark] (7-0.28,0.55) rectangle (7+0.28,1.15);
\node at (7,0.85) {\small S};
\draw[decorate, decoration={brace, amplitude=5pt}, thick, popdark] (0.72,1.35) -- (6.28,1.35);
\node[popdark] at (3.5,1.85) {\small $r=6$ consecutive failures: probability $q^{6}$};
\node at (8.3,0.85) {\small $\cdots$};
\end{tikzpicture}
\legende{$P(X>r)=q^{r}$ is the probability that the first $r$ trials are all failures (F); the first success (S) then occurs later.}
\end{center}
\end{popfigure}

\courssub{Expectation and Variance}

\begin{propriete}[Expectation and variance of $\mathrm{Geo}(p)$]
If $X\sim\mathrm{Geo}(p)$ with $q=1-p$, then
\[ \boxed{E(X)=\frac{1}{p}} \qquad\text{and}\qquad \boxed{\mathrm{Var}(X)=\frac{q}{p^{2}}}. \]
\emph{(The derivations use differentiation of geometric series; they are an extension, not required at GCE, but the results must be known and used.)}
\end{propriete}

\textbf{Sketch of the derivation (extension).} $E(X)=\sum_{r\ge1} r q^{r-1}p = p\sum_{r\ge1} r q^{r-1}$. Differentiating $\sum_{r\ge0}q^{r}=\dfrac{1}{1-q}$ term by term gives $\sum_{r\ge1} r q^{r-1}=\dfrac{1}{(1-q)^2}=\dfrac{1}{p^2}$, hence $E(X)=\dfrac{1}{p}$. A similar second differentiation yields $E(X^2)$ and then $\mathrm{Var}(X)=\dfrac{q}{p^2}$.

The expectation matches intuition perfectly: if the probability of a six is $\tfrac16$, we \emph{expect} to wait $6$ throws for the first six; if $p$ is small, the expected wait is long.

\begin{exoresolu}[First defective item on a production line]
A machine at a soap factory in Douala produces defective items independently with probability $0.04$. Let $X$ be the position of the \emph{first} defective item inspected. Find:
\begin{enumerate}
\item the probability that the first defective item is the 5th inspected;
\item the probability that more than 20 items are inspected before the first defective one;
\item the expected number of items inspected and the standard deviation of $X$.
\end{enumerate}
\tcblower
Here $X\sim\mathrm{Geo}(0.04)$, so $p=0.04$ and $q=0.96$.
\begin{enumerate}
\item $P(X=5)=q^{4}p=(0.96)^{4}(0.04)\approx 0.8493\times 0.04\approx 0.0340$.
\item $P(X>20)=q^{20}=(0.96)^{20}\approx 0.4420$.
\item $E(X)=\dfrac{1}{p}=\dfrac{1}{0.04}=25$ items, and
\[ \mathrm{Var}(X)=\frac{q}{p^{2}}=\frac{0.96}{0.0016}=600, \qquad \sigma=\sqrt{600}\approx 24.5. \]
\end{enumerate}
Note the large standard deviation: the waiting time is very spread out, as the long right tail suggests.
\end{exoresolu}

\courssub{The Memoryless Property}

Suppose you have already thrown the die $10$ times without a six. Does the die ``owe'' you a six? Of course not --- the remaining wait has exactly the same distribution as at the start. This is the \cle{memoryless property}, and the geometric distribution is the \emph{only} discrete distribution with this property.

\begin{propriete}[Memoryless property]
If $X\sim\mathrm{Geo}(p)$, then for all positive integers $m$ and $n$,
\[ \boxed{P(X>m+n \mid X>m)=P(X>n)}. \]
\end{propriete}

\textbf{Verification.} Using conditional probability and $P(X>r)=q^r$:
\[ P(X>m+n\mid X>m)=\frac{P(X>m+n)}{P(X>m)}=\frac{q^{m+n}}{q^{m}}=q^{n}=P(X>n). \]
So knowledge that the first $m$ trials were failures changes nothing about the additional waiting time.

\begin{exemplebox}[Checking the memoryless property numerically]
Let $X\sim\mathrm{Geo}(0.25)$ (e.g.\ waiting for the first ``head'' when spinning a biased coin with $P(\text{success})=0.25$). Take $m=4$, $n=3$:
\[ P(X>7\mid X>4)=\frac{(0.75)^{7}}{(0.75)^{4}}=(0.75)^{3}=0.4219=P(X>3). \]
Given that $4$ failures have already occurred, the probability of waiting \emph{more than 3 further trials} is the same as the original probability of waiting more than $3$ trials.
\end{exemplebox}

\courssub{Geometric versus Binomial: Do Not Confuse Them}

\begin{attention}[Common mistake: $\mathrm{Geo}(p)$ versus $B(n,p)$]
The two distributions answer \emph{different questions} about the \emph{same} kind of trials:
\begin{itemize}
\item $B(n,p)$: the number of trials $n$ is \textbf{fixed}; we count the \textbf{number of successes}. The formula involves $\binom{n}{r}$.
\item $\mathrm{Geo}(p)$: the number of trials is \textbf{not fixed}; we count the \textbf{trial number of the first success}. The formula is $q^{r-1}p$, with \emph{no binomial coefficient}.
\end{itemize}
Exam trap: ``the probability that the \emph{first} success occurs on the 5th trial'' is $q^{4}p$ (geometric), \textbf{not} $\binom{5}{1}p q^{4}$ (which is the probability of \emph{exactly one} success in 5 trials).
\end{attention}

\begin{methode}[Exam tip (GCE): ``first success after the $k$-th trial'']
When a GCE question asks for ``the probability that the first success occurs \emph{after} the $k$-th trial'' (or ``more than $k$ trials are needed''):
\begin{enumerate}
\item Identify $p$ and compute $q=1-p$.
\item Recognise that ``first success after the $k$-th trial'' means the first $k$ trials all fail.
\item Write directly $P(X>k)=q^{k}$ --- no summation needed.
\end{enumerate}
Similarly, ``first success \emph{on or before} the $k$-th trial'' gives $P(X\le k)=1-q^{k}$.
\end{methode}

\begin{exoresolu}[GCE-style problem: waiting for a goal]
A footballer scores from a free kick with probability $0.15$, independently of previous attempts. Let $X$ be the number of free kicks taken up to and including the first goal.
\begin{enumerate}
\item State the distribution of $X$ and the conditions required.
\item Find the probability that his first goal comes on his 4th free kick.
\item Find the probability that he needs more than 10 free kicks to score.
\item Find the expected number of free kicks needed and $\mathrm{Var}(X)$.
\item Given that he has already missed 6 free kicks, find the probability that he needs more than 4 further attempts.
\end{enumerate}
\tcblower
\begin{enumerate}
\item $X\sim\mathrm{Geo}(0.15)$, provided attempts are independent and the probability of scoring is constant at $p=0.15$ for each attempt. Here $q=0.85$.
\item $P(X=4)=q^{3}p=(0.85)^{3}(0.15)\approx 0.6141\times 0.15\approx 0.0921$.
\item $P(X>10)=q^{10}=(0.85)^{10}\approx 0.1969$.
\item $E(X)=\dfrac{1}{0.15}\approx 6.67$ free kicks; $\mathrm{Var}(X)=\dfrac{0.85}{0.15^{2}}=\dfrac{0.85}{0.0225}\approx 37.78$.
\item By the memoryless property,
\[ P(X>6+4\mid X>6)=P(X>4)=q^{4}=(0.85)^{4}\approx 0.5220. \]
\end{enumerate}
\end{exoresolu}

\courssub{Summary and Practice}

\begin{aretenir}[Key points on the geometric distribution]
\begin{itemize}
\item $X\sim\mathrm{Geo}(p)$: $X$ = number of independent trials up to and including the first success; $P(X=r)=q^{r-1}p$, $r=1,2,3,\dots$
\item $P(X>r)=q^{r}$ and $P(X\le r)=1-q^{r}$.
\item $E(X)=\dfrac{1}{p}$, $\mathrm{Var}(X)=\dfrac{q}{p^{2}}$, mode $=1$.
\item Memoryless: $P(X>m+n\mid X>m)=P(X>n)$.
\item Do not confuse with $B(n,p)$: fixed $n$ (binomial) versus waiting time for the first success (geometric).
\end{itemize}
\end{aretenir}

\begin{exercice}[On the geometric distribution]
\begin{enumerate}
\item A fair die is thrown repeatedly. Let $X$ be the throw on which the first six appears. Find (a) $P(X=3)$; (b) $P(X>6)$; (c) $E(X)$ and $\mathrm{Var}(X)$.
\item In Yaound\'e, a call centre agent makes a sale on each call with probability $0.1$, independently of other calls. Find the probability that her first sale (a) occurs on the 8th call; (b) occurs after the 15th call. What is her expected number of calls for the first sale?
\item A student guesses every answer of a multiple-choice quiz with 4 options per question. Find the probability that the first correct guess comes (a) on the 2nd question; (b) after the 5th question. Given that the first 4 answers were wrong, find the probability that more than 3 further questions are needed for the first correct answer.
\end{enumerate}
\end{exercice}

\begin{corrige}[Exercise]
\begin{enumerate}
\item $X\sim\mathrm{Geo}\!\left(\tfrac16\right)$, $q=\tfrac56$. (a) $P(X=3)=\left(\tfrac56\right)^{2}\tfrac16=\tfrac{25}{216}\approx0.1157$. (b) $P(X>6)=\left(\tfrac56\right)^{6}\approx0.3349$. (c) $E(X)=6$, $\mathrm{Var}(X)=\dfrac{5/6}{1/36}=30$.
\item $X\sim\mathrm{Geo}(0.1)$, $q=0.9$. (a) $P(X=8)=(0.9)^{7}(0.1)\approx0.0478$. (b) $P(X>15)=(0.9)^{15}\approx0.2059$. $E(X)=\dfrac{1}{0.1}=10$ calls.
\item $p=\tfrac14$, $q=\tfrac34$. (a) $P(X=2)=\tfrac34\cdot\tfrac14=\tfrac{3}{16}=0.1875$. (b) $P(X>5)=\left(\tfrac34\right)^{5}\approx0.2373$. Memoryless property: $P(X>4+3\mid X>4)=P(X>3)=\left(\tfrac34\right)^{3}\approx0.4219$.
\end{enumerate}
\end{corrige}

\begin{savaistu}[Why ``geometric''?]
The distribution takes its name from the \emph{geometric series}: the probabilities $p,\ qp,\ q^{2}p,\ q^{3}p,\dots$ form a geometric progression, and their sum $p(1+q+q^{2}+\cdots)=\dfrac{p}{1-q}=1$ is exactly the geometric series formula you met in Pure Mathematics. The geometric distribution is also the discrete analogue of the \emph{exponential distribution}, which models waiting times in continuous time --- both share the memoryless property.
\end{savaistu}

\coursec{The Poisson Distribution: Characteristics, Expectation and Variance}

In the previous section we studied the \cle{geometric distribution}, which counts the number of \emph{trials} needed to obtain the first success. Trials there were clearly separated: one attempt, then another, then another. But many real-life situations do not come in separate trials at all. Think of the phone calls arriving at the MTN call centre in Douala during one hour, the accidents occurring along the Douala--Yaoundé highway during one month, or the misprints on the pages of a textbook. There is no fixed number $n$ of ``trials''; events simply \emph{happen}, scattered at random along a continuous stretch of time or space. For such situations we need a new model: the \cle{Poisson distribution}, named after the French mathematician Siméon Denis Poisson.

\courssub{Random Events in a Continuum: the Poisson Process}

\begin{experience}[Observing calls at a call centre]
Imagine you sit next to the switchboard of a small business in Yaoundé and you record the exact times at which calls arrive during one hour, say from 9:00 to 10:00. On a typical morning you might observe something like this: calls at 9:03, 9:11, 9:12, 9:27, 9:41, 9:58 --- six calls in all. The next day, same hour, you might count only three calls; the day after, nine. The calls do not arrive at regular intervals like a clock ticking; they arrive \emph{at random}. Yet, over many days, the \emph{average} number of calls per hour is remarkably stable, say about $5$ calls per hour. This combination --- individual unpredictability but long-term stability of the average rate --- is exactly the signature of a \cle{Poisson process}.
\end{experience}

\begin{popfigure}
\begin{center}
\begin{tikzpicture}[scale=1.0]
% timeline
\draw[->, thick, popdark] (0,0) -- (10.6,0) node[below left, font=\small] {time};
\foreach \x/\t in {0/9:00, 2.5/9:15, 5/9:30, 7.5/9:45, 10/10:00}
  \draw[thick, popdark] (\x,0.12) -- (\x,-0.12) node[below, font=\small] {\t};
% calls
\foreach \x in {0.6, 2.2, 2.7, 5.4, 8.2, 9.6}
  { \draw[very thick, poporange] (\x,0) -- (\x,0.85);
    \fill[poporange] (\x,0.85) circle (2.2pt); }
\node[font=\small, text=poporange, align=center] at (5,1.5) {incoming calls (6 calls in one hour)};
% rate annotation
\draw[<->, thick, popblue] (0,-1.1) -- (10,-1.1);
\node[font=\small, text=popblue] at (5,-1.5) {constant average rate: $\lambda = 6$ calls per hour};
\end{tikzpicture}
\legende{Random phone calls over one hour: the instants are unpredictable, calls arrive one at a time and independently, but the average rate per hour stays constant.}
\end{center}
\end{popfigure}

Other classic examples of events occurring randomly in a continuum include:
\begin{itemize}
\item accidents on a given stretch of the Douala--Yaoundé highway per month;
\item misprints per page of a printed book (here the ``continuum'' is space, not time);
\item radioactive decays per second detected by a counter;
\item customers joining the queue at a bank agency per ten minutes;
\item goals scored per football match in a championship.
\end{itemize}

\begin{methode}[The four conditions for a Poisson process]
Before modelling a situation by a Poisson distribution, check that the events occur:
\begin{enumerate}
\item \textbf{Randomly} --- the exact instant (or position) of each event is unpredictable;
\item \textbf{Independently} --- the occurrence of one event does not affect the probability of another (one accident does not ``cause'' or ``prevent'' the next);
\item \textbf{Singly} --- events occur one at a time; two events cannot happen at exactly the same instant (or same point);
\item \textbf{At a constant average rate} --- the mean number of events per \cle{unit interval}, denoted $\lambda$, stays the same throughout the period considered.
\end{enumerate}
If all four conditions hold, the number of events in a fixed interval follows a Poisson distribution.
\end{methode}

\begin{attention}[When Poisson does \emph{not} apply]
Buses arriving at a stop \emph{according to a fixed timetable} fail condition 1 (not random). Goals in a match where scoring a goal encourages the attacking team may fail condition 2 (not independent). Customers arriving in \emph{groups} at a restaurant fail condition 3 (not single). Calls at a call centre over a whole day fail condition 4 if the rate is much higher at midday than at midnight --- in that case, split the day into periods with roughly constant rates.
\end{attention}

\courssub{Definition and Probability Formula}

\begin{definition}[Poisson distribution]
A discrete random variable $X$ follows a \cle{Poisson distribution} with parameter $\lambda > 0$, written $X \sim \mathrm{Po}(\lambda)$, if $X$ takes values in $\{0, 1, 2, 3, \dots\}$ with
\[ P(X = r) = \frac{e^{-\lambda}\,\lambda^{r}}{r!}, \qquad r = 0, 1, 2, 3, \dots \]
Here $\lambda$ is the \emph{mean number of events per unit interval} (per hour, per page, per kilometre, \dots). The \cle{unit interval} is the fixed stretch of time or space (e.g. one hour, one page, one kilometre) for which $\lambda$ is defined.
\end{definition}

Note immediately a useful special case:
\[ P(X = 0) = \frac{e^{-\lambda}\lambda^{0}}{0!} = e^{-\lambda}, \]
since $\lambda^{0} = 1$ and $0! = 1$. So the probability of \emph{no event at all} in the interval is simply $e^{-\lambda}$.

\begin{propriete}[The probabilities sum to 1]
For $X \sim \mathrm{Po}(\lambda)$, $\displaystyle\sum_{r=0}^{\infty} P(X = r) = 1$.
\end{propriete}

\textbf{Proof (examinable and instructive).} Recall the exponential series, valid for every real number $x$:
\[ e^{x} = \sum_{r=0}^{\infty} \frac{x^{r}}{r!} = 1 + x + \frac{x^{2}}{2!} + \frac{x^{3}}{3!} + \cdots \]
Then
\[ \sum_{r=0}^{\infty} P(X = r) = \sum_{r=0}^{\infty} \frac{e^{-\lambda}\lambda^{r}}{r!} = e^{-\lambda}\sum_{r=0}^{\infty} \frac{\lambda^{r}}{r!} = e^{-\lambda}\, e^{\lambda} = e^{0} = 1. \]
Moreover each term is positive, so the formula defines a genuine probability distribution. \hfill$\blacksquare$

\begin{exemplebox}[First calculations with $X \sim \mathrm{Po}(3)$]
A switchboard receives on average $3$ calls per ten minutes. Let $X$ be the number of calls in ten minutes, so $X \sim \mathrm{Po}(3)$.
\begin{align*}
P(X = 0) &= e^{-3} \approx 0.0498,\\
P(X = 2) &= \frac{e^{-3}\, 3^{2}}{2!} = \frac{9}{2}e^{-3} \approx 0.2240,\\
P(X = 5) &= \frac{e^{-3}\, 3^{5}}{5!} = \frac{243}{120}e^{-3} \approx 0.1008.
\end{align*}
\end{exemplebox}

\begin{attention}[The forgotten factorial]
The most frequent error at the GCE is to write $P(X = r) = e^{-\lambda}\lambda^{r}$ and to \textbf{forget the $r!$ in the denominator}. Without $r!$ the probabilities do not sum to 1 and every numerical answer is wrong. A safe habit: write the formula \emph{in full} on your script before substituting any number, and compute $r!$ explicitly ($3! = 6$, $4! = 24$, $5! = 120$).
\end{attention}

\courssub{Expectation and Variance: the Distinctive Feature $\mathrm{E}(X) = \mathrm{Var}(X) = \lambda$}

\begin{propriete}[Mean and variance of the Poisson distribution]
If $X \sim \mathrm{Po}(\lambda)$, then
\[ \mathrm{E}(X) = \lambda \qquad \text{and} \qquad \mathrm{Var}(X) = \lambda. \]
The proof of $\mathrm{E}(X) = \lambda$ is examinable; the variance proof follows the same technique and is also examinable.
\end{propriete}

\textbf{Proof that $\mathrm{E}(X) = \lambda$.} By definition,
\[ \mathrm{E}(X) = \sum_{r=0}^{\infty} r\, P(X = r) = \sum_{r=0}^{\infty} r\, \frac{e^{-\lambda}\lambda^{r}}{r!}. \]
The term with $r = 0$ vanishes, and for $r \geq 1$ we simplify $\dfrac{r}{r!} = \dfrac{1}{(r-1)!}$:
\[ \mathrm{E}(X) = e^{-\lambda}\sum_{r=1}^{\infty} \frac{\lambda^{r}}{(r-1)!} = e^{-\lambda}\,\lambda \sum_{r=1}^{\infty} \frac{\lambda^{r-1}}{(r-1)!}. \]
Putting $s = r - 1$, the sum becomes $\displaystyle\sum_{s=0}^{\infty} \frac{\lambda^{s}}{s!} = e^{\lambda}$. Hence
\[ \mathrm{E}(X) = e^{-\lambda}\,\lambda\, e^{\lambda} = \lambda. \quad\blacksquare \]

\textbf{Proof that $\mathrm{Var}(X) = \lambda$.} We first compute $\mathrm{E}\bigl(X(X-1)\bigr)$:
\[ \mathrm{E}\bigl(X(X-1)\bigr) = \sum_{r=0}^{\infty} r(r-1)\frac{e^{-\lambda}\lambda^{r}}{r!}. \]
Terms for $r=0,1$ are zero. For $r \geq 2$, $\dfrac{r(r-1)}{r!} = \dfrac{1}{(r-2)!}$:
\[ \mathrm{E}\bigl(X(X-1)\bigr) = e^{-\lambda}\sum_{r=2}^{\infty} \frac{\lambda^{r}}{(r-2)!} = e^{-\lambda}\,\lambda^{2} \sum_{r=2}^{\infty} \frac{\lambda^{r-2}}{(r-2)!} = e^{-\lambda}\,\lambda^{2}\, e^{\lambda} = \lambda^{2}. \]
Now $\mathrm{E}(X^{2}) = \mathrm{E}\bigl(X(X-1)\bigr) + \mathrm{E}(X) = \lambda^{2} + \lambda$. Therefore
\[ \mathrm{Var}(X) = \mathrm{E}(X^{2}) - \bigl(\mathrm{E}(X)\bigr)^{2} = \lambda^{2} + \lambda - \lambda^{2} = \lambda. \quad\blacksquare \]

\begin{aretenir}[A quick diagnostic for Poisson data]
The equality $\text{mean} = \text{variance}$ is the \emph{signature} of the Poisson distribution. If you count events in many equal intervals and find a sample mean of about $4.1$ and a sample variance of about $4.0$, a Poisson model with $\lambda \approx 4$ is plausible. If the variance is much \emph{larger} than the mean (over-dispersion) or much \emph{smaller}, the Poisson model should be rejected. This is the first test used when \emph{fitting} a distribution to real data.
\end{aretenir}

\courssub{Recurrence Formula and Mode of the Poisson Distribution}

\begin{propriete}[Recurrence formula]
For $X \sim \mathrm{Po}(\lambda)$ with $\lambda > 0$, the probabilities satisfy
\[ P(X = r) = \frac{\lambda}{r}\, P(X = r-1), \qquad r = 1, 2, 3, \dots \]
with $P(X = 0) = e^{-\lambda}$.
\end{propriete}

\textbf{Proof.}
\[ \frac{P(X = r)}{P(X = r-1)} = \frac{e^{-\lambda}\lambda^{r}/r!}{e^{-\lambda}\lambda^{r-1}/(r-1)!} = \frac{\lambda}{r}. \quad\blacksquare \]

This recurrence allows rapid computation of successive probabilities without recalculating factorials and powers each time.

\begin{propriete}[Mode of the Poisson distribution]
The \cle{mode} (most probable value) of $X \sim \mathrm{Po}(\lambda)$ is:
\begin{itemize}
\item $\lambda$ and $\lambda - 1$ (two modes) if $\lambda$ is a positive integer;
\item $\lfloor \lambda \rfloor$ (the greatest integer less than $\lambda$) if $\lambda$ is not an integer.
\end{itemize}
\end{propriete}

\textbf{Justification.} From the recurrence, $P(X = r) \geq P(X = r-1)$ iff $\frac{\lambda}{r} \geq 1$, i.e. $r \leq \lambda$. So probabilities increase while $r < \lambda$, are equal at $r = \lambda$ (if $\lambda$ is integer), and decrease for $r > \lambda$.

\begin{exemplebox}[Using the recurrence and finding the mode]
Let $X \sim \mathrm{Po}(4.5)$.
\begin{enumerate}
\item Compute $P(X = 0), P(X = 1), P(X = 2), P(X = 3)$ using the recurrence.
\item State the mode of the distribution.
\end{enumerate}
\tcblower
\textbf{Solution.}
\begin{enumerate}
\item $P(X = 0) = e^{-4.5} \approx 0.0111$.
$P(X = 1) = \frac{4.5}{1} \times 0.0111 = 0.0500$.
$P(X = 2) = \frac{4.5}{2} \times 0.0500 = 0.1125$.
$P(X = 3) = \frac{4.5}{3} \times 0.1125 = 0.1688$.
\item Since $\lambda = 4.5$ is not an integer, the mode is $\lfloor 4.5 \rfloor = 4$.
\end{enumerate}
\end{exemplebox}

\courssub{The Shape of the Distribution as $\lambda$ Grows}

For small $\lambda$, the Poisson distribution is strongly skewed to the right: small values of $X$ dominate. As $\lambda$ increases, the distribution spreads out and becomes more and more symmetric, looking increasingly like a bell shape centred near $\lambda$.

\begin{popfigure}
\begin{center}
\begin{tikzpicture}
\begin{axis}[
  width=0.33\linewidth, height=5.2cm,
  ybar, bar width=4pt,
  ymin=0, ymax=0.45, xmin=-0.6, xmax=8.6,
  xlabel={$r$}, ylabel={$P(X=r)$},
  title={$\mathrm{Po}(1)$},
  xtick={0,2,4,6,8},
  tick label style={font=\scriptsize},
  label style={font=\scriptsize},
  title style={font=\small},
]
\addplot+[popblue, fill=popblueL] coordinates {(0,0.3679) (1,0.3679) (2,0.1839) (3,0.0613) (4,0.0153) (5,0.0031) (6,0.0005) (7,0.0001) (8,0.0000)};
\end{axis}
\end{tikzpicture}\hfill
\begin{tikzpicture}
\begin{axis}[
  width=0.33\linewidth, height=5.2cm,
  ybar, bar width=4pt,
  ymin=0, ymax=0.45, xmin=-0.6, xmax=12.6,
  xlabel={$r$},
  title={$\mathrm{Po}(4)$},
  xtick={0,3,6,9,12},
  tick label style={font=\scriptsize},
  label style={font=\scriptsize},
  title style={font=\small},
]
\addplot+[poporange, fill=popgold] coordinates {(0,0.0183) (1,0.0733) (2,0.1465) (3,0.1954) (4,0.1954) (5,0.1563) (6,0.1042) (7,0.0595) (8,0.0298) (9,0.0132) (10,0.0053) (11,0.0019) (12,0.0006)};
\end{axis}
\end{tikzpicture}\hfill
\begin{tikzpicture}
\begin{axis}[
  width=0.33\linewidth, height=5.2cm,
  ybar, bar width=4pt,
  ymin=0, ymax=0.45, xmin=-0.6, xmax=19.6,
  xlabel={$r$},
  title={$\mathrm{Po}(9)$},
  xtick={0,5,10,15},
  tick label style={font=\scriptsize},
  label style={font=\scriptsize},
  title style={font=\small},
]
\addplot+[popgreen, fill=popgreenL] coordinates {(0,0.0001) (1,0.0011) (2,0.0050) (3,0.0150) (4,0.0337) (5,0.0607) (6,0.0911) (7,0.1171) (8,0.1318) (9,0.1318) (10,0.1186) (11,0.0970) (12,0.0728) (13,0.0504) (14,0.0324) (15,0.0194) (16,0.0109) (17,0.0058) (18,0.0029) (19,0.0014)};
\end{axis}
\end{tikzpicture}
\legende{Bar charts of $\mathrm{Po}(1)$, $\mathrm{Po}(4)$ and $\mathrm{Po}(9)$: as $\lambda$ grows, the distribution moves right, spreads out and becomes more symmetric.}
\end{center}
\end{popfigure}

\courssub{Uses of the Poisson Distribution}

The Poisson distribution is the standard model for \emph{rare or random events in a continuum}:
\begin{itemize}
\item \textbf{Telecommunications:} calls arriving at an exchange, text messages per minute at a network node;
\item \textbf{Traffic:} vehicles passing a checkpoint on the Douala--Yaoundé road per minute in light traffic, accidents per month at a junction;
\item \textbf{Biology and medicine:} bacteria per drop of water under a microscope, mutations per generation, patients arriving at the emergency ward of the Laquintinie Hospital per hour;
\item \textbf{Insurance:} claims of a given type per week received by an insurance company;
\item \textbf{Quality control:} defects per metre of cloth, misprints per page.
\end{itemize}

\begin{exoresolu}[Accidents on the highway]
On a certain stretch of the Douala--Yaoundé highway, accidents occur randomly, independently, singly and at a constant average rate of $2$ accidents per week. Let $X$ be the number of accidents in a given week.
\begin{enumerate}
\item State the distribution of $X$ and give $\mathrm{E}(X)$ and $\mathrm{Var}(X)$.
\item Find the probability that in a given week there are (i) no accidents; (ii) exactly $3$ accidents.
\item Find the probability of \emph{at least one} accident in a week.
\end{enumerate}
\tcblower
\textbf{Solution.}
\begin{enumerate}
\item The four Poisson conditions are given, with rate $\lambda = 2$ per week, so $X \sim \mathrm{Po}(2)$. Hence $\mathrm{E}(X) = 2$ and $\mathrm{Var}(X) = 2$.
\item Using $P(X = r) = \dfrac{e^{-2}\,2^{r}}{r!}$:
\begin{enumerate}
\item[(i)] $P(X = 0) = e^{-2} \approx 0.1353$.
\item[(ii)] $P(X = 3) = \dfrac{e^{-2}\,2^{3}}{3!} = \dfrac{8}{6}e^{-2} \approx 0.1804$.
\end{enumerate}
\item Using the complement:
\[ P(X \geq 1) = 1 - P(X = 0) = 1 - e^{-2} \approx 1 - 0.1353 = 0.8647. \]
\end{enumerate}
So there is about an $86.5\%$ chance of at least one accident in any given week on this stretch.
\end{exoresolu}

\begin{exercice}[Misprints in a textbook]
Misprints in a mathematics textbook occur randomly and independently at an average rate of $1.5$ misprints per page. A page is chosen at random. Find the probability that it contains (a) no misprint; (b) exactly two misprints; (c) more than two misprints. Give your answers to 4 decimal places.
\end{exercice}

\begin{corrige}[Exercise 1]
Let $X \sim \mathrm{Po}(1.5)$, so $P(X = r) = \dfrac{e^{-1.5}(1.5)^{r}}{r!}$ and $e^{-1.5} \approx 0.2231$.
\begin{itemize}
\item (a) $P(X = 0) = e^{-1.5} \approx 0.2231$.
\item (b) $P(X = 2) = \dfrac{e^{-1.5}(1.5)^{2}}{2!} = 1.125\,e^{-1.5} \approx 0.2510$.
\item (c) $P(X = 1) = 1.5\,e^{-1.5} \approx 0.3347$, so
\[ P(X > 2) = 1 - \bigl(P(0) + P(1) + P(2)\bigr) \approx 1 - (0.2231 + 0.3347 + 0.2510) = 0.1912. \]
\end{itemize}
\end{corrige}

\begin{savaistu}[A law born from miscarriages of justice]
Siméon Denis Poisson introduced this distribution in 1837 in a work on the probability of judgements in criminal cases --- he wanted to model the (rare) wrongful convictions in French courts. The distribution later became famous through a striking data set: the number of Prussian cavalry soldiers killed each year by horse kicks, recorded over 20 years in several regiments by Ladislaus Bortkiewicz (1898). The observed frequencies matched a Poisson distribution remarkably well, convincing statisticians of its power to describe rare, random events.
\end{savaistu}

\begin{aretenir}[Key points of this section]
\begin{itemize}
\item A Poisson process requires events that are \cle{random}, \cle{independent}, \cle{single} and occur at a \cle{constant average rate} $\lambda$ per \cle{unit interval}.
\item $X \sim \mathrm{Po}(\lambda)$: $\;P(X = r) = \dfrac{e^{-\lambda}\lambda^{r}}{r!}$, $r = 0, 1, 2, \dots$; in particular $P(X = 0) = e^{-\lambda}$.
\item The probabilities sum to 1 because $\sum \lambda^{r}/r! = e^{\lambda}$.
\item $\mathrm{E}(X) = \mathrm{Var}(X) = \lambda$: mean equals variance, the Poisson signature.
\item Recurrence: $P(X = r) = \frac{\lambda}{r} P(X = r-1)$ for $r \ge 1$.
\item Mode: $\lambda$ and $\lambda-1$ if $\lambda \in \mathbb{N}^*$; otherwise $\lfloor \lambda \rfloor$.
\item Never forget the $r!$ in the denominator, and use $P(X \geq 1) = 1 - e^{-\lambda}$ for ``at least one event''.
\end{itemize}
\end{aretenir}

\coursec{Random Events, Unit Interval, Recurrence Formula and Mode of the Poisson Distribution}

In the previous section we defined the Poisson distribution $\mathrm{Po}(\lambda)$, established its probability mass function
\[
P(X = r) = \frac{e^{-\lambda}\lambda^{r}}{r!}, \qquad r = 0, 1, 2, \ldots
\]
and proved that both its expectation and its variance equal $\lambda$. We now go deeper into the \emph{structure} of the Poisson model: what ``random events'' really look like, how the parameter $\lambda$ behaves when we change the interval of observation, a powerful recurrence trick for computing probabilities quickly, and how to locate the \cle{mode} of the distribution. These ideas are heavily examined at GCE Advanced Level.

\courssub{The Distribution of Random Events}

\begin{experience}[Are random events regularly spaced?]
Suppose calls arrive at an MTN call centre in Douala at an average rate of $4$ calls per minute. Many students imagine that the calls arrive ``evenly'': roughly one call every $15$ seconds. This is wrong. Divide the observation period into blocks of $6$ seconds and record, for each block, how many calls arrived. A typical simulation of a Poisson process with mean $4$ calls per minute (i.e.\ mean $0.4$ call per $6$-second block) gives a pattern like:
\begin{center}
\begin{tabular}{|l|c|c|c|c|c|c|c|c|c|c|}
\hline
\rowcolor{popblueL} Second block & 1--6 & 7--12 & 13--18 & 19--24 & 25--30 & 31--36 & 37--42 & 43--48 & 49--54 & 55--60 \\
\hline
Calls & 1 & 0 & 2 & 0 & 0 & 1 & 0 & 3 & 1 & 0 \\
\hline
\end{tabular}
\end{center}
The total is $8$ calls in this simulated stretch of $2$ minutes (an average of $4$ per minute), but the arrivals are \emph{clustered and irregular}: sometimes three calls arrive within six seconds, sometimes nothing happens for a long stretch.
\end{experience}

The lesson is fundamental: a \cle{constant average rate} does \textbf{not} mean \cle{regular spacing}. The Poisson model assumes that events occur:

\begin{itemize}
\item \textbf{singly} — two events cannot occur at exactly the same instant;
\item \textbf{independently} — the occurrence of one event does not affect the probability of the next;
\item \textbf{at a constant average rate} $\lambda$ per unit interval — the rate does not change over time or space;
\item \textbf{randomly} — the numbers of events in disjoint intervals are independent random variables.
\end{itemize}

Under these conditions, the number $X$ of events in a fixed interval follows $\mathrm{Po}(\lambda)$. The irregularity we observe (gaps and clusters) is precisely what randomness looks like; regular spacing would actually \emph{contradict} the Poisson model.

\begin{attention}[Uniform rate $\neq$ uniform spacing]
Saying ``events occur at a uniform rate'' means the \emph{average} rate $\lambda$ is constant. It does \textbf{not} mean events are evenly spaced. Examiners love the question: ``State the conditions for a Poisson model.'' The answer is: events occur singly, independently, randomly, and at a constant \emph{mean} rate.
\end{attention}

\courssub{The Unit Interval: Scaling the Parameter}

A Poisson variable is always defined \emph{relative to an interval}: per minute, per hour, per metre, per page. If we change the interval, the parameter must change proportionally. Intuitively: if a road has on average $2$ potholes per kilometre, then a $5$ km stretch has on average $10$ potholes.

\begin{propriete}[Scaling rule for the Poisson parameter]
If $X \sim \mathrm{Po}(\lambda)$ counts events per \textbf{unit} interval, then the number $Y$ of events in an interval of length $t$ units satisfies
\[
Y \sim \mathrm{Po}(\lambda t).
\]
The Poisson rate is \cle{proportional} to the length (or duration, or area) of the interval. (Proof not examinable; it follows from the independence of events in disjoint intervals.)
\end{propriete}

\begin{popfigure}
\begin{tikzpicture}[scale=1.0]
\draw[thick, popink, ->] (0,0) -- (12.6,0) node[below left=2pt and -30pt, popink]{\small time};
\foreach \x/\lab in {0/0, 2/1, 4/2, 6/3, 8/4, 10/5, 12/6}{
  \draw[popink] (\x,0.12) -- (\x,-0.12);
  \node[below, popink] at (\x,-0.15) {\small $\lab$ min};
}
\foreach \x in {1,2,3,4,5,6}{
  \node[popblue] at (2*\x-1,0.55) {\small $\mathrm{Po}(\lambda)$};
  \draw[popblue, thick] (2*\x-1.75,0.32) -- (2*\x-0.25,0.32);
}
\draw[poporange, very thick, decorate, decoration={brace, amplitude=6pt, mirror}] (0.05,-1.05) -- (5.95,-1.05);
\node[poporange, align=center] at (3,-1.7) {\small $3$-minute block: $\mathrm{Po}(3\lambda)$};
\draw[popgreen, very thick, decorate, decoration={brace, amplitude=6pt, mirror}] (6.05,-1.05) -- (11.95,-1.05);
\node[popgreen, align=center] at (9,-1.7) {\small next $3$ minutes: $\mathrm{Po}(3\lambda)$};
\end{tikzpicture}
\legende{A time axis split into unit (1-minute) intervals, each carrying a $\mathrm{Po}(\lambda)$ count. Over $t$ consecutive units the counts add, giving $\mathrm{Po}(t\lambda)$.}
\end{popfigure}

\begin{exemplebox}[Changing the unit interval]
Calls arrive at an office at an average rate of $3$ per minute, following a Poisson distribution.
\begin{itemize}
\item Per minute: $X \sim \mathrm{Po}(3)$.
\item Per $5$ minutes: $Y \sim \mathrm{Po}(3 \times 5) = \mathrm{Po}(15)$.
\item Per hour: $Z \sim \mathrm{Po}(3 \times 60) = \mathrm{Po}(180)$.
\item Per $30$ seconds: $W \sim \mathrm{Po}(3 \times 0.5) = \mathrm{Po}(1.5)$.
\end{itemize}
Similarly, if defects occur along a cable at $0.4$ per metre, then per $10$ metres the count is $\mathrm{Po}(0.4 \times 10) = \mathrm{Po}(4)$, and per $25$ cm it is $\mathrm{Po}(0.4 \times 0.25) = \mathrm{Po}(0.1)$.
\end{exemplebox}

\begin{attention}[The classic trap: forgetting to scale $\lambda$]
If the question gives a rate \emph{per minute} but asks about a \emph{five-minute} period, you must use $\lambda = 5 \times$ (rate per minute). Using the original $\lambda$ with the new interval is the single most common error in Poisson questions. Always write, as your first line: ``In a $\ldots$ interval, $\lambda = \ldots$''.
\end{attention}

\begin{propriete}[Sum of independent Poisson variables]
If $X \sim \mathrm{Po}(\lambda_1)$ and $Y \sim \mathrm{Po}(\lambda_2)$ are \textbf{independent}, then
\[
X + Y \sim \mathrm{Po}(\lambda_1 + \lambda_2).
\]
\end{propriete}

\textbf{Justification (instructive).} Using independence,
\begin{align*}
P(X+Y = n) &= \sum_{r=0}^{n} P(X=r)\,P(Y=n-r)
= \sum_{r=0}^{n} \frac{e^{-\lambda_1}\lambda_1^{r}}{r!}\cdot\frac{e^{-\lambda_2}\lambda_2^{\,n-r}}{(n-r)!} \\
&= \frac{e^{-(\lambda_1+\lambda_2)}}{n!}\sum_{r=0}^{n} \frac{n!}{r!(n-r)!}\,\lambda_1^{r}\lambda_2^{\,n-r}
= \frac{e^{-(\lambda_1+\lambda_2)}}{n!}\sum_{r=0}^{n} \binom{n}{r}\lambda_1^{r}\lambda_2^{\,n-r} \\
&= \frac{e^{-(\lambda_1+\lambda_2)}(\lambda_1+\lambda_2)^{n}}{n!},
\end{align*}
by the binomial theorem. This is exactly the $\mathrm{Po}(\lambda_1+\lambda_2)$ probability mass function. $\blacksquare$

\begin{exemplebox}[Combining two Poisson sources]
A hospital receives emergency cases from Bonaberi at $2$ per hour and from Deido at $3$ per hour, independently. The total number of emergencies per hour is $T \sim \mathrm{Po}(2+3) = \mathrm{Po}(5)$. Hence
\[
P(T = 0) = e^{-5} \approx 0.0067, \qquad P(T \geq 1) = 1 - e^{-5} \approx 0.9933.
\]
\end{exemplebox}

\courssub{The Recurrence Formula}

Computing $P(X=r) = \dfrac{e^{-\lambda}\lambda^{r}}{r!}$ from scratch for each $r$ is tedious. Fortunately, consecutive probabilities are linked by a simple ratio.

\begin{propriete}[Recurrence formula for the Poisson distribution]
If $X \sim \mathrm{Po}(\lambda)$, then for all $r \geq 0$,
\[
P(X = r+1) = \frac{\lambda}{r+1}\,P(X = r).
\]
(The derivation is examinable and takes two lines.)
\end{propriete}

\textbf{Derivation.}
\[
\frac{P(X=r+1)}{P(X=r)} = \frac{e^{-\lambda}\lambda^{r+1}/(r+1)!}{e^{-\lambda}\lambda^{r}/r!} = \frac{\lambda^{r+1}}{\lambda^{r}}\cdot\frac{r!}{(r+1)!} = \lambda\cdot\frac{1}{r+1} = \frac{\lambda}{r+1}. \qquad \blacksquare
\]

\begin{methode}[Generating a Poisson probability table by recurrence]
\begin{enumerate}
\item Compute the starting value $P(X=0) = e^{-\lambda}$.
\item Multiply successively: $P(1) = \lambda P(0)$, then $P(2) = \dfrac{\lambda}{2}P(1)$, then $P(3) = \dfrac{\lambda}{3}P(2)$, and so on.
\item Stop when probabilities become negligible (less than $0.0001$, say).
\item Check: the probabilities should sum to $1$ (up to rounding).
\end{enumerate}
\end{methode}

\begin{exoresolu}[Building the table of $\mathrm{Po}(4)$]
Let $X \sim \mathrm{Po}(4)$. Using the recurrence formula, find $P(X=r)$ for $r = 0, 1, \ldots, 8$, giving answers to $4$ decimal places.
\tcblower
\textbf{Solution.} Start with $P(0) = e^{-4} = 0.01832$. Then apply $P(r+1) = \dfrac{4}{r+1}P(r)$:
\begin{align*}
P(1) &= \tfrac{4}{1}(0.01832) = 0.07326, &
P(2) &= \tfrac{4}{2}(0.07326) = 0.14653, \\
P(3) &= \tfrac{4}{3}(0.14653) = 0.19537, &
P(4) &= \tfrac{4}{4}(0.19537) = 0.19537, \\
P(5) &= \tfrac{4}{5}(0.19537) = 0.15629, &
P(6) &= \tfrac{4}{6}(0.15629) = 0.10420, \\
P(7) &= \tfrac{4}{7}(0.10420) = 0.05954, &
P(8) &= \tfrac{4}{8}(0.05954) = 0.02977.
\end{align*}
Notice that $P(3) = P(4)$: the probabilities stop increasing exactly at $r = 4$. This is no accident — it reveals the \emph{mode}, which we study next.
\end{exoresolu}

\courssub{The Mode of the Poisson Distribution}

The \cle{mode} is the value of $r$ with the greatest probability — the \emph{most likely} number of events. The recurrence formula tells us exactly where the probabilities rise and where they fall:
\[
P(X=r+1) \geq P(X=r) \iff \frac{\lambda}{r+1} \geq 1 \iff r+1 \leq \lambda.
\]
So the probabilities \textbf{increase} while $r + 1 < \lambda$ and \textbf{decrease} once $r + 1 > \lambda$. Two cases arise.

\begin{propriete}[Mode of the Poisson distribution]
Let $X \sim \mathrm{Po}(\lambda)$.
\begin{itemize}
\item If $\lambda$ is \textbf{not} an integer, there is a unique mode at $r = \lfloor \lambda \rfloor$ (the greatest integer not exceeding $\lambda$).
\item If $\lambda$ \textbf{is} an integer, there are \textbf{two equal modes}, at $r = \lambda - 1$ and $r = \lambda$, with $P(X=\lambda-1) = P(X=\lambda)$.
\end{itemize}
(Proof examinable: it is the inequality $r+1 \lesseqgtr \lambda$ above.)
\end{propriete}

Indeed, when $\lambda$ is an integer, $P(X=\lambda) = \dfrac{\lambda}{\lambda}P(X=\lambda-1) = P(X=\lambda-1)$, giving the double mode.

\begin{methode}[Finding the mode in an exam]
\begin{enumerate}
\item Look at $\lambda$.
\item If $\lambda$ is an integer: the modes are $\lambda - 1$ and $\lambda$.
\item If $\lambda$ is not an integer: the mode is $\lfloor \lambda \rfloor$.
\item To \emph{justify} the answer (when asked), compare ratios: show $\dfrac{P(X=r+1)}{P(X=r)} = \dfrac{\lambda}{r+1}$ is $\geq 1$ up to the mode and $\leq 1$ after it.
\end{enumerate}
\end{methode}

\begin{popfigure}
\begin{center}
\begin{tikzpicture}
\begin{axis}[
    ybar, bar width=10pt,
    width=12.5cm, height=7cm,
    xlabel={$r$}, ylabel={$P(X=r)$},
    xlabel style={popink}, ylabel style={popink},
    xmin=-0.7, xmax=11.7, ymin=0, ymax=0.23,
    xtick={0,1,2,3,4,5,6,7,8,9,10,11},
    tick label style={popink},
    axis lines=left, axis line style={popline},
    every axis plot/.append style={fill=popblue, draw=popdark},
]
\addplot+ coordinates {(0,0.0183) (1,0.0733) (2,0.1465) (5,0.1563) (6,0.1042) (7,0.0595) (8,0.0298) (9,0.0132) (10,0.0053) (11,0.0019)};
\addplot+[fill=poporange, draw=popdark] coordinates {(3,0.1954) (4,0.1954)};
\node[poporange, align=center] at (axis cs:6.6,0.205) {\small modes: $r=3$ and $r=4$};
\draw[poporange, ->] (axis cs:6.2,0.198) -- (axis cs:4.2,0.196);
\end{axis}
\end{tikzpicture}
\end{center}
\legende{Bar chart of $\mathrm{Po}(4)$. Since $\lambda = 4$ is an integer, the two modes $r = 3$ and $r = 4$ (highlighted) share the greatest probability, $0.1954$.}
\end{popfigure}

\begin{exoresolu}[Locating the most likely number of events]
The number of accidents per week on a stretch of the Douala--Yaound\'e highway follows a Poisson distribution with mean $2.5$.
\begin{enumerate}
\item[(a)] Find the most likely number of accidents in a week, justifying your answer.
\item[(b)] Find the most likely number of accidents in a \emph{fortnight} (two weeks).
\end{enumerate}
\tcblower
\textbf{Solution.}
(a) Here $\lambda = 2.5$, not an integer, so the mode is $\lfloor 2.5 \rfloor = 2$. Justification using the recurrence ratio:
\[
\frac{P(X=r+1)}{P(X=r)} = \frac{2.5}{r+1} \geq 1 \iff r+1 \leq 2.5 \iff r \leq 1.5.
\]
So probabilities increase for $r = 0 \to 1 \to 2$ and decrease from $r = 2$ onwards: the mode is $r = 2$, with $P(X=2) = \dfrac{e^{-2.5}(2.5)^2}{2!} \approx 0.2565$.

(b) In two weeks the parameter scales: $Y \sim \mathrm{Po}(2 \times 2.5) = \mathrm{Po}(5)$. Since $5$ is an integer, there are \textbf{two modes}: $r = 4$ and $r = 5$, with
\[
P(Y=4) = P(Y=5) = \frac{e^{-5}5^{4}}{4!} \approx 0.1755.
\]
Note how the change of unit interval turns a single-mode problem into a double-mode problem.
\end{exoresolu}

\begin{exercice}[Practice]
\begin{enumerate}
\item Bacteria occur in a water sample at an average rate of $1.5$ per millilitre. Write down the distribution of the number of bacteria in (a) $1$ ml, (b) $4$ ml, (c) $200$ ml.
\item Given $X \sim \mathrm{Po}(6)$, use the recurrence formula to compute $P(X=r)$ for $r = 0, 1, \ldots, 7$ (take $e^{-6} = 0.002479$), and state the mode(s).
\item The number of goals scored per match by a team is $\mathrm{Po}(1.8)$, independently from match to match. Find the distribution of the total number of goals in $3$ matches and its most likely value.
\end{enumerate}
\end{exercice}

\begin{corrige}[Exercise 1]
(a) $\mathrm{Po}(1.5)$; (b) $\mathrm{Po}(1.5 \times 4) = \mathrm{Po}(6)$; (c) $\mathrm{Po}(1.5 \times 200) = \mathrm{Po}(300)$.
\end{corrige}

\begin{corrige}[Exercise 2]
$P(0) = 0.002479$; then $P(1) = 6P(0) = 0.014874$, $P(2) = 3P(1) = 0.044622$, $P(3) = 2P(2) = 0.089244$, $P(4) = 1.5P(3) = 0.133866$, $P(5) = 1.2P(4) = 0.160639$, $P(6) = 1 \times P(5) = 0.160639$, $P(7) = \tfrac{6}{7}P(6) = 0.137691$. Since $\lambda = 6$ is an integer, the modes are $r = 5$ and $r = 6$ (equal probabilities, $\approx 0.1606$).
\end{corrige}

\begin{corrige}[Exercise 3]
Total goals $T \sim \mathrm{Po}(3 \times 1.8) = \mathrm{Po}(5.4)$. Since $5.4$ is not an integer, the unique mode is $\lfloor 5.4 \rfloor = 5$ goals.
\end{corrige}

\begin{aretenir}[Key points of this section]
\begin{itemize}
\item A constant \emph{average} rate does not imply regular spacing: Poisson events occur singly, independently and randomly, producing natural gaps and clusters.
\item Scaling: if $X \sim \mathrm{Po}(\lambda)$ per unit interval, then over $t$ units the count is $\mathrm{Po}(\lambda t)$. \textbf{Always rescale $\lambda$ first.}
\item Sums: independent $X \sim \mathrm{Po}(\lambda_1)$, $Y \sim \mathrm{Po}(\lambda_2)$ give $X+Y \sim \mathrm{Po}(\lambda_1+\lambda_2)$.
\item Recurrence: $P(X=r+1) = \dfrac{\lambda}{r+1}P(X=r)$, starting from $P(X=0) = e^{-\lambda}$.
\item Mode: $\lfloor \lambda \rfloor$ if $\lambda \notin \mathbb{Z}$; double mode at $\lambda - 1$ and $\lambda$ if $\lambda \in \mathbb{Z}$.
\end{itemize}
\end{aretenir}

\coursec{Poisson Approximation to the Binomial and Applications of the Poisson Distribution}

In the previous section we studied the Poisson distribution in its own right: how it models random events in a unit interval, its recurrence formula $P(X=k+1)=\dfrac{\lambda}{k+1}P(X=k)$ and its mode. We now close the loop by answering a natural question: \emph{where does the Poisson distribution come from?} The answer is beautiful and extremely useful at the GCE: the Poisson distribution is what the binomial distribution \emph{becomes} when trials are very many but successes are very rare. This gives us a powerful shortcut, and it also explains why the Poisson model fits so many real-life phenomena, from machine breakdowns in a Douala factory to goals scored in an Elite One football match.

\courssub{Motivation: When Binomial Computations Become Heavy}

Suppose a machine produces $n=200$ items per day, each defective independently with probability $p=0.02$, and we want the probability of exactly $5$ defectives. The exact binomial answer is
\[
P(X=5)=\binom{200}{5}(0.02)^5(0.98)^{195}.
\]
The binomial coefficient $\binom{200}{5}$ has $13$ digits, and $(0.98)^{195}$ must be handled with logarithms or a calculator. For $n=1000$ the situation is worse. Yet notice something: the \emph{mean} $np=200\times 0.02=4$ is a small, comfortable number. Although $200$ items are produced, only about $4$ defectives are expected. The action happens entirely in the small values $k=0,1,2,\dots,12$; the probabilities for large $k$ are negligible.

This suggests an idea: when $n$ is large and $p$ is small, perhaps the complicated binomial formula can be replaced by a simpler law that depends only on the mean $\lambda=np$. That simpler law is precisely the Poisson distribution.

\begin{propriete}[Poisson approximation to the binomial]
If $X\sim B(n,p)$ where $n$ is \cle{large} and $p$ is \cle{small}, so that $\lambda=np$ is moderate, then
\[
B(n,p)\approx \mathrm{Po}(\lambda),\qquad\text{i.e.}\qquad P(X=k)\approx e^{-\lambda}\frac{\lambda^{k}}{k!},\quad \lambda=np.
\]
The standard working conditions for the GCE Advanced Level examination are $n\geq 50$ and $p\leq 0.1$ (or $q=1-p\leq 0.1$ if we count failures instead of successes).
\end{propriete}

\textbf{Why the approximation works (matching the means).}\quad A rigorous limit proof is beyond the syllabus, but the key idea is instructive and easy to grasp. Fix $\lambda=np$ and write $p=\lambda/n$. Then
\begin{align*}
P(X=k)&=\binom{n}{k}p^{k}(1-p)^{n-k}\\
&=\frac{n(n-1)\cdots(n-k+1)}{k!}\left(\frac{\lambda}{n}\right)^{k}\left(1-\frac{\lambda}{n}\right)^{n-k}\\
&=\frac{\lambda^{k}}{k!}\cdot\underbrace{\frac{n(n-1)\cdots(n-k+1)}{n^{k}}}_{\to\,1\text{ as }n\to\infty}\cdot\underbrace{\left(1-\frac{\lambda}{n}\right)^{n}}_{\to\,e^{-\lambda}}\cdot\underbrace{\left(1-\frac{\lambda}{n}\right)^{-k}}_{\to\,1}.
\end{align*}
As $n\to\infty$ with $\lambda$ fixed, the three bracketed factors tend to $1$, $e^{-\lambda}$ and $1$ respectively, leaving exactly the Poisson formula $e^{-\lambda}\lambda^{k}/k!$. Notice also that the two distributions \emph{agree on the mean}: $E(X)=np=\lambda$ in both cases, and the binomial variance $np(1-p)=\lambda(1-p)\approx\lambda$ when $p$ is small, matching the Poisson variance. The approximation is therefore consistent in mean \emph{and} in variance.

\begin{methode}[Approximation checklist]
To approximate $X\sim B(n,p)$ by a Poisson distribution:
\begin{enumerate}
\item \textbf{Check} $n\geq 50$ (large number of trials).
\item \textbf{Check} $p\leq 0.1$ (success is rare). Equivalently, if failures are rare ($q=1-p\leq 0.1$), work with $q$ and count failures.
\item \textbf{Compute} $\lambda=np$ and state $X\approx Y$ where $Y\sim \mathrm{Po}(np)$.
\item \textbf{Compute} Poisson probabilities using $P(Y=k)=e^{-\lambda}\lambda^{k}/k!$, the recurrence formula $P(Y=k+1)=\frac{\lambda}{k+1}P(Y=k)$, or cumulative tables.
\end{enumerate}
\end{methode}

\begin{attention}[When NOT to use the Poisson approximation]
A very common mistake is to approximate $B(n,p)$ by a Poisson law when $p$ is close to $0.5$. For example $B(100,0.5)$ is \emph{not} rare-event territory: its mean is $50$ and its spread is large, and $\mathrm{Po}(50)$ fits it poorly (the binomial is nearly symmetric, the Poisson is right-skewed). For large $n$ with moderate $p$, the appropriate approximation is the \emph{normal} distribution (met later in the course), not the Poisson. Reserve the Poisson approximation for \cle{rare} events: $p$ small, $\lambda=np$ moderate.
\end{attention}

\courssub{Comparing Exact and Approximate Values}

\begin{exemplebox}[B(100, 0.04) versus Po(4)]
Let $X\sim B(100,0.04)$: $n=100\geq 50$, $p=0.04\leq 0.1$, so the approximation with $\lambda=np=4$ is legitimate. The table compares exact binomial probabilities (computed from the binomial formula, rounded to 4 decimal places) with the Poisson approximations $e^{-4}4^{k}/k!$.
\begin{center}
\begin{tabular}{|c|c|c|c|}
\hline
\rowcolor{popblueL} $k$ & $B(100,0.04)$ exact & $\mathrm{Po}(4)$ approx. & Absolute error \\
\hline
0 & 0.0169 & 0.0183 & 0.0014 \\
1 & 0.0703 & 0.0733 & 0.0030 \\
2 & 0.1450 & 0.1465 & 0.0015 \\
3 & 0.1973 & 0.1954 & 0.0019 \\
4 & 0.1994 & 0.1954 & 0.0040 \\
5 & 0.1595 & 0.1563 & 0.0032 \\
6 & 0.1053 & 0.1042 & 0.0011 \\
7 & 0.0589 & 0.0595 & 0.0006 \\
8 & 0.0285 & 0.0298 & 0.0013 \\
9 & 0.0122 & 0.0132 & 0.0010 \\
10 & 0.0046 & 0.0053 & 0.0007 \\
\hline
\end{tabular}
\end{center}
The agreement is excellent: every error is below $0.005$, and cumulative probabilities agree to two decimal places. The figure below shows the two distributions superimposed.
\end{exemplebox}

\begin{popfigure}
\begin{center}
\begin{tikzpicture}
\begin{axis}[
    width=12cm, height=7cm,
    xlabel={$k$}, ylabel={$P(X=k)$},
    xmin=-0.5, xmax=12.5, ymin=0, ymax=0.23,
    xtick={0,1,2,3,4,5,6,7,8,9,10,11,12},
    legend style={at={(0.97,0.95)}, anchor=north east, font=\small},
    axis lines=left,
]
\addplot+[ybar, bar width=0.35cm, fill=popblueL, draw=popblue] coordinates {
(0,0.0169) (1,0.0703) (2,0.1450) (3,0.1973) (4,0.1994) (5,0.1595)
(6,0.1053) (7,0.0589) (8,0.0285) (9,0.0122) (10,0.0046) (11,0.0016) (12,0.0005)
};
\addlegendentry{$B(100,0.04)$ (exact)}
\addplot+[only marks, mark=*, mark size=2.2pt, color=poporange] coordinates {
(0,0.0183) (1,0.0733) (2,0.1465) (3,0.1954) (4,0.1954) (5,0.1563)
(6,0.1042) (7,0.0595) (8,0.0298) (9,0.0132) (10,0.0053) (11,0.0019) (12,0.0006)
};
\addlegendentry{$\mathrm{Po}(4)$ (approximation)}
\end{axis}
\end{tikzpicture}
\end{center}
\legende{Bars: exact binomial probabilities for $B(100,0.04)$. Dots: Poisson approximation $\mathrm{Po}(4)$. The two are almost indistinguishable.}
\end{popfigure}

\begin{exoresolu}[Quality control in a soap factory]
A factory in Bonab\'eri produces $400$ tablets of soap per batch. Each tablet is defective with probability $0.005$, independently of the others. Using a suitable approximation, find the probability that a batch contains (a) exactly $3$ defectives; (b) at least one defective.
\tcblower
\textbf{Solution.}\quad Let $X$ be the number of defectives, $X\sim B(400,0.005)$. Since $n=400\geq 50$ and $p=0.005\leq 0.1$, we approximate by $Y\sim\mathrm{Po}(\lambda)$ with
\[
\lambda=np=400\times 0.005=2.
\]
(a) $\displaystyle P(X=3)\approx P(Y=3)=e^{-2}\frac{2^{3}}{3!}=0.1353\times\frac{8}{6}\approx 0.1804.$

(b) Using the ``at least one'' trick:
\[
P(X\geq 1)\approx P(Y\geq 1)=1-P(Y=0)=1-e^{-2}\approx 1-0.1353=0.8647.
\]
So a batch contains at least one defective about $86\%$ of the time.
\end{exoresolu}

\courssub{Applications of the Poisson Distribution}

Because rare-event binomial situations are everywhere, the Poisson distribution is one of the most applied laws in statistics. Typical modelling situations (all with $n$ large, $p$ small, events independent, in a fixed interval of time, length, area or volume):

\begin{itemize}
\item \textbf{Rare diseases:} number of cases of a rare illness per year in a region of Cameroon, with $\lambda$ = average annual number of cases.
\item \textbf{Machine breakdowns:} number of breakdowns of a generator per month in a factory; maintenance teams use $P(X\geq k)$ to plan spare parts.
\item \textbf{Typing and printing errors:} number of misprints per page of a textbook; a proofreader expects $\lambda$ errors per page and computes $P(X=0)=e^{-\lambda}$, the probability of a clean page.
\item \textbf{Customer arrivals:} number of customers joining the queue at an MTN service point per $10$ minutes, or calls arriving at a call centre per minute.
\item \textbf{Goals in football:} the number of goals scored by a team in a match is well modelled by a Poisson law with $\lambda$ equal to the team's average goals per match.
\end{itemize}

\begin{savaistu}[A law born from data]
The Poisson distribution was introduced by Sim\'eon Denis Poisson in the 1830s, and became famous when Ladislaus Bortkiewicz showed in 1898 that it fitted the number of soldiers killed by horse kicks per year in Prussian army corps: many corps (large $n$), tiny probability of such a death (small $p$), and a small average $\lambda$. It remains the standard model for rare events in insurance, epidemiology and telecommunications.
\end{savaistu}

\courssub{Multi-Step Problems: Scaling, Cumulative Probabilities and ``At Least One''}

GCE questions often combine three ideas in one problem: the \cle{scaling rule} (if $X\sim\mathrm{Po}(\lambda)$ counts events in one unit interval, then the count in $t$ units is $\mathrm{Po}(\lambda t)$), \cle{cumulative probabilities} $P(X\leq k)$ obtained by summing, and the \cle{at-least-one event} $P(X\geq 1)=1-e^{-\lambda}$.

\begin{exoresolu}[Calls at a call centre]
Calls arrive at a customer-care centre in Yaound\'e at an average rate of $3$ calls per $10$ minutes, independently, following a Poisson process.
\begin{enumerate}
\item Find the probability of exactly $2$ calls in a $10$-minute period.
\item Find the probability of at least one call in a $5$-minute period.
\item Find the probability of at most $2$ calls in a $20$-minute period.
\end{enumerate}
\tcblower
\textbf{Solution.}
\begin{enumerate}
\item Here $\lambda=3$: $\displaystyle P(X=2)=e^{-3}\frac{3^{2}}{2!}=0.04979\times 4.5\approx 0.2240.$
\item Scale the interval: in $5$ minutes, $\lambda=3\times\frac{5}{10}=1.5$. Then
\[
P(X\geq 1)=1-e^{-1.5}\approx 1-0.2231=0.7769.
\]
\item In $20$ minutes, $\lambda=3\times 2=6$. Then
\begin{align*}
P(X\leq 2)&=e^{-6}\left(1+6+\frac{6^{2}}{2!}\right)=e^{-6}(1+6+18)=25e^{-6}\approx 0.0620.
\end{align*}
\end{enumerate}
\end{exoresolu}

\begin{attention}[Scale $\lambda$ before computing]
The most frequent error in such questions is forgetting to adjust $\lambda$ when the interval changes. If the rate is given per $10$ minutes and the question concerns $5$ minutes, you \emph{must} halve $\lambda$ first. Writing $1-e^{-3}$ instead of $1-e^{-1.5}$ loses all the marks for that part.
\end{attention}

\courssub{Using the Recurrence Formula for Efficient Computation}

When calculating several consecutive Poisson probabilities, the recurrence formula avoids repeated evaluation of $e^{-\lambda}$ and factorials.

\begin{exemplebox}[Using the recurrence formula]
For the call centre with $\lambda=3$ (per 10 minutes), compute $P(X=0), P(X=1), P(X=2), P(X=3)$ using the recurrence $P(X=k+1)=\frac{\lambda}{k+1}P(X=k)$.
\tcblower
\textbf{Solution.}
\begin{align*}
P(X=0) &= e^{-3} \approx 0.04979. \\
P(X=1) &= \frac{3}{1} \times 0.04979 = 0.14936. \\
P(X=2) &= \frac{3}{2} \times 0.14936 = 0.22404. \\
P(X=3) &= \frac{3}{3} \times 0.22404 = 0.22404.
\end{align*}
Check: $P(X=3) = e^{-3}\frac{27}{6} = 0.04979 \times 4.5 = 0.22404$. The recurrence is faster and reduces rounding errors.
\end{exemplebox}

\courssub{Estimating $\lambda$ from Data and Checking the Fit}

In practice $\lambda$ is unknown and must be \cle{estimated} from data. Since $E(X)=\lambda$ for a Poisson distribution, the natural estimate is the \emph{sample mean}:
\[
\hat{\lambda}=\bar{x}=\frac{\sum f_k\, k}{\sum f_k}.
\]
A quick diagnostic for whether a Poisson model is reasonable exploits the signature property of the law: \textbf{mean $\approx$ variance}. Compute the sample mean $\bar{x}$ and the sample variance $s^{2}$; if they are close, a Poisson fit is plausible. If the variance is much larger than the mean (``overdispersion''), the Poisson model should be rejected.

\begin{exemplebox}[Goals scored per match]
Over $40$ matches, a team scores the following numbers of goals: $0$ goals in $14$ matches, $1$ goal in $14$ matches, $2$ goals in $8$ matches, $3$ goals in $3$ matches, $4$ goals in $1$ match. Then
\[
\hat{\lambda}=\bar{x}=\frac{0(14)+1(14)+2(8)+3(3)+4(1)}{40}=\frac{43}{40}=1.075.
\]
The sample variance is
\[
s^2 = \frac{\sum f_k k^2}{\sum f_k} - \bar{x}^2 = \frac{0+14+32+27+16}{40} - (1.075)^2 = \frac{89}{40} - 1.155625 = 2.225 - 1.155625 = 1.069375 \approx 1.07,
\]
very close to the mean, so $\mathrm{Po}(1.075)$ is a sensible model. For instance the fitted probability of scoring no goal is $e^{-1.075}\approx 0.341$, predicting about $40\times 0.341\approx 13.6$ goalless matches, against $14$ observed: an excellent fit.
\end{exemplebox}

\courssub{Comparing the Binomial, Poisson and Normal Approximations}

Lesson 158 of the syllabus requires comparing distributions. The table below summarises when to use each approximation for a binomial $B(n,p)$.

\begin{center}
\begin{tabularx}{\linewidth}{|l|X|X|X|}
\hline
\rowcolor{popblueL} Situation & Approximation & Conditions & Parameter(s) \\
\hline
$n$ large, $p$ small (rare successes) & \textbf{Poisson} & $n\geq 50,\ p\leq 0.1$ & $\lambda = np$ \\
$n$ large, $q=1-p$ small (rare failures) & \textbf{Poisson} (on failures) & $n\geq 50,\ q\leq 0.1$ & $\lambda = nq$ \\
$n$ large, $p$ moderate (not near 0 or 1) & \textbf{Normal} & $np>5,\ nq>5$ & $\mu=np,\ \sigma^2=npq$ \\
\hline
\end{tabularx}
\end{center}

\begin{attention}[Continuity correction for Normal, not for Poisson]
When using the Normal approximation to the binomial, a \cle{continuity correction} is mandatory (e.g. $P(X\leq k) \approx P(Z \leq \frac{k+0.5-\mu}{\sigma})$). The Poisson approximation is \emph{discrete-to-discrete}: no continuity correction is needed or allowed. Simply compute $P(Y=k)$ or $P(Y\leq k)$ for $Y\sim\mathrm{Po}(np)$.
\end{attention}

\courssub{Real-Life Problem Solving: A Comprehensive Example}

\begin{exoresolu}[Inspection of electrical cables]
A factory produces electrical cables. Experience shows that $0.8\%$ of the cables have a microscopic defect. A quality controller selects a random sample of $500$ cables.
\begin{enumerate}
\item State the exact distribution of the number of defective cables in the sample. Justify why a Poisson approximation is appropriate and give its parameter.
\item Use the Poisson approximation to find the probability that the sample contains (i) no defective cables, (ii) exactly two defective cables, (iii) at most one defective cable.
\item The factory produces $10\,000$ cables per day. Estimate the number of days in a $300$-day working year on which a sample of $500$ cables would contain more than $3$ defectives.
\end{enumerate}
\tcblower
\textbf{Solution.}
\begin{enumerate}
\item Let $X$ be the number of defective cables in the sample. Exact distribution: $X\sim B(500,\,0.008)$. Since $n=500\geq 50$ and $p=0.008\leq 0.1$, the Poisson approximation is valid with $\lambda = np = 500\times 0.008 = 4$. So $X\approx Y\sim\mathrm{Po}(4)$.
\item Using $Y\sim\mathrm{Po}(4)$:
\begin{itemize}
\item[(i)] $P(Y=0)=e^{-4}\approx 0.0183$.
\item[(ii)] $P(Y=2)=e^{-4}\frac{4^2}{2!}=0.0183\times 8 = 0.1465$.
\item[(iii)] $P(Y\leq 1)=P(Y=0)+P(Y=1)=e^{-4}(1+4)=5e^{-4}\approx 5\times 0.0183 = 0.0916$.
\end{itemize}
\item For a daily sample of $500$, $P(\text{more than 3 defectives}) = P(Y>3) = 1-P(Y\leq 3)$.
$P(Y\leq 3) = e^{-4}\left(1+4+\frac{16}{2}+\frac{64}{6}\right) = e^{-4}(1+4+8+10.6667) = 23.6667e^{-4} \approx 0.4335$.
So $P(Y>3) \approx 1-0.4335 = 0.5665$.
Expected number of days in $300$ days $= 300 \times 0.5665 \approx 170$ days.
\end{enumerate}
\end{exoresolu}

\begin{methode}[GCE exam technique: ``use a suitable approximation'']
When a question says \emph{``use a suitable approximation''}:
\begin{enumerate}
\item \textbf{Name} the approximating distribution explicitly: ``$X\sim B(n,p)\approx \mathrm{Po}(\lambda)$''.
\item \textbf{Justify} the conditions: state $n\geq 50$ and $p\leq 0.1$ (or $q\leq 0.1$).
\item \textbf{Give the parameter}: $\lambda=np$, with its numerical value.
\item \textbf{Compute} with the Poisson formula, recurrence, or tables, and round sensibly (usually $4$ decimal places).
\end{enumerate}
Examiners award method marks for the statement of the approximating distribution and its parameter, even before any computation.
\end{methode}

\begin{exercice}[Insurance claims]
An insurance company receives claims from $600$ policyholders; each policyholder makes a claim in a given month with probability $0.004$, independently.
\begin{enumerate}
\item Justify the use of a Poisson approximation and state its parameter.
\item Find the probability that exactly $2$ claims are made in a month.
\item Find the probability that at least $3$ claims are made in a month.
\end{enumerate}
\end{exercice}

\begin{corrige}[Exercise 1]
\begin{enumerate}
\item $n=600\geq 50$, $p=0.004\leq 0.1$: the Poisson approximation is valid with $\lambda=np=600\times 0.004=2.4$, so $X\approx Y\sim\mathrm{Po}(2.4)$.
\item $P(Y=2)=e^{-2.4}\dfrac{2.4^{2}}{2!}=0.0907\times 2.88\approx 0.2613$.
\item $P(Y\geq 3)=1-P(Y\leq 2)=1-e^{-2.4}\left(1+2.4+\dfrac{2.4^{2}}{2}\right)=1-e^{-2.4}(1+2.4+2.88)=1-0.0907\times 6.28\approx 1-0.5697=0.4303.$
\end{enumerate}
\end{corrige}

\begin{aretenir}[Key points]
\begin{itemize}
\item If $n\geq 50$ and $p\leq 0.1$, then $B(n,p)\approx\mathrm{Po}(np)$; set $\lambda=np$ and say so explicitly.
\item The approximation matches the mean exactly ($np=\lambda$) and the variance approximately ($np(1-p)\approx\lambda$).
\item Never use the Poisson approximation when $p$ is near $0.5$; that is normal-approximation territory.
\item ``At least one event'': $P(X\geq 1)=1-e^{-\lambda}$. Change of interval: scale $\lambda$ in proportion to the interval length.
\item Estimate $\lambda$ by the sample mean, and check the Poisson fit by verifying mean $\approx$ variance.
\item Use the recurrence formula $P(X=k+1)=\frac{\lambda}{k+1}P(X=k)$ for fast, accurate computation of several probabilities.
\item No continuity correction for Poisson approximation (discrete to discrete).
\end{itemize}
\end{aretenir}

\coursec{Comparing Distributions and Real-Life Problem Solving}

We have now studied three fundamental discrete probability distributions: the \cle{binomial}, the \cle{geometric} and the \cle{Poisson}. Each arises from a different random experiment, yet they share a common language of parameters, expectations and variances. In this final section we bring them together, develop a systematic method for choosing the right model in a worded problem, and practise the kind of multi-part questions that appear in GCE Paper 2.

\courssub{Summary Comparison of the Three Distributions}

Before solving problems, it is essential to have a clear mental picture of the differences and similarities. The table below summarises the key features. Note that for the Geometric distribution, we adopt the convention $X =$ number of trials up to and including the first success (support $x=1,2,\dots$), which is the standard for the GCE Advanced Level syllabus.

\begin{aretenir}[Summary Table: Binomial vs Geometric vs Poisson]
\begin{center}
\begin{tabularx}{\linewidth}{|l|X|X|X|}
\hline
\rowcolor{popblueL}
\textbf{Feature} & \textbf{Binomial $B(n,p)$} & \textbf{Geometric $Geo(p)$} & \textbf{Poisson $Po(\lambda)$} \\ \hline
\cle{Random experiment} & Fixed number $n$ of independent Bernoulli trials. & Repeated independent Bernoulli trials until the \emph{first} success. & Counting random events occurring in a fixed interval of time/space. \\ \hline
\cle{Parameters} & $n \in \mathbb{N}^*$ (number of trials), $p \in (0,1)$ (success probability). & $p \in (0,1)$ (success probability). & $\lambda > 0$ (average rate per unit interval). \\ \hline
\cle{Support (possible values)} & $x = 0,1,2,\dots,n$ & $x = 1,2,3,\dots$ & $x = 0,1,2,\dots$ \\ \hline
\cle{Probability mass function} & $P(X=x) = \binom{n}{x}p^x(1-p)^{n-x}$ & $P(X=x) = (1-p)^{x-1}p$ & $P(X=x) = \dfrac{e^{-\lambda}\lambda^x}{x!}$ \\ \hline
\cle{Expectation} & $E(X) = np$ & $E(X) = \dfrac{1}{p}$ & $E(X) = \lambda$ \\ \hline
\cle{Variance} & $Var(X) = np(1-p)$ & $Var(X) = \dfrac{1-p}{p^2}$ & $Var(X) = \lambda$ \\ \hline
\cle{Key condition} & Fixed $n$, constant $p$, independence. & No fixed $n$, constant $p$, independence, memoryless property. & Events occur singly, independently, at constant average rate. \\ \hline
\cle{Shape} & Symmetric if $p=0.5$, skewed otherwise; bounded. & Strictly decreasing; unbounded right tail. & Skewed for small $\lambda$, approx. symmetric for large $\lambda$; unbounded. \\ \hline
\end{tabularx}
\end{center}
\end{aretenir}

\courssub{Decision Flowchart: Which Distribution to Use?}

When faced with a worded problem, follow the logical steps in the flowchart below. It guides you from the description of the random experiment to the appropriate distribution.

\begin{popfigure}
\begin{tikzpicture}[
    node distance=1.8cm and 2.5cm,
    startstop/.style={rectangle, rounded corners, minimum width=3cm, minimum height=1cm, text centered, draw=popblue, fill=popblueL, thick},
    decision/.style={diamond, aspect=2, minimum width=3.5cm, minimum height=1.2cm, text centered, draw=poporange, fill=poporange!20, thick, text width=3.5cm, align=center},
    process/.style={rectangle, minimum width=3.5cm, minimum height=1cm, text centered, draw=popgreen, fill=popgreenL, thick, text width=3.5cm, align=center},
    arrow/.style={thick,->,>=stealth, popdark}
]
% Nodes
\node (start) [startstop] {Start: Read the problem carefully.};
\node (q1) [decision, below of=start] {Is there a \textbf{fixed number of trials} $n$?};
\node (bin) [process, left of=q1, xshift=-3.5cm, align=center]{\textbf{Binomial $B(n,p)$}\\ Check: constant $p$, independence.};
\node (q2) [decision, right of=q1, xshift=3.5cm] {Are we counting \textbf{events in a fixed interval} (time, area, length)?};
\node (geo) [process, below of=q1, yshift=-1.5cm, align=center]{\textbf{Geometric $Geo(p)$}\\ (waiting for first success)\\ Check: constant $p$, independence.};
\node (poi) [process, below of=q2, yshift=-1.5cm, align=center]{\textbf{Poisson $Po(\lambda)$}\\ Check: events independent, constant average rate.};
\node (end) [startstop, below of=geo, yshift=-1.5cm] {Model chosen. Write down parameters, formula, compute and \textbf{interpret}.};

% Arrows
\draw [arrow] (start) -- (q1);
\draw [arrow] (q1) -- node[above, pos=0.5] {Yes} (bin);
\draw [arrow] (q1) -- node[above, pos=0.5] {No} (q2);
\draw [arrow] (q2) -- node[above, pos=0.5] {Yes} (poi);
\draw [arrow] (q2) -- node[above, pos=0.5] {No} (geo);
\draw [arrow] (bin) |- (end);
\draw [arrow] (poi) |- (end);
\draw [arrow] (geo) -- (end);
\end{tikzpicture}
\legende{Flowchart for choosing between Binomial, Geometric and Poisson distributions.}
\end{popfigure}

\courssub{Comparing Binomial and Poisson Distributions}

The binomial and Poisson distributions are closely linked. The Poisson distribution arises as a limiting case of the binomial when $n \to \infty$, $p \to 0$ while $np = \lambda$ remains constant. This is the basis of the \cle{Poisson approximation to the binomial}, which we studied in the previous section.

\begin{propriete}[When to Use Poisson Approximation to Binomial]
The approximation $B(n,p) \approx Po(\lambda)$ with $\lambda = np$ is appropriate when:
\begin{itemize}
    \item $n$ is large (typically $n \ge 50$),
    \item $p$ is small (typically $p \le 0.1$),
    \item $np = \lambda$ is moderate (typically $\lambda \le 5$ or $\lambda \le 10$).
\end{itemize}
The approximation works because the binomial's fixed upper bound $n$ becomes irrelevant when $p$ is tiny; the Poisson's infinite tail matches the binomial's effective range.
\end{propriete}

\begin{attention}[Common Mistake: Applying Poisson When There Is a Fixed Maximum]
A frequent error is to model a situation with a \emph{fixed maximum number of occurrences} using a Poisson distribution. For example, the number of defective items in a batch of $100$ cannot exceed $100$. If $n=100$ and $p=0.02$, then $np=2$. Although the Poisson approximation $Po(2)$ is numerically excellent, the \emph{correct} model is $B(100,0.02)$ because the experiment has a fixed number of trials. In GCE exams, you must \textbf{state the exact distribution first}, then mention the approximation if appropriate.
\end{attention}

\courssub{Modelling Strategy for Real-Life Problems}

GCE Paper 2 questions often present a real-life scenario and ask you to "state the distribution", "calculate the probability that...", and "interpret your result". Use this four-step method.

\begin{methode}[4-Step Modelling Strategy]
\begin{enumerate}
    \item \textbf{Identify the random variable.} Define $X$ clearly in words (e.g., "$X$ = number of faulty bulbs in a box of 20").
    \item \textbf{Check the conditions.} Match the scenario to the assumptions of a distribution (fixed $n$? waiting for first success? events in an interval?).
    \item \textbf{State the distribution and its parameters.} Write, for example, "$X \sim B(20, 0.05)$" or "$X \sim Po(3.2)$".
    \item \textbf{Compute and interpret.} Perform the calculation (using formula, tables, or approximation) and finish with a sentence in context: "The probability that a box contains at most 1 faulty bulb is 0.736, so about 74\% of boxes meet the quality standard."
\end{enumerate}
\end{methode}

\courssub{Mixed Problems: Combining Binomial and Poisson}

Many exam questions combine two distributions in a single problem. A classic pattern: a Poisson process generates a random number of "items", and each item is then subjected to a binomial trial (e.g., each item is inspected and may be defective).

\begin{exoresolu}[Mixed Binomial--Poisson Problem]
\textbf{Statement:} In a factory, the number of machines that break down in a week follows a Poisson distribution with mean $2.5$. Each broken machine is repaired independently; the probability that a repair is completed within the same week is $0.6$. Find the probability that in a given week, exactly $2$ machines break down and both are repaired within the week.

\tcblower
\textbf{Detailed solution:}
\begin{enumerate}
    \item \textbf{Define variables.}
    Let $N$ = number of breakdowns in a week. $N \sim Po(2.5)$.
    Given $N=n$, let $R$ = number of those $n$ machines repaired within the week. Then $R \mid N=n \sim B(n, 0.6)$.
    \item \textbf{We want $P(N=2 \text{ and } R=2)$.}
    By the multiplication rule:
    \[
    P(N=2, R=2) = P(N=2) \times P(R=2 \mid N=2).
    \]
    \item \textbf{Compute each factor.}
    \[
    P(N=2) = \frac{e^{-2.5}(2.5)^2}{2!} = \frac{6.25}{2}e^{-2.5} = 3.125 e^{-2.5}.
    \]
    Using $e^{-2.5} \approx 0.082085$:
    \[
    P(N=2) \approx 3.125 \times 0.082085 = 0.2565.
    \]
    \[
    P(R=2 \mid N=2) = \binom{2}{2}(0.6)^2(0.4)^0 = 0.36.
    \]
    \item \textbf{Multiply.}
    \[
    P(N=2, R=2) \approx 0.2565 \times 0.36 = 0.09234.
    \]
    \item \textbf{Interpret.}
    The probability that exactly two machines break down and both are repaired within the same week is approximately $0.0923$ (or $9.23\%$).
\end{enumerate}
\end{exoresolu}

\courssub{Interpreting Results in Context}

A numerical answer without interpretation is incomplete in GCE Pure Mathematics with Statistics. Always relate the probability or expectation back to the real situation.

\begin{aretenir}[What "Interpret" Means in GCE Questions]
\begin{itemize}
    \item If you compute $P(X \le k) = 0.85$, say: "There is an $85\%$ chance that at most $k$ events occur."
    \item If you find $E(X) = 4.2$, say: "On average, we expect $4.2$ events per interval (e.g., per week)."
    \item If you compare observed and expected frequencies (goodness-of-fit), say: "The model fits the data well / does not fit the data well because..."
    \item If you use an approximation, say: "Using the Poisson approximation to the binomial, the probability is approximately $0.123$."
\end{itemize}
\end{aretenir}

\courssub{Visual Comparison of Shapes}

To reinforce intuition, the figure below plots the probability mass functions of $B(10, 0.5)$, $Geo(0.3)$ (counting trials until first success) and $Po(3)$ on aligned axes. Notice the bounded symmetric binomial, the strictly decreasing geometric, and the skewed Poisson with a mode at $3$.

\begin{popfigure}
\begin{tikzpicture}
\begin{axis}[
    width=0.95\linewidth,
    height=8cm,
    axis lines=middle,
    xlabel={$x$},
    ylabel={$P(X=x)$},
    xmin=-0.5, xmax=12.5,
    ymin=0, ymax=0.3,
    xtick={0,1,2,3,4,5,6,7,8,9,10,11,12},
    ytick={0,0.05,0.1,0.15,0.2,0.25,0.3},
    tick label style={font=\small},
    label style={font=\small},
    legend style={at={(0.98,0.98)}, anchor=north east, font=\small, fill=white, draw=popline},
    clip=false
]
% Binomial B(10,0.5)
\addplot+[popblue, only marks, mark=*, mark size=2.5, thick] coordinates {
    (0,0.0009766) (1,0.0097656) (2,0.0439453) (3,0.1171875) (4,0.2050781)
    (5,0.2460938) (6,0.2050781) (7,0.1171875) (8,0.0439453) (9,0.0097656) (10,0.0009766)
};
\addlegendentry{$B(10,0.5)$}
% Geometric Geo(0.3) for x=1..12
\addplot+[poporange, only marks, mark=square*, mark size=2.5, thick] coordinates {
    (1,0.3) (2,0.21) (3,0.147) (4,0.1029) (5,0.07203) (6,0.050421) (7,0.0352947) (8,0.0247063) (9,0.0172944) (10,0.0121061) (11,0.0084743) (12,0.005932)
};
\addlegendentry{$Geo(0.3)$}
% Poisson Po(3)
\addplot+[popgreen, only marks, mark=triangle*, mark size=2.5, thick] coordinates {
    (0,0.049787) (1,0.149361) (2,0.224042) (3,0.224042) (4,0.168031) (5,0.100819) (6,0.050409) (7,0.021604) (8,0.008101) (9,0.002700) (10,0.000810) (11,0.000221) (12,0.000055)
};
\addlegendentry{$Po(3)$}
\end{axis}
\end{tikzpicture}
\legende{Comparison of probability mass functions: $B(10,0.5)$ (blue dots), $Geo(0.3)$ (orange squares), $Po(3)$ (green triangles).}
\end{popfigure}

\courssub{Chapter Synthesis and GCE Exam Technique}

This chapter has equipped you with three powerful discrete distributions. Mastery comes from recognising the \emph{structure} of the random experiment, not just memorising formulas. Here is a checklist for exam success.

\begin{methode}[GCE Exam Technique Checklist for Discrete Distributions]
\begin{enumerate}
    \item \textbf{Read twice.} Underline key phrases: "fixed number of trials", "until the first success", "average rate per hour", "randomly scattered".
    \item \textbf{Define $X$ explicitly.} "$X$ = number of ..." --- this earns a mark and focuses your thinking.
    \item \textbf{State the distribution with parameters.} Write $X \sim B(n,p)$, $X \sim Geo(p)$ or $X \sim Po(\lambda)$. If using an approximation, write: "$X \sim B(100,0.02) \approx Po(2)$".
    \item \textbf{Quote the formula.} For binomial: $P(X=x)=\binom{n}{x}p^x(1-p)^{n-x}$. For Poisson: $P(X=x)=e^{-\lambda}\lambda^x/x!$. For geometric: $P(X=x)=(1-p)^{x-1}p$. This shows you know the theory.
    \item \textbf{Show substitution.} Write the numbers into the formula before using a calculator or tables.
    \item \textbf{Use tables correctly.} If using cumulative binomial tables, state clearly: "From tables with $n=20, p=0.3$, $P(X \le 5) = 0.4164$." For Poisson tables, similarly.
    \item \textbf{Round appropriately.} Probabilities to 3 or 4 significant figures (or as requested). Expectations/variances to 1 or 2 decimal places.
    \item \textbf{Conclude in context.} \emph{Always} end with a sentence interpreting the result. This is often a separate mark in the mark scheme.
    \item \textbf{Check variance vs mean.} For Poisson, variance $=$ mean. For binomial, variance $<$ mean. For geometric, variance $>$ mean. This can be a quick sanity check.
\end{enumerate}
\end{methode}

\begin{savaistu}[Historical Note: The Poisson Distribution in Cameroon]
The Poisson distribution is named after Siméon Denis Poisson (1781--1840), a French mathematician who introduced it in 1837. In Cameroon, it is widely used in \emph{queueing theory} for modelling customer arrivals at banks (e.g., Ecobank, Afriland First Bank) and in \emph{telecommunications} for call arrivals at MTN or Orange switching centres. It also models the distribution of rare diseases in epidemiology (e.g., cases of cholera in a health district per month) and the number of defects in manufactured goods (e.g., bottles at the Cameroon Breweries plant in Douala).
\end{savaistu}

\begin{exoresolu}[Real-Life Problem: Choosing and Justifying the Model]
\textbf{Statement:} A call centre receives an average of $12$ calls per hour. The centre has $5$ operators. Each call lasts an average of $5$ minutes. Assume calls arrive randomly and independently.
\begin{enumerate}
    \item State the distribution of the number of calls received in a $30$-minute period.
    \item Find the probability that more than $8$ calls are received in a $30$-minute period.
    \item The manager claims that the probability of all $5$ operators being busy at a given instant is negligible. Comment on this claim using a suitable model.
\end{enumerate}

\tcblower
\textbf{Detailed solution:}
\begin{enumerate}
    \item \textbf{Identify:} Calls arrive randomly, independently, at a constant average rate. This is a Poisson process.
    Rate per hour $\lambda = 12$. For $30$ minutes ($0.5$ hour), $\lambda_{30} = 12 \times 0.5 = 6$.
    Let $X$ = number of calls in $30$ minutes. Then $X \sim Po(6)$.
    \item $P(X > 8) = 1 - P(X \le 8)$.
    Using Poisson cumulative tables with $\lambda = 6$:
    $P(X \le 8) = 0.8472$ (from tables).
    So $P(X > 8) = 1 - 0.8472 = 0.1528$.
    \textbf{Interpretation:} There is about a $15.3\%$ chance that more than $8$ calls arrive in a half-hour period.
    \item The number of operators busy at an instant depends on the number of calls that have arrived in the last $5$ minutes and are still ongoing. We approximate the number of simultaneous calls by the number of arrivals in a $5$-minute interval.
    Average calls in $5$ minutes $= 12 \times (5/60) = 1$.
    Let $Y$ = number of calls in a $5$-minute interval. Then $Y \sim Po(1)$.
    $P(\text{all 5 busy}) \approx P(Y \ge 5) = 1 - P(Y \le 4)$.
    From Poisson tables with $\lambda = 1$: $P(Y \le 4) = 0.9963$.
    Thus $P(Y \ge 5) = 1 - 0.9963 = 0.0037$.
    This is very small ($0.37\%$), so the manager's claim is \textbf{reasonable} under the Poisson assumption. However, real call centres often have bursty traffic (not perfectly Poisson), so the actual probability may be higher.
\end{enumerate}
\end{exoresolu}

\begin{exercice}[Practice: Mixed and Comparison Problems]
\begin{enumerate}
    \item A factory produces bolts. The probability a bolt is defective is $0.02$. Bolts are packed in boxes of $100$.
    \begin{enumerate}
        \item State the exact distribution of the number of defective bolts in a box.
        \item Use a Poisson approximation to find the probability that a box contains at most $3$ defective bolts.
        \item A quality inspector checks $10$ boxes. Find the probability that exactly $2$ boxes contain more than $3$ defectives. (Hint: two-stage binomial--Poisson).
    \end{enumerate}
    \item In a large field, daisies are randomly distributed with an average density of $4$ per square metre. A botanist places a circular hoop of radius $0.5$ m at a random location.
    \begin{enumerate}
        \item State the distribution of the number of daisies inside the hoop.
        \item Find the probability that the hoop contains exactly $2$ daisies.
        \item The botanist repeats this in $20$ independent locations. Find the probability that at least $15$ hoops contain exactly $2$ daisies.
    \end{enumerate}
    \item A student takes a multiple-choice test with $20$ questions, each with $4$ options. She guesses all answers.
    \begin{enumerate}
        \item State the distribution of the number of correct answers.
        \item Find the probability she passes (score $\ge 10$).
        \item If she repeats the test until she passes, what is the distribution of the number of attempts needed? Find the expected number of attempts.
    \end{enumerate}
\end{enumerate}
\end{exercice}

\begin{corrige}[Exercise 1]
\begin{enumerate}
    \item[a)] $X \sim B(100, 0.02)$.
    \item[b)] $\lambda = np = 2$. $X \approx Po(2)$.
    $P(X \le 3) = e^{-2}\left(1 + 2 + \frac{2^2}{2!} + \frac{2^3}{3!}\right) = e^{-2}\left(1 + 2 + 2 + \frac{8}{6}\right) = e^{-2}\left(5 + 1.3333\right) = 6.3333 \times 0.13534 = 0.8571$.
    (Tables give $0.8571$.)
    \item[c)] Let $p = P(\text{box has } >3 \text{ defectives}) = 1 - 0.8571 = 0.1429$.
    Let $Y$ = number of such boxes out of $10$. $Y \sim B(10, 0.1429)$.
    $P(Y=2) = \binom{10}{2}(0.1429)^2(0.8571)^8 = 45 \times 0.02042 \times 0.291 \approx 0.268$.
\end{enumerate}
\end{corrige}

\begin{corrige}[Exercise 2]
\begin{enumerate}
    \item[a)] Area of hoop $= \pi r^2 = \pi \times 0.25 \approx 0.7854$ m$^2$.
    Mean daisies $\lambda = 4 \times 0.7854 = 3.1416$.
    $X \sim Po(3.1416)$.
    \item[b)] $P(X=2) = e^{-3.1416} \frac{3.1416^2}{2!} \approx 0.216$.
    \item[c)] $p = 0.216$. $Y \sim B(20, 0.216)$.
    $P(Y \ge 15) = 1 - P(Y \le 14)$. This is extremely small (use tables or normal approximation). $P(Y \ge 15) \approx 1.2 \times 10^{-7}$.
\end{enumerate}
\end{corrige}

\begin{corrige}[Exercise 3]
\begin{enumerate}
    \item[a)] $X \sim B(20, 0.25)$.
    \item[b)] $P(X \ge 10) = 1 - P(X \le 9)$. Using binomial tables with $n=20, p=0.25$: $P(X \le 9) = 0.9861$, so $P(X \ge 10) = 0.0139$.
    \item[c)] Let $N$ = number of attempts until first pass. $N \sim Geo(p)$ with $p = 0.0139$.
    $E(N) = 1/p = 1/0.0139 \approx 72$ attempts.
\end{enumerate}
\end{corrige}

\begin{attention}[Final Warning: "State the Distribution" Means Write the Symbol]
In GCE exams, "State the distribution of $X$" requires the answer "$X \sim B(n,p)$" or "$X \sim Po(\lambda)$" with the correct numerical parameters. Writing only "binomial" or "Poisson" without parameters loses the mark. Always include the parameters.
\end{attention}

With this section, we complete the chapter on Special Discrete Probability Distributions. You now possess the toolkit to model a wide variety of random phenomena --- from quality control in a Douala factory to call arrivals at a Yaoundé call centre, from the number of attempts to pass a driving test in Buea to the distribution of malaria cases in a health district. The key is not just computation, but \emph{modelling}: translating a real situation into a mathematical framework, solving it, and translating the answer back. That is the essence of statistics at Advanced Level. Good luck in your revision and in the GCE examination!

\newpage

\coursec{Integration activity}

\coursec{Comparing Distributions and Real-Life Problem Solving}

\courssub{Integration Activity: The Biyem-Assi Health Centre Case Study}

\begin{experience}[Context and Objectives]
This integration activity consolidates your understanding of the \cle{binomial} and \cle{Poisson} distributions by applying them to a realistic health-centre scenario in Cameroon. You will:
\begin{itemize}
\item Identify the appropriate probability model for a given random phenomenon (binomial for a fixed number of independent trials with constant success probability; Poisson for events occurring randomly in time/space at a known average rate).
\item Compute and interpret expectation, variance, and probabilities (both exact and using cumulative tables).
\item Apply the \cle{Poisson approximation to the binomial distribution} and critically assess its validity using the standard rule of thumb.
\item Use the \cle{scaling property} of the Poisson process to change the time interval.
\item Determine the \cle{mode} of a Poisson distribution via its recurrence formula.
\item Communicate findings clearly in a concise report suitable for a non-specialist (the head nurse).
\end{itemize}
\end{experience}

\begin{experience}[Situation: Biyem-Assi District Health Centre]
The \emph{Biyem-Assi District Health Centre} in the Centre Region of Cameroon serves a semi-urban population. The centre's laboratory and reception desk generate two distinct streams of data that the head nurse wishes to understand better for staffing and resource planning.

\textbf{Part A -- Malaria Rapid Diagnostic Tests (RDTs).}  
Historical records show that, on average, \textbf{6\%} of patients tested for malaria are positive. Each morning, the laboratory tests exactly \textbf{25 patients} (a fixed number). The result for each patient (positive / negative) is assumed independent of the others, and the probability of a positive result remains constant at $0.06$ throughout the morning session.

\textbf{Part B -- Incoming Telephone Calls at the Reception Desk.}  
The reception desk receives telephone calls randomly throughout the day. Over a long observation period, the average rate is \textbf{4 calls per 15-minute interval}. Calls are independent of each other, and the rate is constant during working hours.

The head nurse has asked your group to produce a one-page statistical brief answering the following operational questions:
\begin{enumerate}
\item For the morning malaria testing session: what is the expected number of positive cases? What is the variability? What is the probability that at most 2 patients test positive? (Use the cumulative binomial probability table.)
\item The laboratory technician suggests using a Poisson approximation to simplify calculations. Compute the Poisson approximation for the probability in (1) and comment on whether the approximation is justified.
\item For the reception desk: what is the probability that no call arrives during a given 15-minute period?
\item What is the probability that more than 6 calls arrive in a 30-minute period? (Hint: scale the rate.)
\item What is the most likely number of calls per hour (i.e., the mode of the hourly distribution)? Use the recurrence formula for the Poisson distribution to justify your answer.
\end{enumerate}
\end{experience}

\begin{experience}[Tasks for Group Work]
Work in groups of 3--4. Each group will submit a single report (maximum one A4 page) containing:
\begin{enumerate}
\item \textbf{Model Justification:} For each part (A and B), state the name of the distribution used, list the assumptions that make it appropriate, and define the parameter(s) clearly.
\item \textbf{Calculations for Part A:}
      \begin{enumerate}
      \item Write the probability mass function of the binomial distribution for the number $X$ of positive tests.
      \item Compute $E(X)$ and $Var(X)$.
      \item Using the cumulative binomial table (provided in the examination booklet), find $P(X \le 2)$.
      \item Compute the Poisson approximation $\tilde{P}(X \le 2)$ with $\lambda = np$. Show the steps.
      \item Compare the two probabilities and state whether the approximation is acceptable (use the rule of thumb: $n \ge 50$, $p \le 0.1$, $np \le 5$).
      \end{enumerate}
\item \textbf{Calculations for Part B:}
      \begin{enumerate}
      \item Write the probability mass function of the Poisson distribution for the number $Y$ of calls in a 15-minute interval.
      \item Find $P(Y = 0)$.
      \item For a 30-minute interval, define the new random variable and its parameter. Compute $P(\text{more than 6 calls})$.
      \item For a 60-minute (1-hour) interval, define the random variable $Z$ and its parameter $\lambda_Z$. Use the recurrence formula $P(Z=k) = \frac{\lambda_Z}{k} P(Z=k-1)$ to determine the mode (the value of $k$ that maximises $P(Z=k)$).
      \end{enumerate}
\item \textbf{Interpretation and Recommendations:} In plain language, explain what the probabilities mean for the head nurse. Suggest one practical action for the laboratory and one for the reception desk based on your findings.
\end{enumerate}
\end{experience}

% Figure: Comparison of Binomial and Poisson PMFs for Part A
\begin{popfigure}
\begin{tikzpicture}
\begin{axis}[
    width=\linewidth,
    height=7cm,
    axis lines=middle,
    xlabel={$k$ (number of positive tests / calls)},
    ylabel={Probability},
    xtick={0,1,2,3,4,5,6,7,8},
    ytick={0,0.1,0.2,0.3,0.4},
    ymin=0, ymax=0.35,
    xmin=-0.5, xmax=8.5,
    legend style={at={(0.98,0.98)}, anchor=north east, font=\small},
    mark size=2.5,
]
% Binomial PMF: n=25, p=0.06
\addplot+[popblue, only marks, mark=*] coordinates {
    (0,0.2129)
    (1,0.3389)
    (2,0.2588)
    (3,0.1310)
    (4,0.0467)
    (5,0.0124)
    (6,0.0026)
};
\addlegendentry{Binomial $B(25, 0.06)$}

% Poisson PMF: lambda=1.5
\addplot+[poporange, only marks, mark=square*] coordinates {
    (0,0.2231)
    (1,0.3347)
    (2,0.2510)
    (3,0.1255)
    (4,0.0471)
    (5,0.0141)
    (6,0.0035)
};
\addlegendentry{Poisson $Po(1.5)$}
\end{axis}
\end{tikzpicture}
\legende{Comparison of the exact Binomial $B(25,0.06)$ and the Poisson $Po(1.5)$ approximation for the number of positive malaria tests. The approximation is visually excellent.}
\end{popfigure}

\begin{exoresolu}[Model Solution and Examiner's Commentary]

\textbf{1. Model Justification}

\begin{itemize}
\item \textbf{Part A (Malaria Tests):} The number of positive tests $X$ follows a \textbf{binomial distribution} $B(n=25, p=0.06)$.  
      \textbf{Assumptions verified:}
      \begin{itemize}
      \item Fixed number of trials: $n = 25$ patients tested each morning.
      \item Two outcomes per trial: Positive (success) / Negative (failure).
      \item Constant success probability: $p = 0.06$ (historical average).
      \item Independence: The result for one patient does not affect another.
      \end{itemize}
\item \textbf{Part B (Phone Calls):} The number of calls in a given time interval follows a \textbf{Poisson distribution}. For a 15-minute interval, $Y \sim Po(\lambda = 4)$.  
      \textbf{Assumptions verified:}
      \begin{itemize}
      \item Events occur randomly and independently in time.
      \item Constant average rate: $4$ calls per 15 minutes.
      \item Two events cannot occur at exactly the same instant (negligible probability of simultaneous calls).
      \end{itemize}
\end{itemize}

\textbf{2. Calculations for Part A}

\begin{enumerate}
\item \textbf{Probability Mass Function (PMF):}
      \[
      P(X = k) = \binom{25}{k} (0.06)^k (0.94)^{25-k}, \quad k = 0,1,\dots,25.
      \]

\item \textbf{Expectation and Variance:}
      \begin{align*}
      E(X) &= np = 25 \times 0.06 = \mathbf{1.5}, \\
      Var(X) &= np(1-p) = 25 \times 0.06 \times 0.94 = \mathbf{1.41}.
      \end{align*}
      The standard deviation is $\sigma_X = \sqrt{1.41} \approx 1.19$.

\item \textbf{Cumulative Probability $P(X \le 2)$ using Tables:}  
      From the cumulative binomial table for $n=25$, $p=0.06$ (or $p=0.05$ and $p=0.10$ with interpolation if necessary; here we assume the exact table for $p=0.06$ is available):
      \[
      P(X \le 2) = \mathbf{0.8106}.
      \]
      \textit{(Verification by direct calculation: $P(X=0)=0.94^{25}\approx 0.2129$, $P(X=1)=25\times0.06\times0.94^{24}\approx 0.3389$, $P(X=2)=\binom{25}{2}0.06^2 0.94^{23}\approx 0.2588$, $\text{Sum} \approx 0.8106$.)}

\item \textbf{Poisson Approximation:} Take $\lambda = np = 1.5$. Then $X \approx \tilde{X} \sim Po(1.5)$.
      \[
      \tilde{P}(X \le 2) = e^{-1.5}\left(1 + 1.5 + \frac{1.5^2}{2!}\right)
      = e^{-1.5}\left(1 + 1.5 + 1.125\right)
      = e^{-1.5} \times 3.625.
      \]
      Using $e^{-1.5} \approx 0.223130$, we get
      \[
      \tilde{P}(X \le 2) \approx 0.223130 \times 3.625 = \mathbf{0.8088}.
      \]

\item \textbf{Comparison and Validity Assessment:}
      \begin{center}
      \begin{tabular}{|l|c|c|}\hline
      \rowcolor{popblueL}
      \textbf{Measure} & \textbf{Exact Binomial} & \textbf{Poisson Approx.} \\ \hline
      $P(X \le 2)$ & $0.8106$ & $0.8088$ \\ \hline
      Absolute Error & \multicolumn{2}{c|}{$0.0018$} \\ \hline
      Relative Error & \multicolumn{2}{c|}{$0.22\%$} \\ \hline
      \end{tabular}
      \end{center}
      \textbf{Rule of Thumb Check:} $n = 25$ (\textbf{not} $\ge 50$), $p = 0.06 \le 0.1$, $np = 1.5 \le 5$.  
      \textbf{Conclusion:} Although $n < 50$, the approximation is \textbf{excellent} (error $< 0.3\%$) because $p$ is small and $np$ is moderate. It is acceptable for practical operational planning.
\end{enumerate}

\textbf{3. Calculations for Part B}

\begin{enumerate}
\item \textbf{PMF for 15-minute interval:} $Y \sim Po(4)$.
      \[
      P(Y = k) = \frac{e^{-4} 4^k}{k!}, \quad k = 0,1,2,\dots
      \]

\item \textbf{Probability of No Call in 15 Minutes:}
      \[
      P(Y = 0) = e^{-4} \approx \mathbf{0.0183} \quad (1.83\%).
      \]

\item \textbf{30-Minute Interval -- Scaling the Rate:}  
      The Poisson rate scales linearly with the length of the interval. For 30 minutes (2 units of 15 minutes), $\lambda_{30} = 4 \times 2 = 8$.  
      Let $W \sim Po(8)$ be the number of calls in 30 minutes.
      \[
      P(W > 6) = 1 - P(W \le 6).
      \]
      Using the cumulative Poisson table for $\lambda = 8$:
      $P(W \le 6) = 0.3134$. Hence
      \[
      P(W > 6) = 1 - 0.3134 = \mathbf{0.6866} \quad (68.66\%).
      \]

\item \textbf{1-Hour (60-Minute) Interval and Mode:}  
      For 60 minutes (4 units of 15 minutes), $\lambda_{60} = 4 \times 4 = 16$. Let $Z \sim Po(16)$.  
      The \cle{recurrence formula} for Poisson probabilities is:
      \[
      P(Z = k) = \frac{\lambda}{k} P(Z = k-1), \quad k \ge 1.
      \]
      Probabilities increase while $\frac{\lambda}{k} > 1 \iff k < \lambda$, and decrease when $k > \lambda$.  
      Since $\lambda = 16$ is an \textbf{integer}, the maximum probability occurs at two adjacent values:
      \[
      k = \lambda - 1 = 15 \quad \text{and} \quad k = \lambda = 16.
      \]
      Both give the same probability. Thus the distribution has \textbf{two modes: 15 and 16 calls per hour}.
\end{enumerate}

\textbf{4. Interpretation and Recommendations for the Head Nurse}

\begin{itemize}
\item \textbf{Laboratory (Malaria Testing):} On average, $1.5$ patients test positive each morning. However, there is a $1 - 0.8106 = 0.1894 \approx \mathbf{19\%}$ chance that \textbf{more than 2 patients} are positive. The Poisson approximation ($19.12\%$) confirms this risk.
      \begin{attention}[Operational Implication]
      Preparing only 2 treatment kits each morning would leave the centre unprepared roughly 1 morning in 5.
      \end{attention}
      \textbf{Recommendation:} Prepare \textbf{at least 3 treatment kits} every morning to cover the vast majority ($>81\%$) of sessions. Keep a small buffer stock for the remaining $19\%$ of days.

\item \textbf{Reception Desk (Phone Calls):} The probability of a completely quiet 15-minute period is very low ($1.8\%$). In any 30-minute window, there is a $\mathbf{69\%}$ chance of receiving \textbf{more than 6 calls}. The hourly mode of $15$--$16$ calls indicates a consistently high and predictable workload.
      \textbf{Recommendation:} Schedule a \textbf{second receptionist during peak hours} (e.g., 8\,am--12\,pm) to manage the high call volume and keep patient waiting times acceptable. Consider an automated call-back system for non-urgent queries.
\end{itemize}
\end{exoresolu}

\begin{aretenir}[Key Examination Takeaways for PMS12]
\begin{itemize}
\item \textbf{Binomial $B(n,p)$:} Fixed $n$, independent trials, constant $p$. $E=np$, $Var=np(1-p)$. Use tables for cumulative probabilities.
\item \textbf{Poisson $Po(\lambda)$:} Random events in time/space, constant rate, independence. $E=Var=\lambda$. PMF: $P(X=k)=e^{-\lambda}\lambda^k/k!$.
\item \textbf{Poisson Approximation to Binomial:} Valid when $n$ is large ($\ge 50$), $p$ is small ($\le 0.1$), $np$ is moderate ($\le 5$). Use $\lambda = np$. Always check conditions and comment on accuracy.
\item \textbf{Scaling Poisson Rate:} If rate is $\lambda$ per unit time, then for $t$ units the rate is $t\lambda$. The random variable for the new interval is $Po(t\lambda)$.
\item \textbf{Mode of $Po(\lambda)$:} Use recurrence $P(k)=\frac{\lambda}{k}P(k-1)$. If $\lambda$ not integer, mode $= \lfloor \lambda \rfloor$. If $\lambda$ integer, modes are $\lambda$ and $\lambda-1$.
\item \textbf{Contextual Justification:} Always explicitly link the assumptions of the distribution to the wording of the problem (fixed $n$ vs. rate in time/space).
\item \textbf{Communication:} Translate probabilities into plain language (e.g., ``19\% chance of shortage'') and propose concrete actions.
\end{itemize}
\end{aretenir}

\begin{methode}[Choosing Between Binomial and Poisson Models]
\begin{enumerate}
\item \textbf{Identify the random variable:} Is it a count of successes in a fixed number of trials? $\to$ Binomial. Is it a count of events occurring randomly in a continuous interval (time, area, length)? $\to$ Poisson.
\item \textbf{Check for a fixed $n$:} If the problem states ``out of 25 patients'' or ``in a sample of 100'', $n$ is fixed $\to$ Binomial. If it states ``average 4 calls per 15 minutes'' with no upper limit, $n$ is not fixed $\to$ Poisson.
\item \textbf{Check for constant $p$ vs. constant rate $\lambda$:} Binomial needs constant success probability per trial. Poisson needs constant average rate per unit interval.
\item \textbf{Check independence:} Required for both.
\item \textbf{Approximation?} If Binomial conditions hold but $n$ is large and $p$ small, Poisson approximation may simplify calculations. State the rule of thumb and verify.
\end{enumerate}
\end{methode}

\begin{attention}[Common Pitfalls in Examinations]
\begin{itemize}
\item \textbf{Confusing parameters:} Writing $Po(np)$ for the approximation but forgetting to define $\lambda = np$ explicitly.
\item \textbf{Forgetting to scale $\lambda$:} Using $\lambda=4$ for a 30-minute or 1-hour interval without multiplying by the number of 15-minute units.
\item \textbf{Mode of Poisson:} Stating the mode is $\lambda$ (e.g., 16) without checking if $\lambda$ is integer. If $\lambda$ is integer, you \textbf{must} give both $\lambda$ and $\lambda-1$.
\item \textbf{Table usage:} Writing $P(X < 2)$ instead of $P(X \le 2)$ when reading cumulative tables. Remember tables give $P(X \le x)$.
\item \textbf{Continuity correction:} Do \textbf{not} apply continuity correction for Poisson approximation to Binomial (that is for Normal approximation).
\item \textbf{Units in final answer:} Always include units (e.g., ``1.5 patients'', ``16 calls per hour'').
\end{itemize}
\end{attention}

\begin{savaistu}[Historical Note: The Poisson Distribution in Cameroon]
The Poisson distribution is named after Siméon Denis Poisson (1781--1840), a French mathematician. In Cameroon, it is widely used in \emph{queueing theory} to model customer arrivals at banks (e.g., BICEC, Afriland First Bank), call centres (MTN, Orange), and health facilities like the one in this activity. The \emph{Erlang formula}, derived from the Poisson process, is the mathematical backbone of modern telecommunications network dimensioning used by Camtel and mobile operators to determine the number of lines needed to keep call blocking probability low.
\end{savaistu}

\newpage

\coursec{Methods and worked examples}

\courssub{Methods and Worked Examples}

\begin{methode}[Recognising and Computing Binomial Probabilities]
\begin{enumerate}
\item \textbf{Identify the binomial setting:} a fixed number $n$ of independent trials, each with two outcomes (Success/Failure), and constant success probability $p$.
\item \textbf{Define the random variable:} $X \sim B(n,p)$ where $X$ = number of successes in the $n$ trials.
\item \textbf{Write the probability mass function:} $\mathbb{P}(X = k) = \binom{n}{k} p^k (1-p)^{n-k}$ for $k = 0,1,\dots,n$.
\item \textbf{Compute:} substitute $n$, $p$ and the required $k$; evaluate $\binom{n}{k}$ and the powers.
\item \textbf{Check:} probabilities must sum to $1$; each $\mathbb{P}(X=k) \in [0,1]$.
\end{enumerate}
\end{methode}

\begin{exoresolu}[Binomial Probability: Defective Items in a Batch]
A factory produces light bulbs. The probability that a bulb is defective is $0.03$. A quality control officer selects a random sample of $20$ bulbs.
\begin{enumerate}
\item Find the probability that exactly $2$ bulbs are defective.
\item Find the probability that at most $1$ bulb is defective.
\item Find the probability that at least $3$ bulbs are defective.
\end{enumerate}
\tcblower
\textbf{Solution.}

Let $X$ be the number of defective bulbs in the sample of $20$. The trials are independent, $n=20$ is fixed and $p=0.03$ is constant, so $X \sim B(20, 0.03)$.
\begin{enumerate}
\item 
\[
\mathbb{P}(X = 2) = \binom{20}{2} (0.03)^2 (0.97)^{18}
= 190 \times 0.0009 \times (0.97)^{18}.
\]
Since $(0.97)^{18} \approx 0.5780$,
\[
\mathbb{P}(X = 2) \approx 190 \times 0.0009 \times 0.5780 \approx 0.0988.
\]
\item $\mathbb{P}(X \le 1) = \mathbb{P}(X = 0) + \mathbb{P}(X = 1)$.
\[
\mathbb{P}(X = 0) = (0.97)^{20} \approx 0.5438,
\]
\[
\mathbb{P}(X = 1) = \binom{20}{1}(0.03)(0.97)^{19} = 20 \times 0.03 \times 0.5616 \approx 0.3370.
\]
Hence $\mathbb{P}(X \le 1) \approx 0.5438 + 0.3370 = 0.8808$.
\item Using the complement,
\[
\mathbb{P}(X \ge 3) = 1 - \mathbb{P}(X \le 2) = 1 - \big[\mathbb{P}(X \le 1) + \mathbb{P}(X = 2)\big].
\]
$\mathbb{P}(X \le 2) \approx 0.8808 + 0.0988 = 0.9796$, so
\[
\mathbb{P}(X \ge 3) \approx 1 - 0.9796 = 0.0204.
\]
\end{enumerate}
\end{exoresolu}

\begin{attention}[``At least'', ``at most'', ``more than'']
For a discrete variable, $\mathbb{P}(X \ge k) = 1 - \mathbb{P}(X \le k-1)$ and $\mathbb{P}(X > k) = 1 - \mathbb{P}(X \le k)$. A very common GCE error is to subtract $\mathbb{P}(X \le k)$ when ``at least $k$'' is required, thus losing the term $\mathbb{P}(X=k)$.
\end{attention}

\begin{methode}[Expectation, Variance and Recurrence Formula for the Binomial Distribution]
\begin{enumerate}
\item \textbf{Expectation:} $\mathbb{E}(X) = np$.
\item \textbf{Variance:} $\mathrm{Var}(X) = np(1-p) = npq$ where $q = 1-p$.
\item \textbf{Standard deviation:} $\sigma = \sqrt{npq}$.
\item \textbf{Recurrence formula:} $\mathbb{P}(X = k+1) = \dfrac{n-k}{k+1} \cdot \dfrac{p}{q} \cdot \mathbb{P}(X = k)$ for $k = 0,1,\dots,n-1$. It is obtained by dividing the expression for $\mathbb{P}(X=k+1)$ by that for $\mathbb{P}(X=k)$ and simplifying the binomial coefficients.
\item \textbf{Mode:} the value(s) of $k$ maximising $\mathbb{P}(X=k)$. If $(n+1)p$ is an integer, there are two modes, $(n+1)p$ and $(n+1)p-1$; otherwise the mode is $\lfloor (n+1)p \rfloor$.
\end{enumerate}
\end{methode}

\begin{exoresolu}[Binomial Expectation, Variance, Recurrence and Mode]
The random variable $X \sim B(12, 0.4)$.
\begin{enumerate}
\item Find $\mathbb{E}(X)$ and $\mathrm{Var}(X)$.
\item Given that $\mathbb{P}(X = 3) = 0.1419$ (to 4 d.p.), use the recurrence formula to find $\mathbb{P}(X = 4)$ and $\mathbb{P}(X = 5)$.
\item Determine the mode of the distribution.
\end{enumerate}
\tcblower
\textbf{Solution.}
\begin{enumerate}
\item $\mathbb{E}(X) = np = 12 \times 0.4 = 4.8$.
\[
\mathrm{Var}(X) = npq = 12 \times 0.4 \times 0.6 = 2.88, \qquad \sigma = \sqrt{2.88} \approx 1.697.
\]
\item Here $\dfrac{p}{q} = \dfrac{0.4}{0.6} = \dfrac{2}{3}$, so the recurrence formula gives
\[
\mathbb{P}(X = k+1) = \frac{12-k}{k+1} \times \frac{2}{3} \times \mathbb{P}(X = k).
\]
For $k=3$:
\[
\mathbb{P}(X = 4) = \frac{9}{4} \times \frac{2}{3} \times 0.1419 = 1.5 \times 0.1419 = 0.2129.
\]
For $k=4$:
\[
\mathbb{P}(X = 5) = \frac{8}{5} \times \frac{2}{3} \times 0.2129 = \frac{16}{15} \times 0.2129 \approx 0.2271.
\]
\item $(n+1)p = 13 \times 0.4 = 5.2$, which is not an integer, so the mode is $\lfloor 5.2 \rfloor = 5$.

Check: $\mathbb{P}(X=6) = \dfrac{7}{6}\times\dfrac{2}{3}\times 0.2271 \approx 0.1766 < \mathbb{P}(X=5)$, and $\mathbb{P}(X=4) = 0.2129 < \mathbb{P}(X=5)$. Since the probabilities increase up to the mode and decrease after it, $k=5$ is indeed the mode.
\end{enumerate}
\end{exoresolu}

\begin{methode}[Using Cumulative Binomial Tables and Complementary Probabilities]
\begin{enumerate}
\item \textbf{Identify the parameters} $n$ and $p$. If $p > 0.5$, use the symmetry: $X \sim B(n,p) \Rightarrow Y = n-X \sim B(n,1-p)$, and tables for $1-p$.
\item \textbf{Read cumulative probabilities:} tables give $\mathbb{P}(X \le x)$.
\item \textbf{Convert the required probability:}
\begin{itemize}
\item $\mathbb{P}(X \le k)$: read directly from the table;
\item $\mathbb{P}(X < k) = \mathbb{P}(X \le k-1)$;
\item $\mathbb{P}(X \ge k) = 1 - \mathbb{P}(X \le k-1)$;
\item $\mathbb{P}(X > k) = 1 - \mathbb{P}(X \le k)$;
\item $\mathbb{P}(a \le X \le b) = \mathbb{P}(X \le b) - \mathbb{P}(X \le a-1)$.
\end{itemize}
\item \textbf{If $p > 0.5$:} with $Y = n - X \sim B(n, 1-p)$, use $\mathbb{P}(X = k) = \mathbb{P}(Y = n-k)$ and $\mathbb{P}(X \ge k) = \mathbb{P}(Y \le n-k)$.
\end{enumerate}
\end{methode}

\begin{exoresolu}[Cumulative Binomial Tables with $p > 0.5$]
The probability that a student passes a test is $0.75$. In a group of $15$ students, find the probability that:
\begin{enumerate}
\item exactly $10$ pass,
\item at least $12$ pass,
\item between $8$ and $11$ inclusive pass.
\end{enumerate}
Use cumulative tables for $B(15, 0.25)$, which give $\mathbb{P}(Y \le 3) = 0.4613$, $\mathbb{P}(Y \le 4) = 0.6865$, $\mathbb{P}(Y \le 5) = 0.8516$ and $\mathbb{P}(Y \le 7) = 0.9827$.
\tcblower
\textbf{Solution.}

Let $X$ = number of passes, so $X \sim B(15, 0.75)$. Since $p = 0.75 > 0.5$, set $Y = 15 - X \sim B(15, 0.25)$ (the number of failures).
\begin{enumerate}
\item $\mathbb{P}(X = 10) = \mathbb{P}(Y = 5) = \mathbb{P}(Y \le 5) - \mathbb{P}(Y \le 4) = 0.8516 - 0.6865 = 0.1651$.
\item $\mathbb{P}(X \ge 12) = \mathbb{P}(Y \le 3) = 0.4613$.
\item $\mathbb{P}(8 \le X \le 11) = \mathbb{P}(X \le 11) - \mathbb{P}(X \le 7)$.
\[
\mathbb{P}(X \le 11) = \mathbb{P}(Y \ge 4) = 1 - \mathbb{P}(Y \le 3) = 1 - 0.4613 = 0.5387,
\]
\[
\mathbb{P}(X \le 7) = \mathbb{P}(Y \ge 8) = 1 - \mathbb{P}(Y \le 7) = 1 - 0.9827 = 0.0173.
\]
Hence $\mathbb{P}(8 \le X \le 11) = 0.5387 - 0.0173 = 0.5214$.
\end{enumerate}
\end{exoresolu}

\begin{methode}[Fitting a Binomial Distribution to Observed Data]
\begin{enumerate}
\item \textbf{Calculate the sample mean:} $\bar{x} = \dfrac{\sum f_i x_i}{\sum f_i}$.
\item \textbf{Estimate $p$:} for $B(n,p)$, $\mathbb{E}(X) = np$, so set $\hat{p} = \dfrac{\bar{x}}{n}$ (the value of $n$ is known from the context).
\item \textbf{Compute expected frequencies:} $E_i = N \times \mathbb{P}(X = x_i)$, where $N = \sum f_i$ and $\mathbb{P}(X = x_i) = \binom{n}{x_i} \hat{p}^{x_i} (1-\hat{p})^{n-x_i}$.
\item \textbf{Pool adjacent classes} if any expected frequency is less than $5$ (a requirement for the $\chi^2$ goodness-of-fit test).
\item \textbf{Compare observed and expected frequencies} to judge the fit.
\end{enumerate}
\end{methode}

\begin{exoresolu}[Fitting a Binomial Distribution]
The following data show the number of boys in $200$ families, each with $5$ children.
\begin{center}
\begin{tabular}{|c|cccccc|}\hline
Number of boys ($x$) & 0 & 1 & 2 & 3 & 4 & 5 \\ \hline
Number of families ($f$) & 8 & 32 & 70 & 60 & 25 & 5 \\ \hline
\end{tabular}
\end{center}
Fit a binomial distribution and calculate the expected frequencies.
\tcblower
\textbf{Solution.}

Here $n = 5$ is fixed and $N = 200$.
\[
\bar{x} = \frac{0(8) + 1(32) + 2(70) + 3(60) + 4(25) + 5(5)}{200}
= \frac{0 + 32 + 140 + 180 + 100 + 25}{200} = \frac{477}{200} = 2.385.
\]
Estimate $p$: $\hat{p} = \dfrac{\bar{x}}{n} = \dfrac{2.385}{5} = 0.477$, so $\hat{q} = 1 - 0.477 = 0.523$.

Compute $\mathbb{P}(X = x) = \binom{5}{x} (0.477)^x (0.523)^{5-x}$, using
$(0.523)^2 = 0.2735$, $(0.523)^3 = 0.1431$, $(0.523)^4 = 0.0748$, $(0.523)^5 = 0.0391$.
\begin{itemize}
\item $x=0$: $(0.523)^5 = 0.0391$; \quad $E_0 = 200 \times 0.0391 = 7.83$.
\item $x=1$: $5(0.477)(0.523)^4 = 5 \times 0.477 \times 0.0748 = 0.1784$; \quad $E_1 = 35.69$.
\item $x=2$: $10(0.477)^2(0.523)^3 = 10 \times 0.2275 \times 0.1431 = 0.3255$; \quad $E_2 = 65.10$.
\item $x=3$: $10(0.477)^3(0.523)^2 = 10 \times 0.1085 \times 0.2735 = 0.2968$; \quad $E_3 = 59.35$.
\item $x=4$: $5(0.477)^4(0.523) = 5 \times 0.0518 \times 0.523 = 0.1354$; \quad $E_4 = 27.08$.
\item $x=5$: $(0.477)^5 = 0.0247$; \quad $E_5 = 200 \times 0.0247 = 4.94$.
\end{itemize}
The expected frequencies sum to $7.83 + 35.69 + 65.10 + 59.35 + 27.08 + 4.94 \approx 200$, as required.
\begin{center}
\begin{tabular}{|c|cccccc|}\hline
$x$ & 0 & 1 & 2 & 3 & 4 & 5 \\ \hline
Observed $f$ & 8 & 32 & 70 & 60 & 25 & 5 \\ \hline
Expected $E$ & 7.8 & 35.7 & 65.1 & 59.4 & 27.1 & 4.9 \\ \hline
\end{tabular}
\end{center}
The agreement is close, so a binomial model is reasonable. Note that $E_5 = 4.94 < 5$: for a $\chi^2$ goodness-of-fit test, the last two classes ($x=4$ and $x=5$) should be pooled, giving observed $30$ and expected $32.0$.
\end{exoresolu}

\begin{methode}[Geometric Distribution: Probabilities, Expectation and Variance]
\begin{enumerate}
\item \textbf{Setting:} independent Bernoulli trials with success probability $p$; $X$ = number of trials \textbf{up to and including} the first success. Support: $x = 1,2,3,\dots$
\item \textbf{PMF:} $\mathbb{P}(X = x) = (1-p)^{x-1}p = q^{x-1}p$.
\item \textbf{Expectation:} $\mathbb{E}(X) = \dfrac{1}{p}$.
\item \textbf{Variance:} $\mathrm{Var}(X) = \dfrac{q}{p^2}$.
\item \textbf{Tail probabilities:} $\mathbb{P}(X > x) = q^x$ (the first $x$ trials all fail); $\mathbb{P}(X \ge x) = q^{x-1}$.
\item \textbf{Memoryless property:} $\mathbb{P}(X > m+n \mid X > m) = \mathbb{P}(X > n)$: past failures do not change future probabilities.
\end{enumerate}
\end{methode}

\begin{exoresolu}[Geometric Distribution: Drilling for Oil]
An oil company drills wells independently; each well has probability $0.2$ of striking oil.
\begin{enumerate}
\item Find the probability that the first strike occurs on the $4$th well drilled.
\item Find the probability that the first strike occurs on or before the $3$rd well.
\item Find the expected number of wells drilled until the first strike, and the variance.
\item Given that the first $5$ wells are dry, find the probability that the first strike occurs on the $8$th well.
\end{enumerate}
\tcblower
\textbf{Solution.}

$X \sim \mathrm{Geo}(0.2)$, with $p = 0.2$ and $q = 0.8$.
\begin{enumerate}
\item $\mathbb{P}(X = 4) = q^3 p = (0.8)^3 \times 0.2 = 0.512 \times 0.2 = 0.1024$.
\item $\mathbb{P}(X \le 3) = 1 - \mathbb{P}(X > 3) = 1 - q^3 = 1 - 0.512 = 0.488$.

(Check: $\mathbb{P}(X=1)+\mathbb{P}(X=2)+\mathbb{P}(X=3) = 0.2 + 0.16 + 0.128 = 0.488$.)
\item $\mathbb{E}(X) = \dfrac{1}{p} = \dfrac{1}{0.2} = 5$ wells; \quad $\mathrm{Var}(X) = \dfrac{q}{p^2} = \dfrac{0.8}{0.04} = 20$.
\item By the memoryless property,
\[
\mathbb{P}(X = 8 \mid X > 5) = \mathbb{P}(X = 3) = q^2 p = (0.8)^2 \times 0.2 = 0.128.
\]
(Directly: $\mathbb{P}(X=8 \mid X>5) = \dfrac{\mathbb{P}(X=8)}{\mathbb{P}(X>5)} = \dfrac{q^7 p}{q^5} = q^2 p$.)
\end{enumerate}
\end{exoresolu}

\begin{methode}[Poisson Distribution: Characteristics, Expectation, Variance, Recurrence and Mode]
\begin{enumerate}
\item \textbf{Setting:} random events occurring independently, singly, and at a constant average rate $\lambda$ per unit interval (of time, length, area, volume, \dots).
\item \textbf{PMF:} $\mathbb{P}(X = x) = \dfrac{e^{-\lambda} \lambda^x}{x!}$, for $x = 0,1,2,\dots$
\item \textbf{Expectation and variance:} $\mathbb{E}(X) = \lambda$ and $\mathrm{Var}(X) = \lambda$ (a hallmark of the Poisson distribution: mean equals variance).
\item \textbf{Recurrence formula:} $\mathbb{P}(X = x+1) = \dfrac{\lambda}{x+1}\,\mathbb{P}(X = x)$, obtained by dividing consecutive terms of the PMF.
\item \textbf{Mode:} if $\lambda$ is an integer, there are two modes, $\lambda$ and $\lambda - 1$; otherwise the mode is $\lfloor \lambda \rfloor$.
\item \textbf{Sum of independent Poisson variables:} if $X \sim \mathrm{Po}(\lambda_1)$ and $Y \sim \mathrm{Po}(\lambda_2)$ are independent, then $X+Y \sim \mathrm{Po}(\lambda_1 + \lambda_2)$. In particular, changing the unit interval simply scales $\lambda$.
\end{enumerate}
\end{methode}

\begin{exoresolu}[Poisson Distribution: Calls at a Call Centre]
A call centre in Douala receives an average of $6$ calls per hour. Assume a Poisson distribution.
\begin{enumerate}
\item Find the probability of exactly $4$ calls in a given hour.
\item Find the probability of at most $2$ calls in a $30$-minute period.
\item Find the probability of more than $10$ calls in a $2$-hour period, given that $\mathbb{P}(Z \le 10) = 0.3472$ for $Z \sim \mathrm{Po}(12)$.
\item Determine the mode of the distribution for a $1$-hour period.
\end{enumerate}
\tcblower
\textbf{Solution.}
\begin{enumerate}
\item For one hour, $X \sim \mathrm{Po}(6)$.
\[
\mathbb{P}(X = 4) = \frac{e^{-6}\, 6^4}{4!} = e^{-6} \times \frac{1296}{24} = 54\,e^{-6} \approx 54 \times 0.002479 = 0.1339.
\]
\item For $30$ minutes the rate is $\lambda = 6 \times 0.5 = 3$, so $Y \sim \mathrm{Po}(3)$.
\[
\mathbb{P}(Y \le 2) = e^{-3}\left(1 + 3 + \frac{3^2}{2!}\right) = 8.5\,e^{-3} \approx 8.5 \times 0.049787 = 0.4232.
\]
\item For $2$ hours, $\lambda = 12$, so $Z \sim \mathrm{Po}(12)$.
\[
\mathbb{P}(Z > 10) = 1 - \mathbb{P}(Z \le 10) = 1 - 0.3472 = 0.6528.
\]
\item $\lambda = 6$ is an integer, so there are two modes: $5$ and $6$.

Check with the recurrence formula: $\mathbb{P}(X=6) = \dfrac{6}{6}\,\mathbb{P}(X=5) = \mathbb{P}(X=5)$, so the two probabilities are equal; and $\mathbb{P}(X=7) = \dfrac{6}{7}\mathbb{P}(X=6) < \mathbb{P}(X=6)$, confirming the maximum.
\end{enumerate}
\end{exoresolu}

\begin{methode}[Poisson Approximation to the Binomial Distribution]
\begin{enumerate}
\item \textbf{Conditions:} $n$ large (typically $n \ge 50$), $p$ small (typically $p \le 0.1$), and $\lambda = np$ moderate (typically $\lambda \le 5$).
\item \textbf{Approximation:} $X \sim B(n,p) \approx Y \sim \mathrm{Po}(\lambda)$ with $\lambda = np$.
\item \textbf{No continuity correction} is needed (discrete approximated by discrete).
\item \textbf{Compute:} $\mathbb{P}(X = k) \approx \dfrac{e^{-\lambda}\lambda^k}{k!}$, and similarly for cumulative probabilities.
\item \textbf{Decide:} if $n$ is small, use the exact binomial; if $p$ is not small but $np > 5$ and $nq > 5$, a normal approximation is preferable.
\end{enumerate}
\end{methode}

\begin{exoresolu}[Poisson Approximation to the Binomial]
A machine produces $2\%$ defective items. In a batch of $200$ items, use a Poisson approximation to find the probability that:
\begin{enumerate}
\item there are exactly $3$ defectives,
\item there are at most $2$ defectives,
\item there are more than $5$ defectives.
\end{enumerate}
Comment on the suitability of the approximation.
\tcblower
\textbf{Solution.}

$n = 200$, $p = 0.02$, so $\lambda = np = 4$. The conditions hold: $n = 200 \ge 50$, $p = 0.02 \le 0.1$, $\lambda = 4 \le 5$. Thus $X \sim B(200, 0.02) \approx Y \sim \mathrm{Po}(4)$.
\begin{enumerate}
\item $\mathbb{P}(X = 3) \approx \mathbb{P}(Y = 3) = \dfrac{e^{-4}\,4^3}{3!} = \dfrac{64}{6}e^{-4} \approx 10.6667 \times 0.018316 = 0.1954$.

(Exact binomial value: $\binom{200}{3}(0.02)^3(0.98)^{197} \approx 0.1964$. Very close.)
\item $\mathbb{P}(X \le 2) \approx \mathbb{P}(Y \le 2) = e^{-4}\left(1 + 4 + \dfrac{4^2}{2!}\right) = 13\,e^{-4} \approx 13 \times 0.018316 = 0.2381$.

(Exact binomial value: $\approx 0.2351$.)
\item $\mathbb{P}(X > 5) \approx 1 - \mathbb{P}(Y \le 5)$. Using the recurrence formula $\mathbb{P}(Y=x+1) = \dfrac{4}{x+1}\mathbb{P}(Y=x)$:
\[
\mathbb{P}(Y=3) = 0.1954, \quad \mathbb{P}(Y=4) = \tfrac{4}{4}(0.1954) = 0.1954, \quad \mathbb{P}(Y=5) = \tfrac{4}{5}(0.1954) = 0.1563.
\]
\[
\mathbb{P}(Y \le 5) \approx 0.2381 + 0.1954 + 0.1954 + 0.1563 = 0.7852,
\]
\[
\mathbb{P}(X > 5) \approx 1 - 0.7852 = 0.2148.
\]
(Exact binomial value: $\approx 0.2131$.)
\end{enumerate}
The approximation is excellent here because $n$ is large, $p$ is small and $\lambda = 4$ is moderate; the errors are of the order of $0.003$ or less.
\end{exoresolu}

\begin{methode}[Comparing Distributions and Solving Real-Life Problems]
\begin{enumerate}
\item \textbf{Identify the random process:}
\begin{itemize}
\item fixed number $n$ of trials, two outcomes, constant $p$ $\rightarrow$ Binomial;
\item number of trials up to the first success $\rightarrow$ Geometric;
\item events occurring randomly, singly and at a constant average rate in a continuous interval $\rightarrow$ Poisson.
\end{itemize}
\item \textbf{Write down the parameters:} $(n,p)$ for the binomial; $p$ for the geometric; $\lambda = \text{rate} \times \text{interval}$ for the Poisson.
\item \textbf{Choose the appropriate distribution} and compute the required probability (exact, cumulative, recurrence or approximation).
\item \textbf{Use approximations when appropriate:} Poisson to binomial ($n$ large, $p$ small); normal to binomial ($np > 5$ and $nq > 5$, with continuity correction); normal to Poisson ($\lambda$ large, with continuity correction).
\item \textbf{Interpret the result in context}, stating clearly what the probability means for the real situation.
\end{enumerate}
\end{methode}

\begin{exoresolu}[Real-Life Problems: Choosing the Right Distribution]
\begin{enumerate}
\item A typist makes on average $1.5$ errors per page. Find the probability that a $3$-page letter contains no errors.
\item A multiple-choice test has $50$ questions, each with $4$ options. A student guesses all the answers. Using a suitable approximation, find the probability that the student gets at least $15$ correct.
\item A machine breaks down on average once every $10$ days. Find the probability that it breaks down at least twice in a $30$-day month.
\item A box contains $100$ fuses, $5$ of which are defective. A sample of $10$ fuses is taken without replacement. Find the probability that exactly $1$ is defective, naming the exact distribution and a suitable approximation.
\end{enumerate}
\tcblower
\textbf{Solution.}
\begin{enumerate}
\item Errors occur randomly and independently at a constant average rate over a continuous ``interval'' (pages): \textbf{Poisson}. For $3$ pages, $\lambda = 1.5 \times 3 = 4.5$, so $X \sim \mathrm{Po}(4.5)$.
\[
\mathbb{P}(X = 0) = e^{-4.5} \approx 0.0111.
\]
\item Fixed $n = 50$ trials, two outcomes, constant $p = \tfrac{1}{4} = 0.25$, independent guesses: \textbf{binomial} $X \sim B(50, 0.25)$.

Here $np = 12.5 > 5$ and $nq = 37.5 > 5$, but $p$ is not small, so the Poisson approximation is unsuitable; use the \textbf{normal approximation} with continuity correction:
\[
X \approx Y \sim N(12.5,\ 9.375), \qquad \sigma = \sqrt{9.375} \approx 3.062.
\]
\[
\mathbb{P}(X \ge 15) \approx \mathbb{P}\!\left(Z > \frac{14.5 - 12.5}{3.062}\right) = \mathbb{P}(Z > 0.653) \approx 1 - 0.7431 = 0.2569.
\]
\item Breakdowns occur randomly in time at a constant average rate: \textbf{Poisson}. Rate $= \tfrac{1}{10}$ per day, so for $30$ days $\lambda = 3$ and $X \sim \mathrm{Po}(3)$.
\[
\mathbb{P}(X \ge 2) = 1 - \mathbb{P}(X=0) - \mathbb{P}(X=1) = 1 - e^{-3} - 3e^{-3} = 1 - 4e^{-3} \approx 1 - 0.1992 = 0.8008.
\]
\item Sampling \textbf{without replacement} from a finite population: the exact model is the \textbf{hypergeometric distribution}:
\[
\mathbb{P}(X=1) = \frac{\binom{5}{1}\binom{95}{9}}{\binom{100}{10}} \approx 0.3394.
\]
Since the sample fraction $\tfrac{n}{N} = 0.1$ is small, the \textbf{binomial approximation} $B(10, 0.05)$ is suitable:
\[
\mathbb{P}(X=1) \approx \binom{10}{1}(0.05)(0.95)^9 = 10 \times 0.05 \times 0.6302 = 0.3151.
\]
(A Poisson approximation with $\lambda = np = 0.5$ gives $\mathbb{P}(Y=1) = 0.5e^{-0.5} \approx 0.3033$, less accurate here because the binomial approximation is already good.)
\end{enumerate}
\end{exoresolu}

\begin{attention}[Choosing the wrong model]
Two frequent examination errors: (i) using the Poisson approximation when $p$ is not small (e.g. $p = 0.25$) — the results are then badly wrong; (ii) forgetting to rescale $\lambda$ when the interval changes (e.g. from $1$ hour to $30$ minutes, $\lambda$ must be halved). Always state the distribution and its parameters before computing.
\end{attention}

\newpage

\coursec{Exercises}

\courssub{I check my knowledge}

\begin{exercice}[MCQ: Identifying a binomial experiment]
Which of the following scenarios can be modelled by a binomial distribution?
\begin{enumerate}
\item The number of heads obtained when a fair coin is tossed 10 times.
\item The number of throws of a fair die until a six appears.
\item The number of defective items in a random sample of 20 items taken from a large batch where 5\% are defective, without replacement.
\item The number of customers arriving at a shop in a 10-minute interval, given an average of 3 customers per minute.
\end{enumerate}
\end{exercice}

\begin{exercice}[True/False: Properties of the Poisson distribution]
For each statement, state whether it is true or false and justify your answer.
\begin{enumerate}
\item If $X \sim \mathrm{Po}(\lambda)$, then $\mathrm{E}(X) = \mathrm{Var}(X) = \lambda$.
\item The Poisson distribution is always symmetric.
\item The Poisson distribution can be used to approximate a binomial distribution when $n$ is large and $p$ is small.
\item If $X \sim \mathrm{Po}(4)$, then the mode of $X$ is 4.
\end{enumerate}
\end{exercice}

\begin{exercice}[Definition: Geometric distribution]
Define the geometric distribution with parameter $p$. Write down its probability mass function, expectation, and variance. Give a concrete example of a random experiment that follows this distribution.
\end{exercice}

\begin{exercice}[Direct calculation: Binomial probability]
A fair die is thrown 10 times. Let $X$ be the number of sixes obtained.
\begin{enumerate}
\item State the distribution of $X$ with its parameters.
\item Find $\mathrm{E}(X)$ and $\mathrm{Var}(X)$.
\item Calculate $\mathrm{P}(X = 2)$.
\item Calculate $\mathrm{P}(X \ge 1)$.
\end{enumerate}
\end{exercice}

\courssub{I practise}

\begin{exercice}[Binomial distribution: Recurrence formula and mode]
Given $X \sim \mathrm{B}(12, 0.35)$.
\begin{enumerate}
\item Use the recurrence formula $\mathrm{P}(X = r+1) = \frac{n-r}{r+1} \cdot \frac{p}{1-p} \, \mathrm{P}(X = r)$ to compute $\mathrm{P}(X = 0)$, $\mathrm{P}(X = 1)$, and $\mathrm{P}(X = 2)$.
\item Hence find the mode of $X$.
\item Verify your result using the formula for the mode of a binomial distribution.
\end{enumerate}
\end{exercice}

\begin{exercice}[Cumulative binomial tables]
Using cumulative binomial tables for $X \sim \mathrm{B}(20, 0.15)$, find:
\begin{enumerate}
\item $\mathrm{P}(X = 4)$.
\item $\mathrm{P}(X < 3)$.
\item $\mathrm{P}(X \ge 5)$.
\item The smallest integer $k$ such that $\mathrm{P}(X \le k) \ge 0.95$.
\end{enumerate}
\end{exercice}

\begin{exercice}[Fitting a binomial distribution to data]
In an experiment, 200 samples of 5 seeds each were planted. The number of germinating seeds per sample was recorded as follows:
\begin{center}
\begin{tabular}{|c|cccccc|}
\hline
\rowcolor{popblueL} Number of germinating seeds ($x$) & 0 & 1 & 2 & 3 & 4 & 5 \\
\hline
Frequency ($f$) & 12 & 48 & 60 & 52 & 24 & 4 \\
\hline
\end{tabular}
\end{center}
\begin{enumerate}
\item Estimate the probability $p$ that a single seed germinates.
\item Fit a binomial distribution $\mathrm{B}(5, \hat{p})$ to the data and compute the expected frequencies.
\item Compare the observed and expected frequencies. Comment on the goodness of fit.
\end{enumerate}
\end{exercice}

\begin{exercice}[Poisson distribution: Recurrence formula and mode]
Let $X \sim \mathrm{Po}(5)$.
\begin{enumerate}
\item Use the recurrence formula $\mathrm{P}(X = r+1) = \frac{\lambda}{r+1} \mathrm{P}(X = r)$ to find the mode(s) of $X$.
\item Compute the probability(ies) at the mode(s).
\end{enumerate}
Now let $Y \sim \mathrm{Po}(4)$.
\begin{enumerate}
\setcounter{enumi}{2}
\item Find the mode(s) of $Y$ and compute the corresponding probability(ies).
\item Comment on the difference between the modes of $X$ and $Y$.
\end{enumerate}
\end{exercice}

\courssub{I prepare for the GCE}

\begin{exercice}[{Geometric distribution: Penalty kicks] [15 marks}]
A footballer scores on 20\% of his penalty kicks. He continues taking penalties until he scores his first goal.
\begin{enumerate}
\item Name the distribution that models the number of kicks needed to score the first goal, and state its parameter. [2]
\item Find the probability that his first goal occurs on the 4th kick. [3]
\item Find the probability that he needs more than 5 kicks to score his first goal. [3]
\item Find the expected number of kicks required to score the first goal. [2]
\item What is the probability that he scores at least one goal in his first 3 kicks? [3]
\item If he takes 10 penalties, what is the probability that he scores exactly 2 goals? Name the distribution used. [2]
\end{enumerate}
\end{exercice}

\begin{exercice}[{Poisson distribution: Call centre and rare disease] [20 marks}]
\textbf{Part A: Call centre}
Calls arrive at an office at an average rate of 3 per minute. Assume that the number of calls in disjoint time intervals are independent and that the rate is constant.
\begin{enumerate}
\item State the distribution of the number of calls in a one-minute interval, giving its parameter. [2]
\item Find the probability of exactly 2 calls in a given minute. [2]
\item Find the probability of no call in a 30-second interval. [3]
\item Find the probability of more than 8 calls in a 2-minute interval. [3]
\item State two assumptions necessary for the Poisson model to be valid in this context. [2]
\end{enumerate}

\textbf{Part B: Rare disease}
A rare disease affects 0.5\% of a population. In a village of 400 people, a health survey is conducted.
\begin{enumerate}
\setcounter{enumi}{5}
\item State an exact probability model for the number of cases in the village, and a suitable approximation. Justify the approximation. [3]
\item Using the approximation, find the probability of exactly 2 cases. [2]
\item Find the probability of at least one case. [2]
\item Find the expected number of cases and comment on the result. [1]
\end{enumerate}
\end{exercice}

\begin{exercice}[{Comparison of distributions and real-life problems] [20 marks}]
For each of the following situations, choose the most appropriate distribution from: Binomial, Geometric, Poisson. Justify your choice by stating the conditions that are satisfied. Then solve the probability question asked.
\begin{enumerate}
\item \textbf{Exam guessing:} A multiple-choice test has 20 questions, each with 4 options. A student guesses all answers. Find the probability that the student gets exactly 5 questions correct. [4]
\item \textbf{First success in a game:} A player wins a game with probability 0.3 each time they play, independently. Find the probability that the player's first win occurs on the 4th game. [4]
\item \textbf{Accidents per week:} Accidents occur at a factory at an average rate of 2 per week. Find the probability of exactly 3 accidents in a given week. [4]
\item \textbf{Defective items in a sample:} A batch of 100 items contains 5 defectives. A sample of 10 items is selected at random without replacement. Find the probability that the sample contains exactly 1 defective item. [4]
\item \textbf{Emails per hour:} Emails arrive at a server at an average rate of 12 per hour. Find the probability that more than 15 emails arrive in a given hour. [4]
\end{enumerate}
\end{exercice}

\newpage

\coursec{Solutions to the exercises}

\begin{corrige}[Exercise 1 — MCQ: Identifying a binomial experiment]
\begin{enumerate}
\item \textbf{Yes.} Fixed number of trials ($n = 10$), two outcomes (head/tail), constant probability $p = 0.5$, independent tosses. $X \sim \mathrm{B}(10, 0.5)$.
\item \textbf{No.} The number of trials is not fixed; we count trials until the first success. This is a \cle{geometric} distribution with $p = \tfrac{1}{6}$.
\item \textbf{Yes (as an approximation).} Sampling is without replacement, so strictly the distribution is hypergeometric; however the batch is \emph{large}, so $p = 0.05$ stays essentially constant and $X \approx \mathrm{B}(20, 0.05)$ is an excellent model.
\item \textbf{No.} There is no fixed number of trials; events occur randomly in continuous time at an average rate. This is a \cle{Poisson} distribution, $\mathrm{Po}(30)$ for a 10-minute interval.
\end{enumerate}
\end{corrige}

\begin{corrige}[Exercise 2 — True/False: Properties of the Poisson distribution]
\begin{enumerate}
\item \textbf{True.} For $X \sim \mathrm{Po}(\lambda)$, both the expectation and the variance equal the parameter: $\mathrm{E}(X) = \mathrm{Var}(X) = \lambda$. This equality is a signature of the Poisson model.
\item \textbf{False.} The Poisson distribution is skewed to the right (positively skewed), especially for small $\lambda$. It only becomes \emph{approximately} symmetric (bell-shaped) when $\lambda$ is large.
\item \textbf{True.} When $n$ is large, $p$ is small and $\lambda = np$ is moderate, $\mathrm{B}(n,p) \approx \mathrm{Po}(np)$. The usual rule of thumb is $n \ge 50$ and $p \le 0.1$.
\item \textbf{True (with a precision).} Since $\lambda = 4$ is an integer, the Poisson distribution has \emph{two} modes: $\lambda - 1 = 3$ and $\lambda = 4$, with $\mathrm{P}(X = 3) = \mathrm{P}(X = 4)$. So 4 is indeed a mode, but it shares this status with 3.
\end{enumerate}
\end{corrige}

\begin{corrige}[Exercise 3 — Definition: Geometric distribution]
\textbf{Definition (Geometric distribution).} A discrete random variable $X$ follows a \cle{geometric distribution} with parameter $p$ ($0 < p < 1$), written $X \sim \mathrm{Geo}(p)$, if $X$ counts the number of \textbf{independent trials needed to obtain the first success}, where each trial has constant success probability $p$.

\textbf{Probability mass function:}
\[ \mathrm{P}(X = r) = (1-p)^{r-1} p, \qquad r = 1, 2, 3, \dots \]
\textbf{Expectation and variance:}
\[ \mathrm{E}(X) = \frac{1}{p}, \qquad \mathrm{Var}(X) = \frac{1-p}{p^{2}}. \]
\textbf{Concrete example:} a candidate at the GCE tries a multiple-choice question by guessing ($p = \tfrac{1}{4}$); the number of guesses needed before the first correct answer follows $\mathrm{Geo}(0.25)$. Another example: the number of throws of a fair die until the first six appears, $X \sim \mathrm{Geo}\!\left(\tfrac{1}{6}\right)$.
\end{corrige}

\begin{corrige}[Exercise 4 — Direct calculation: Binomial probability]
\begin{enumerate}
\item Each throw is an independent trial with success (a six) probability $p = \tfrac{1}{6}$, and there are $n = 10$ fixed trials:
\[ X \sim \mathrm{B}\!\left(10, \tfrac{1}{6}\right). \]
\item 
\[ \mathrm{E}(X) = np = 10 \times \frac{1}{6} = \frac{5}{3} \approx 1.67, \qquad
\mathrm{Var}(X) = np(1-p) = 10 \times \frac{1}{6} \times \frac{5}{6} = \frac{25}{18} \approx 1.39. \]
\item 
\begin{align*}
\mathrm{P}(X = 2) &= \binom{10}{2}\left(\frac{1}{6}\right)^{2}\left(\frac{5}{6}\right)^{8}
= 45 \times \frac{1}{36} \times 0.23257 \\
&\approx 0.2907.
\end{align*}
\item Using the complement:
\[ \mathrm{P}(X \ge 1) = 1 - \mathrm{P}(X = 0) = 1 - \left(\frac{5}{6}\right)^{10} = 1 - 0.16151 \approx 0.8385. \]
\end{enumerate}
\end{corrige}

\begin{corrige}[Exercise 5 — Binomial distribution: Recurrence formula and mode]
\begin{enumerate}
\item With $n = 12$, $p = 0.35$, $\dfrac{p}{1-p} = \dfrac{0.35}{0.65} = \dfrac{7}{13} \approx 0.53846$.
\[ \mathrm{P}(X = 0) = 0.65^{12} \approx 0.005687. \]
\[ \mathrm{P}(X = 1) = \frac{12 - 0}{0 + 1} \times \frac{7}{13} \times \mathrm{P}(X=0) = 12 \times 0.53846 \times 0.005687 \approx 0.03675. \]
\[ \mathrm{P}(X = 2) = \frac{12 - 1}{1 + 1} \times \frac{7}{13} \times \mathrm{P}(X=1) = \frac{11}{2} \times 0.53846 \times 0.03675 \approx 0.10883. \]
\item The ratio $\dfrac{\mathrm{P}(X=r+1)}{\mathrm{P}(X=r)} = \dfrac{12-r}{r+1}\cdot\dfrac{7}{13}$ exceeds 1 while $7(12-r) > 13(r+1)$, i.e. $84 - 7r > 13r + 13$, i.e. $r < 3.55$. So probabilities increase up to $r = 4$ and decrease afterwards: the \textbf{mode is $4$}.
\item Formula check: $(n+1)p = 13 \times 0.35 = 4.55$, not an integer, so the mode is $\lfloor 4.55 \rfloor = 4$. \checkmark
\end{enumerate}
\end{corrige}

\begin{corrige}[Exercise 6 — Cumulative binomial tables]
For $X \sim \mathrm{B}(20, 0.15)$ the cumulative tables give:
\[ \mathrm{P}(X \le 2) = 0.4049, \quad \mathrm{P}(X \le 3) = 0.6477, \quad \mathrm{P}(X \le 4) = 0.8298, \]
\[ \mathrm{P}(X \le 5) = 0.9327, \quad \mathrm{P}(X \le 6) = 0.9781. \]
\begin{enumerate}
\item $\mathrm{P}(X = 4) = \mathrm{P}(X \le 4) - \mathrm{P}(X \le 3) = 0.8298 - 0.6477 = 0.1821.$
\item $\mathrm{P}(X < 3) = \mathrm{P}(X \le 2) = 0.4049.$
\item $\mathrm{P}(X \ge 5) = 1 - \mathrm{P}(X \le 4) = 1 - 0.8298 = 0.1702.$
\item $\mathrm{P}(X \le 5) = 0.9327 < 0.95$ but $\mathrm{P}(X \le 6) = 0.9781 \ge 0.95$, so the smallest such integer is $k = 6$.
\end{enumerate}
\end{corrige}

\begin{corrige}[Exercise 7 — Fitting a binomial distribution to data]
\begin{enumerate}
\item Total number of germinating seeds:
\[ \sum fx = 0(12) + 1(48) + 2(60) + 3(52) + 4(24) + 5(4) = 48 + 120 + 156 + 96 + 20 = 440. \]
Sample mean: $\bar{x} = \dfrac{440}{200} = 2.2$ germinating seeds per sample of 5.
Since $\mathrm{E}(X) = n\hat{p} = 5\hat{p}$, we get $\hat{p} = \dfrac{2.2}{5} = 0.44$.
\item Expected frequencies $E_r = 200\binom{5}{r}(0.44)^r(0.56)^{5-r}$:
\begin{center}
\begin{tabular}{|c|cccccc|}
\hline
\rowcolor{popblueL} $x$ & 0 & 1 & 2 & 3 & 4 & 5 \\
\hline
Observed $f$ & 12 & 48 & 60 & 52 & 24 & 4 \\
\hline
Expected $E$ & 11.0 & 43.3 & 68.0 & 53.4 & 21.0 & 3.3 \\
\hline
\end{tabular}
\end{center}
\item The expected frequencies follow the same pattern as the observed ones: both rise to a peak at $x = 2$ and then fall. The largest discrepancy is at $x = 2$ (60 observed vs 68.0 expected), but the differences are small relative to the frequencies, and $\sum E = 200$ as required. The binomial model $\mathrm{B}(5, 0.44)$ fits the data \textbf{well}: germination behaves like 5 independent trials with constant probability $p \approx 0.44$.
\end{enumerate}
\end{corrige}

\begin{corrige}[Exercise 8 — Poisson distribution: Recurrence formula and mode]
\textbf{Case $X \sim \mathrm{Po}(5)$.}
\begin{enumerate}
\item The ratio of consecutive probabilities is
\[ \frac{\mathrm{P}(X = r+1)}{\mathrm{P}(X = r)} = \frac{5}{r+1}. \]
This is $> 1$ when $r + 1 < 5$, i.e. $r \le 3$ (probabilities increase), and $= 1$ when $r + 1 = 5$, i.e. $\mathrm{P}(X = 5) = \mathrm{P}(X = 4)$. For $r \ge 5$ the probabilities decrease. Hence the \textbf{modes are $4$ and $5$}.
\item 
\[ \mathrm{P}(X = 4) = \mathrm{P}(X = 5) = e^{-5}\,\frac{5^{4}}{4!} = e^{-5} \times \frac{625}{24} \approx 0.0067379 \times 26.0417 \approx 0.1755. \]
\end{enumerate}
\textbf{Case $Y \sim \mathrm{Po}(4)$.}
\begin{enumerate}
\setcounter{enumi}{2}
\item Similarly $\dfrac{\mathrm{P}(Y=r+1)}{\mathrm{P}(Y=r)} = \dfrac{4}{r+1}$, which equals 1 when $r = 3$: the \textbf{modes are $3$ and $4$}, with
\[ \mathrm{P}(Y = 3) = \mathrm{P}(Y = 4) = e^{-4}\,\frac{4^{4}}{4!} = e^{-4} \times \frac{256}{24} \approx 0.0183156 \times 10.6667 \approx 0.1954. \]
\item In both cases $\lambda$ is an integer, so each distribution has \emph{two} consecutive modes, $\lambda - 1$ and $\lambda$. The modal probability is larger for $\mathrm{Po}(4)$ ($0.1954 > 0.1755$): a smaller $\lambda$ concentrates the distribution more sharply near its centre, while a larger $\lambda$ spreads the probability mass over more values.
\end{enumerate}
\end{corrige}

\begin{corrige}[Exercise 9 — Geometric distribution: Penalty kicks]
Let $p = 0.2$ (probability of scoring), $q = 0.8$.
\begin{enumerate}
\item The number of kicks until the first goal follows a \cle{geometric distribution}: $X \sim \mathrm{Geo}(0.2)$. \hfill [2 marks: name 1, parameter 1]
\item $\mathrm{P}(X = 4) = q^{3}p = 0.8^{3} \times 0.2 = 0.512 \times 0.2 = 0.1024.$ \hfill [3 marks: formula 1, substitution 1, answer 1]
\item ``More than 5 kicks'' means the first 5 kicks all fail:
\[ \mathrm{P}(X > 5) = q^{5} = 0.8^{5} = 0.32768 \approx 0.3277. \] \hfill [3 marks]
\item $\mathrm{E}(X) = \dfrac{1}{p} = \dfrac{1}{0.2} = 5$ kicks. \hfill [2 marks]
\item ``At least one goal in the first 3 kicks'' is the complement of ``no goal in 3 kicks'':
\[ \mathrm{P}(X \le 3) = 1 - q^{3} = 1 - 0.512 = 0.488. \] \hfill [3 marks]
\item Now the number of trials is fixed at $n = 10$ and we count successes: $Y \sim \mathrm{B}(10, 0.2)$.
\[ \mathrm{P}(Y = 2) = \binom{10}{2}(0.2)^{2}(0.8)^{8} = 45 \times 0.04 \times 0.16777 \approx 0.3020. \] \hfill [2 marks: binomial named 1, probability 1]
\end{enumerate}
\textbf{Examiner's expectations:} the candidate must clearly distinguish the \emph{geometric} setting (number of trials until first success) from the \emph{binomial} setting (number of successes in a fixed number of trials). Writing $\mathrm{P}(X > 5) = q^5$ directly (memorylessness) earns full credit; a term-by-term sum is accepted but slower.
\end{corrige}

\begin{corrige}[Exercise 10 — Poisson distribution: Call centre and rare disease]
\textbf{Part A: Call centre}
\begin{enumerate}
\item Calls occur randomly, independently, at a constant average rate of 3 per minute, so the number of calls in one minute is $X \sim \mathrm{Po}(3)$. \hfill [2 marks]
\item $\mathrm{P}(X = 2) = e^{-3}\dfrac{3^{2}}{2!} = 0.049787 \times 4.5 \approx 0.2240.$ \hfill [2 marks]
\item In 30 seconds the mean is $\lambda = 3 \times 0.5 = 1.5$, so $N \sim \mathrm{Po}(1.5)$:
\[ \mathrm{P}(N = 0) = e^{-1.5} \approx 0.2231. \] \hfill [3 marks: new $\lambda$ 1, model 1, answer 1]
\item In 2 minutes, $\lambda = 6$, so $M \sim \mathrm{Po}(6)$. From cumulative Poisson tables, $\mathrm{P}(M \le 8) = 0.8472$:
\[ \mathrm{P}(M > 8) = 1 - 0.8472 = 0.1528. \] \hfill [3 marks]
\item Any two of: calls arrive \textbf{independently} of each other; calls arrive at a \textbf{constant average rate}; calls occur \textbf{singly} (two calls cannot arrive at exactly the same instant); calls occur \textbf{randomly} in time. \hfill [2 marks]
\end{enumerate}
\textbf{Part B: Rare disease}
\begin{enumerate}
\setcounter{enumi}{5}
\item Exact model: $Y \sim \mathrm{B}(400, 0.005)$. Since $n = 400$ is large, $p = 0.005$ is small and $\lambda = np = 2$ is moderate, the suitable approximation is
\[ Y \approx \mathrm{Po}(2). \] \hfill [3 marks: exact 1, approximation 1, justification 1]
\item $\mathrm{P}(Y = 2) = e^{-2}\dfrac{2^{2}}{2!} = 2e^{-2} \approx 0.2707.$ \hfill [2 marks]
\item $\mathrm{P}(Y \ge 1) = 1 - \mathrm{P}(Y = 0) = 1 - e^{-2} \approx 1 - 0.1353 = 0.8647.$ \hfill [2 marks]
\item $\mathrm{E}(Y) = \lambda = 2$ cases. In a village of 400 people one expects about 2 cases of the disease, which confirms that the disease is rare but likely to appear in a community of this size. \hfill [1 mark]
\end{enumerate}
\textbf{Examiner's expectations:} marks are awarded for \emph{adjusting the rate to the interval} (questions 3 and 4), for quoting the conditions $n$ large, $p$ small, $np$ moderate when justifying the Poisson approximation, and for using the complement in question 8.
\end{corrige}

\begin{corrige}[Exercise 11 — Comparison of distributions and real-life problems]
\begin{enumerate}
\item \textbf{Binomial.} Fixed number of trials ($n = 20$), two outcomes per question, constant success probability $p = \tfrac{1}{4} = 0.25$, independent guesses: $X \sim \mathrm{B}(20, 0.25)$.
\[ \mathrm{P}(X = 5) = \binom{20}{5}(0.25)^{5}(0.75)^{15} = 15504 \times 0.0009766 \times 0.013363 \approx 0.2023. \] \hfill [4 marks]
\item \textbf{Geometric.} Independent plays, constant win probability $p = 0.3$, counting trials until the first win: $X \sim \mathrm{Geo}(0.3)$.
\[ \mathrm{P}(X = 4) = (0.7)^{3}(0.3) = 0.343 \times 0.3 = 0.1029. \] \hfill [4 marks]
\item \textbf{Poisson.} Accidents occur randomly, independently, at a constant average rate $\lambda = 2$ per week, with no fixed number of trials: $X \sim \mathrm{Po}(2)$.
\[ \mathrm{P}(X = 3) = e^{-2}\,\frac{2^{3}}{3!} = e^{-2} \times \frac{8}{6} \approx 0.1353 \times 1.3333 \approx 0.1804. \] \hfill [4 marks]
\item \textbf{Binomial (approximation).} Strictly, sampling without replacement gives a hypergeometric distribution; however the batch ($100$ items) is large compared with the sample ($10$), so the defective proportion $p = 0.05$ stays nearly constant and $X \approx \mathrm{B}(10, 0.05)$:
\[ \mathrm{P}(X = 1) = \binom{10}{1}(0.05)(0.95)^{9} = 0.5 \times 0.63025 \approx 0.3151. \]
(The exact hypergeometric value is about $0.34$; the binomial approximation is the expected GCE answer when a choice among the three distributions is imposed.) \hfill [4 marks]
\item \textbf{Poisson.} Emails arrive randomly and independently at a constant average rate $\lambda = 12$ per hour: $X \sim \mathrm{Po}(12)$. From cumulative Poisson tables, $\mathrm{P}(X \le 15) \approx 0.8444$:
\[ \mathrm{P}(X > 15) = 1 - \mathrm{P}(X \le 15) \approx 1 - 0.8444 = 0.1556. \] \hfill [4 marks]
\end{enumerate}
\textbf{Examiner's expectations:} for each situation the mark scheme rewards (i) the correct distribution by name, (ii) the parameters, (iii) a justification quoting the defining conditions (fixed $n$ and constant $p$ for binomial; ``until the first success'' for geometric; random independent events at constant rate for Poisson), and (iv) the correct numerical probability rounded to 4 decimal places. In question 4, mentioning that the sampling is without replacement and that the binomial is an approximation shows the deeper understanding that earns full credit.
\end{corrige}

\newpage

\coursec{Summary sheet}

\begin{popfigure}
\begin{tikzpicture}[scale=0.9, every node/.style={font=\small}]
\node[draw, circle, fill=popblue, text=white, minimum size=1.2cm, align=center] (center){Special Discrete\\Distributions};
\foreach \angle/\color/\text in {
    0/poporange/Binomial,
    51.4/popgreen/Geometric,
    102.9/poppurple/Poisson,
    154.3/poppink/Poisson Approx.,
    205.7/popgold/Fitting \& Tables,
    257.1/popblueL/Applications,
    308.6/popmuted/Comparisons
} {
    \draw[thick, \color] (center) -- ++(\angle:3.5cm) node[draw, rounded corners, fill=\color!20, text=\color!80!black, anchor=center] {\text};
}
\end{tikzpicture}
\legende{Mind map of Chapter PMS12: Special Discrete Probability Distributions}
\end{popfigure}

\begin{bilan}
\textbf{Key Definitions}
\begin{itemize}
\item \cle{Binomial Distribution}: $X\sim B(n,p)$ with $n$ independent Bernoulli trials, success probability $p$. PMF: $P(X=k)=\binom{n}{k}p^k(1-p)^{n-k}$, $k=0,\dots,n$.
\item \cle{Geometric Distribution}: $X\sim Geo(p)$ counting trials until first success (including success). PMF: $P(X=k)=p(1-p)^{k-1}$, $k\ge1$. Alternative: failures before first success $P(Y=k)=p(1-p)^k$, $k\ge0$.
\item \cle{Poisson Distribution}: $X\sim Po(\lambda)$ for rare events in a fixed interval. PMF: $P(X=k)=e^{-\lambda}\frac{\lambda^k}{k!}$, $k\ge0$.
\end{itemize}

\textbf{Laws, Formulas, Theorems}
\begin{itemize}
\item \textbf{Expectation \& Variance}:
  \begin{align*}
  B(n,p):&\quad E(X)=np,\quad Var(X)=np(1-p).\\
  Geo(p):&\quad E(X)=\frac1p,\quad Var(X)=\frac{1-p}{p^2}.\\
  Po(\lambda):&\quad E(X)=\lambda,\quad Var(X)=\lambda.
  \end{align*}
\item \textbf{Recurrence Formulae}:
  \begin{align*}
  B(n,p):&\quad P(X=k+1)=\frac{n-k}{k+1}\frac{p}{1-p}P(X=k).\\
  Po(\lambda):&\quad P(X=k+1)=\frac{\lambda}{k+1}P(X=k).
  \end{align*}
\item \textbf{Mode}:
  \begin{itemize}
  \item Binomial: $\lfloor (n+1)p\rfloor$ (or two modes if $(n+1)p$ integer).
  \item Poisson: $\lfloor\lambda\rfloor$ (or $\lambda$ and $\lambda-1$ if $\lambda$ integer).
  \end{itemize}
\item \textbf{Poisson Approximation to Binomial}: If $n$ large, $p$ small, $np=\lambda$ moderate, then $B(n,p)\approx Po(\lambda)$. Rule of thumb: $n\ge50$, $p\le0.05$, $np\le5$.
\item \textbf{Cumulative Distribution Function}: $F(k)=P(X\le k)=\sum_{i=0}^k P(X=i)$. Use binomial tables for $n\le50$; for larger $n$ use normal approximation (not in this chapter).
\end{itemize}

\textbf{Methods}
\begin{enumerate}
\item \textbf{Identify the distribution}: Check for fixed $n$ and two outcomes $\to$ Binomial; waiting for first success $\to$ Geometric; rare events in continuous interval $\to$ Poisson.
\item \textbf{Compute probabilities}: Use PMF directly, recurrence to build successive terms, or cumulative tables.
\item \textbf{Fit a binomial to data}: Estimate $p$ by $\hat p=\frac{\text{total successes}}{n\times\text{frequency}}$; compute expected frequencies $E_i=N\times P(X=i)$; check $\chi^2$ goodness-of-fit (if required).
\item \textbf{Approximate binomial by Poisson}: When conditions hold, replace $B(n,p)$ with $Po(np)$ for easier calculation.
\item \textbf{Solve real-life problems}: Translate word problem into parameters, choose distribution, compute required probability or expectation, interpret answer in context.
\end{enumerate}

\textbf{Common Mistakes}
\begin{itemize}
\item Confusing the two definitions of geometric distribution (trials vs failures). Always state which version you use.
\item Forgetting that binomial requires \emph{independent} trials with constant $p$.
\item Using Poisson approximation when $np$ is not moderate (e.g. $np>5$ or $p$ not small).
\item Misreading cumulative tables: $P(X\le k)$ vs $P(X<k)$.
\item In fitting, not pooling expected frequencies $<5$ before $\chi^2$ test.
\item Mixing up variance formulas: $np(1-p)$ for binomial, $\lambda$ for Poisson, $\frac{1-p}{p^2}$ for geometric.
\end{itemize}

\textbf{GCE Tips}
\begin{itemize}
\item \textbf{Show all steps}: Write the distribution name, parameters, formula, substitution, and final answer.
\item \textbf{Use recurrence} to compute several probabilities quickly in exams; it saves time and reduces arithmetic errors.
\item \textbf{Tables}: If $n\le50$ and $p$ is in the table, use it; otherwise compute manually or approximate.
\item \textbf{Approximation justification}: State clearly why Poisson approximation is valid (large $n$, small $p$, $np$ moderate).
\item \textbf{Interpretation}: Always conclude with a sentence answering the question in context (e.g. ``The probability that at most 2 defects occur is 0.92'').
\item \textbf{Time management}: Binomial and Poisson questions often carry 6-8 marks; allocate 8-10 minutes each.
\end{itemize}
\end{bilan}

\findecours

\end{document}
